%% Multiplicity one for SL(2) -- David Joyner
%% Draft 2 (2026) of the 1995 Scientific Word manuscript aabook.tex and its
%% \input files, assembled into a single file.  See ../revision_log.md.
\documentclass[12pt,letterpaper]{book}
\usepackage[american]{babel}
\usepackage[T1]{fontenc}
\usepackage{amsmath,amssymb,amsthm}
\usepackage{mathptmx}          % Times text and math
\usepackage{courier}           % Courier for typewriter text
\usepackage[letterpaper,margin=1in]{geometry}

% --- theorem environments (formerly from the thmsa option) ---
\theoremstyle{plain}
\newtheorem{theorem}{Theorem}[chapter]
\newtheorem{lemma}[theorem]{Lemma}
\newtheorem{proposition}[theorem]{Proposition}
\newtheorem{corollary}[theorem]{Corollary}
\newtheorem{conjecture}[theorem]{Conjecture}
\newtheorem{hypothesis}[theorem]{Hypothesis}
\theoremstyle{definition}
\newtheorem{definition}[theorem]{Definition}
\newtheorem{example}[theorem]{Example}
\theoremstyle{remark}
\newtheorem{remark}[theorem]{Remark}
\newtheorem{claim}[theorem]{Claim}

\newcommand{\tr}{\operatorname{tr}}

% --- the book class has no abstract environment ---
\newenvironment{abstract}%
  {\begin{center}\bfseries Abstract\end{center}\begin{quote}}%
  {\end{quote}}

\begin{document}

\author{David Joyner}
\title{Multiplicity one for $SL(2)$\\[2ex]
\large Draft 2\\[0.5ex] Corrected and expanded version}
\date{Original draft: 1995. Corrected and expanded draft: 2026}
\maketitle

\begin{abstract}
This revision was prepared with the help of AI systems, whose roles are
described in the Acknowledgments at the end of Chapter~\ref{ch:intro}.
We apply the trace formula to compare cuspidal automorphic
representations of $SL(2,\mathbb{A})$ with genuine cuspidal
automorphic representations of the 2-fold metaplectic cover
$\overline{SL(2,\mathbb{A})}$, where $\mathbb{A}$ is the ring of adeles
of a number field $k$. \emph{The main result is conditional.}
Let $\pi$ be a cuspidal automorphic representation of $SL(2,\mathbb{A})$
and let $S(\pi)$ be the set of archimedean places, dyadic places, and
odd places $v$ at which $\pi_v$ has no vector fixed by $SL(2,{\cal O}_v)$.
Assuming a global trace identity, refined according to archimedean
infinitesimal characters (Hypothesis~\ref{hyp:GH}), and a local
spectral transfer hypothesis (Hypothesis~\ref{hyp:LH}) at each place of
$S(\pi)$, we show that $\pi$ occurs with multiplicity one and
corresponds to a unique genuine cuspidal automorphic representation of
the cover. The local hypothesis is verified at odd places where $\pi_v$
is spherical with respect to a hyperspecial maximal compact subgroup.
Multiplicity one for $SL(2)$ is known unconditionally by the work of
Ramakrishnan; the interest here is the method, which derives it from
Waldspurger's multiplicity one theorem for the metaplectic cover.
The 1995 draft claimed an unconditional proof; this version explains
which steps of that argument are incomplete, and marks as ``Claims''
the local statements of the 1995 draft that have not been verified.
\end{abstract}

\tableofcontents

\chapter{Introduction}\label{ch:intro}

The goal of this text is to use the trace formula to study the
multiplicity one problem for cuspidal automorphic representations of
$SL(2)$, by deducing it from the multiplicity one result for
the 2-fold metaplectic cover of $SL(2)$, where it is a theorem
of Waldspurger. The motivation for this problem
has its origins in several stimulating conversations
with R. Langlands in the late 1980s.

\section{Status of the main result}

Multiplicity one for the cuspidal spectrum of $SL(2)$ was conjectured
by Labesse and Langlands \cite{LL}, \cite{Lab}, and was proved by
Ramakrishnan \cite{Ra}. D. Blasius has shown that the analog of the
multiplicity one result is false for $SL(n)$, $n>2$ \cite{Blas}.
Strong multiplicity one fails for $SL(2)$ \cite{LL}; Chai and Zhang
\cite{CZ} (see also \cite{Zh}) proved a qualified form of it, for pairs
of cuspidal representations which also agree at the archimedean places
and at the places above $2$. We stress that the following four
properties are different and are not used interchangeably here:
multiplicity one (each cuspidal representation occurs at most once),
strong multiplicity one (a cuspidal representation is determined by
almost all of its local components), the local theta correspondence,
and the description of global near-equivalence classes (``Waldspurger
packets'') on the metaplectic cover \cite{Wa1}, \cite{Wa2}, \cite{PS}.

The 1995 draft of this text claimed an unconditional proof of
multiplicity one for $SL(2)$ by a comparison of trace formulas.
On re-examination, several steps of that argument are incomplete:
\begin{itemize}
\item at the dyadic places, the generalized Hecke algebra argument of the
1995 draft contains errors, and the orbital integral identity it
needed was derived from a global argument that presupposed the
equality of the cuspidal spectral sides (see \S\ref{sec:dyadic});
\item at the archimedean places, the matchings of \S\ref{sec:realA} and
\S\ref{sec:realds} control only some of the representations
(for instance, the elliptic matchings see only functions whose
orbital integrals vanish on the hyperbolic set);
\item at odd places, several of the local character relations are
incomplete (see Remark~\ref{rem:sc-caveat} and
Remark~\ref{rem:savin-caveat}), the correspondence of characters of
the diagonal torus in the 1995 draft was inconsistent when
$(\varpi,\varpi)_2=-1$ (corrected in \S\ref{sec:charsA}), and the
spherical Hecke algebra isomorphism gives the matching of hyperbolic,
but not of elliptic, orbital integrals (Remark~\ref{rem:satake-oi});
\item in the global argument, weighted orbital integrals cannot be
matched for arbitrary corresponding pairs
(Remark~\ref{rem:weighted-noninv}); and the final step of the 1995
argument contained two further gaps (see \S\ref{sec:1995-argument}).
\end{itemize}
This revision therefore states the main result conditionally, and
isolates exactly what is assumed.

Informally, the main result (Theorem~\ref{thm:main}) is as follows.
Fix a cuspidal automorphic representation $\pi=\otimes_v\pi_v$ of
$SL(2,\mathbb{A})$, and let $S=S(\pi)$ be the finite set consisting of the
archimedean places, the dyadic places, and the odd places $v$ at which
$\pi_v$ is not spherical with respect to $K_{0,v}=SL(2,{\cal O}_v)$.
Assume:
\begin{itemize}
\item (Hypothesis~\ref{hyp:GH}) for global functions $f$ on
$SL(2,\mathbb{A})$ and genuine global functions $\overline{f}$ on the cover
which correspond at every place (by the Satake correspondence outside
$S$, and through prescribed local classes of pairs at the places of
$S$), the traces of the cuspidal spectra agree after restriction to
each archimedean infinitesimal character $\nu$:
$\tr R_{cusp}^\nu(f)=\tr \overline{R}{}_{cusp}^{\,\nu}(\overline{f})$;
\item (Hypothesis~\ref{hyp:LH}) at each $v\in S$, the local class of
pairs separates characters, detects $\pi_v$, and transfers each genuine
irreducible representation to at most one irreducible representation
of $SL(2,k_v)$, with only finitely many genuine infinitesimal
characters over each infinitesimal character of $SL(2,k_v)$ when $v$
is archimedean.
\end{itemize}
Then $\pi$ occurs with multiplicity one, and $\pi$ corresponds to a
unique genuine cuspidal automorphic representation $\overline{\pi}$ of
the cover. The local hypothesis is verified here at odd places where
$\pi_v$ is spherical with respect to a hyperspecial maximal compact
subgroup (Proposition~\ref{prop:LH-spherical}).
Chapters \ref{ch:sc} and \ref{ch:local} contain partial results toward
the local hypothesis at other places, which we record together with
their limitations in \S\ref{sec:LH-status}.

The refinement by infinitesimal characters in Hypothesis~\ref{hyp:GH} is
needed because the cuspidal spectrum at a fixed level and with fixed
$K_\infty$-types is infinite (Remark~\ref{rem:why-refine}); Draft~1
overlooked this. Hypothesis~\ref{hyp:GH} is a strong assumption: it
contains most of the global comparison that the 1995 project hoped to
establish.

The proof of Theorem~\ref{thm:main} uses Waldspurger's multiplicity one
theorem for the cover, which is a statement about the space ${\cal A}_{00}$
of genuine cusp forms orthogonal to the elementary theta series. We show
(Lemma~\ref{lem:A00}) that the genuine representations which arise in
the proof lie in ${\cal A}_{00}$; this step uses the Labesse--Langlands
description of the cuspidal spectrum of $SL(2)$ by restriction from
$GL(2)$ and the Gelbart--Jacquet bound toward the Ramanujan conjecture
for $GL(2)$. The proof does not use the failure, or any weak form, of
strong multiplicity one for the cover.

\section{Conventions on the status of statements}

Statements labeled Theorem, Proposition, Lemma, or Corollary are either
proved here or quoted, with a reference, from the literature. Statements
of the 1995 draft that have not been verified in this revision, or that
are known to need correction, are labeled \textbf{Claim}; the qualification
is attached to the statement itself. Statements labeled Conjecture or
Hypothesis are assumptions, and are named as such wherever they are used.

\section{Outline}

The table of contents gives some idea about the content of each chapter.
Chapter~\ref{ch:background} recalls the metaplectic covers of $SL(2)$,
locally and globally. Chapter~\ref{ch:reps} surveys the representation
theory of the local cover. Chapter~\ref{ch:sc} discusses a correspondence
between supercuspidal representations for $p>2$.
Chapter~\ref{ch:local} presents the local correspondences needed in the
comparison between the trace formula of $SL(2)$ and its 2-fold cover.
They are divided into three types according to the residual
characteristic of the field: $p=\infty$ (real or complex), $p=2$,
and $p>2$. For $p>2$ and $p=\infty$, the idea is to show that if there
is a certain type of identity between the orbital integral of a smooth
function $f$ on $SL(2,F)$ and the orbital integral of a smooth function
$\overline{f}$ on the two-fold cover $\overline{SL(2,F)}$ (call this a
``correspondence'' or ``matching'' of functions), then there is an
analogous identity between the trace maps
$f^\wedge:\pi\longmapsto \tr(\pi(f))$
and $\overline{f}^\wedge:\overline{\pi}\longmapsto
\tr(\overline{\pi}(\overline{f}))$
on the unitary duals (call this a ``lifting'' or ``transfer''
of genuine virtual representations of $\overline{SL(2,F)}$ to
virtual representations of $SL(2,F)$).
In \cite{L4} and \cite{Fl}, only two types
of correspondences or matchings were needed: the
``Satake correspondence'' of spherical functions and the
``smooth matching'' of Schwartz-Bruhat functions.
In the case of $SL(2)$, some new
types of ``matchings'' are introduced. No matching is established at
$p=2$ in this version (\S\ref{sec:dyadic}).

Chapter~\ref{ch:tf} recalls the trace formulas of $SL(2)$ and its cover.
Chapter~\ref{ch:global} contains the application. The comparison is
modeled on the Corvallis paper of Gelbart and Jacquet \cite{GJ}.
Under the hypotheses above, comparing the cuspidal spectra
and applying Waldspurger's multiplicity one theorem for the cover gives
multiplicity one for $\pi$.

\section{Notation}\label{sec:notation}

The following symbols are used throughout. Local notation refers to a
local field $F$; global notation refers to a number field $k$.

\medskip
\begin{center}
\renewcommand{\arraystretch}{1.15}
\begin{tabular}{p{0.27\textwidth}p{0.66\textwidth}}
\hline
$G=SL(2)$, $G(F)$ & the special linear group and its $F$-points\\
$\overline{G(F)}$, $\overline{H}$ & the 2-fold metaplectic cover, and the
inverse image of a subgroup $H\subset G(F)$\\
$(g,\varsigma)$ & an element of $\overline{G(F)}$, $g\in G(F)$,
$\varsigma\in\mu_2=\{\pm1\}$\\
$\beta$, $(x,y)_2$ & the Kubota cocycle, and the quadratic Hilbert
symbol\\
$p$, $q$, ${\cal O}$, $\varpi$ & residual characteristic ($p=\infty$ if
$F$ is archimedean), order of the residue field, ring of integers, and a
uniformizer of $F$ \\
$\nu$, $|\cdot|$ & normalized valuation and absolute value,
$|\varpi|=q^{-1}$\\
$B=AN$ & the upper triangular Borel subgroup, its diagonal torus $A$ and
unipotent radical $N$\\
$a(x)$, $h(x)$ & the diagonal element ${\rm diag}(x,x^{-1})$\\
$T$, ${\cal T}$ & a Cartan subgroup of $G$; a set of representatives of
the $G(F)$-conjugacy classes of Cartan subgroups\\
$K_0$, $K_1$ & the hyperspecial maximal compact subgroups $G({\cal O})$
and $w^{-1}G({\cal O})w$ ($p$ odd)\\
$\pi$, $\overline{\pi}$ & a representation of $G$, and a genuine
representation of the cover\\
$\Theta_\pi$, $\tr\pi(f)$ & the character of $\pi$, and the trace of
$\pi(f)$\\
$\chi$, $\overline{\chi}$ & a character of $A(F)$, and a genuine
character of $\overline{A(F)}$\\
$O_f$, $F_f$, $\Phi_f$ & orbital integral, normalized hyperbolic orbital
integral, normalized elliptic orbital integral of $f$\\
$S$ & a finite set of places of $k$ (only in Chapters~\ref{ch:tf}
and~\ref{ch:global})\\
$S(\pi)$ & the exceptional set of places of
Definition~\ref{def:exceptional}\\
${\cal C}_v$ & a class of corresponding pairs at a place $v\in S$\\
$t_v$ & the local transfer of genuine representations in
Hypothesis~\ref{hyp:LH}\\
$\nu_\infty(\pi)$ & the archimedean infinitesimal character of $\pi$
(\S\ref{sec:hypotheses})\\
\hline
\end{tabular}
\end{center}

\section*{Acknowledgments}
\addcontentsline{toc}{section}{Acknowledgments}

I thank C. Moen and J. Adams for many helpful conversations,
and some anonymous referees of the 1995 draft for their comments
(one of whom prodded me to treat the case $p=2$ separately).

This 2026 revision was prepared with the help of AI systems.
Carbon (Google Gemini) refereed the 1995 draft. Claude (Anthropic)
assessed those referee reports, merged the original source files into a
single file, made the corrections, and rewrote the Introduction,
\S\ref{sec:dyadic}, and Chapter~\ref{ch:global} so that the main result
is stated conditionally (Draft~1). The author made the decision to state
the result conditionally. A referee report on Draft~1, supplied by the
author, found further errors, among them a false finiteness lemma in
Chapter~\ref{ch:global}; Claude assessed that report and made the
Draft~2 revisions, including the infinitesimal-character refinement of
Hypothesis~\ref{hyp:GH}, Lemma~\ref{lem:A00}, and the corrections to the
local character formulas in Chapter~\ref{ch:local}. Further details of
the AI usage, the list of changes, and the sources that were not
consulted directly during the revision are recorded in the revision log
(\texttt{revision\_log.md} in the directory \texttt{mult1sl2} of the
repository \texttt{github.com/wdjoyner/trace-formulas}).


%% ===== begin aabook1.tex =====
\chapter{Background on covers of $SL(2)$}\label{ch:background}

\section{Local covers of $SL(2)$}\label{sec:local}

Let $F$ be a local field of characteristic zero and let $\mu_2=\{\pm 1\}$.
Let $p=\infty$ if $F$ is archimedean and otherwise
let $p$ denote the residual characteristic of $F$.
If $F$ is a p-adic field, let ${\cal O}$ denote its ring of integers,
$\varpi$ a uniformizer, $q$ the cardinality of the residue field, and
$\nu =\nu _F=\nu _v:F^{\times }\rightarrow {\mathbb{Z}}$ the
normalized valuation, so that $\nu (\varpi )=1$; we write
$|x|_F=q^{-\nu(x)}$, so that $|\varpi|_F=q^{-1}$.


For $g=\left(
\begin{array}{cc}
a & b \\
c & d
\end{array}
\right) \in SL_2(F)$, let
\[
x(g)=\left\{
\begin{array}{c}
c,\;\;\;c\neq 0, \\
d,\;\;\;c=0,
\end{array}
\right.
\]
and let $\beta =\beta _v:SL_2(F)\times SL_2(F)\rightarrow \mu _2$ be defined
by
\begin{align}
\beta (g_1,g_2) &=(x(g_1),x(g_2))_2\cdot (-x(g_1)^{-1}x(g_2),x(g_1g_2))_2 \label{eq:beta}\\
&=(\frac{x(g_1g_2)}{x(g_1)},\frac{x(g_1g_2)}{x(g_2)})_2,  \nonumber
\end{align}
where $(...,...)_2=(...,...)_{2,v}:F^{\times }
\times F^{\times }\rightarrow \mu _2$
denotes the Hilbert symbol \cite{W1}.
Note that this is trivial if $F={\mathbb {C}}$.

\begin{lemma}\label{lem:kubota}
(Kubota \cite{K1}, Moore \cite{M})
The map $\beta $ defines a continuous Borel
2-cocycle (giving $\mu _2$ the discrete topology of course). Its cohomology
class in the continuous cohomology group $H^2(SL_2(F),\mu _2)$ corresponds
to a 2-fold topological central extension $\overline{SL_2(F)}$ which fits
into the short exact sequence
\[
1\rightarrow \mu _2\rightarrow \overline{SL_2(F)}\stackrel{p}{\rightarrow }%
SL_2(F)\rightarrow 1.
\]
\end{lemma}

If $H<SL_2(F)$ is a subgroup then we define $\overline{H}=p^{-1}(H)\subset
\overline{SL_2(F)}$.

If $g\in GL_2(F)$, let
\begin{align*}
p(g) &=\left(
\begin{array}{cc}
1 & 0 \\
0 & \det g
\end{array}
\right) ^{-1}g, \\
d(a,b) &=\left(
\begin{array}{cc}
a & 0 \\
0 & b
\end{array}
\right) , \\
g^x &=x^{-1}gx,
\end{align*}
and define the extension $\beta _1$ of the above cocycle to
$G_1(F)=GL_2(F)$ by
\begin{equation}
\beta _1(g_1,g_2)=\beta (p(g_1)^{d(1,\det g_2)},p(g_2))\cdot v(\det
g_2,p(g_1)),\label{eq:beta1}
\end{equation}
where $v:F^{\times }\times SL_2(F)\rightarrow \mu_2$ is defined by
\[
v(z,g)=\left\{
\begin{array}{c}
1,\;\;\;c\neq 0,\;F\neq {\mathbb R}, \\
(z,d)_2,\;\;\;\;\;\;\;c=0,
\end{array}
\right.
\]
and $v=1$ if $F={\mathbb R}$. This defines a 2-fold metaplectic cover of
$GL_2(F)$.

Elements of $\overline{SL_2(F)}$ will be denoted by $(g,\varsigma )$,
where $g\in SL_2(F)$, $\varsigma \in \mu _2$.
Multiplication on $\overline{SL_2(F)}$ is given by
\begin{equation}
(g_1,\varsigma _1)\cdot (g_2,\varsigma _2)=(g_1g_2,\beta (g_1,g_2)\varsigma
_1\varsigma _2),\label{eq:mult}
\end{equation}
so inverses are given by
\begin{equation}
(g,\varsigma )^{-1}=(g^{-1},\beta (g,g^{-1})^{-1}\varsigma ^{-1})\label{eq:inv}
\end{equation}
and conjugation is given by
\begin{equation}
(g,\varsigma )\cdot (h,\varsigma ^{\prime })\cdot (g,\varsigma
)^{-1}=(ghg^{-1},\beta (gh,g^{-1})\beta (g,h)\beta (g,g^{-1})^{-1}\varsigma
^{\prime }),\label{eq:conj}
\end{equation}
for $g,h,g_1,g_2\in G(F)$.

These identities also hold for $\beta _1$ in place of $\beta$,
for $g,h,g_1,g_2\in G_1(F)$.

We let $B=\{\left(
\begin{array}{cc}
\ast  & * \\
0 & *
\end{array}
\right) \}\subset G=SL_2$ denote the standard Borel, $N$ its unipotent
radical, and $A$ its Levi component. Let
$\theta\in F^\times -F^{\times 2}$ and let
$T=T_\theta\subset G(F)$ denote the non-split Cartan subgroup
\[
T=\{
    \left(
    \begin{array}{cc}
    a&b\theta\\
    b&a
    \end{array}
    \right)
    \ |\ a^2-b^2\theta =1,\ a,b\in F
  \}.
\]
\begin{lemma}\label{lem:Ainv}
For all $g\in A(F)$, 
\[
g=\left(
\begin{array}{cc}
a & 0 \\
0 & a^{-1}
\end{array}
\right) ,
\]
$\varsigma \in \mu _2$, we have $(g,\varsigma
)^{-1}=(g^{-1},(a,a)_2\varsigma ^{-1})$.
\end{lemma}

\begin{lemma}\label{lem:facts}
(a) $\overline{N(F)}=N(F)\times \mu _2$ (direct product),

(b) if $p>2$ then
$\overline{SL_2({\cal O})}\cong SL_2({\cal O})\times
\mu _2$ (direct product but the splitting may be non-trivial),

(c) if $K_0=SL_2({\cal{O}})$ and
\[
K_{m}=\{\left(
\begin{array}{cc}
a & b \\
c & d
\end{array}
\right) \in SL_2({\cal{O}})\;|\;a-1,d-1,b,c\in \varpi ^m{\cal{O}}\},
\;m>0,
\]
then $\overline{K_{m}}=K_{m}\times \mu _2$
(direct product) for all sufficiently large $m$,

(d) if $F={\mathbb {R}}$, and $K=SO(2,{\mathbb {R})}$ then
$\overline{K}$ does not split,

(e) The center of $\overline{SL_2(F)}$ is $Z(\overline{SL_2(F)})=\{\pm
1_2\}\times \mu _2$ (direct product if $(-1,-1)_2=1$).

(f) $\overline{A(F)}$ is abelian.

(g) $\overline{T}$ is abelian.

(h) If $g\in G(F)$ and
$n=
\left(
\begin{array}{cc}
1 & 1\\
0 & 1
\end{array}
\right)
$
then $(g,\zeta)(n,\zeta ')(g,\zeta)^{-1}
=(gng^{-1},\zeta')$.

\end{lemma}

\begin{remark}
These well-known facts
were proven in \cite{K2} or follow from the properties of the Hilbert
symbol \cite{W1}.
For example, (g) is a consequence of a calculation in \cite{Fl}, p. 128.
\end{remark}

\subsection{Local algebras of functions}\label{sec:localalg}

If $H$ is the group of $F$-rational points of a reductive algebraic group
over an archimedean
or p-adic field $F$, let $C_c^\infty (H)$ denote the usual
algebra of compactly supported smooth (if $F$ is p-adic: locally constant)
functions under convolution, and, for an open compact subgroup $K$ of
$H$ ($F$ p-adic), let $C_c^\infty (H//K)$ denote the subalgebra of
functions in $C_c^\infty(H)$ which are bi-$K$-invariant. A
function $f:\overline{G(F)}\rightarrow {\mathbb{C}}$ such that
\begin{equation}
f(g,\varsigma )=\varsigma \cdot f(g,1)\label{eq:genuine}
\end{equation}
will be called \textbf{genuine}. For a function $f$ on
$\overline{G(F)}$ we write $f|_{G(F)}$ for the function $g\mapsto f(g,1)$
on $G(F)$. Let
\begin{align*}
C_c^\infty (\overline{G(F)})_{gen}
&=\{f:\overline{G(F)}\rightarrow {\mathbb{C}}\
|\;(i)\;f|_{G(F)}\in C_c^\infty (G(F)),
\\
&\quad (ii)\;f\;\mathrm{is\;genuine}\}
\end{align*}
denote the Schwartz-Bruhat algebra of smooth genuine functions.
If $F$ is p-adic, $K=G({\cal O})$, and $p>2$, let
$\kappa:K\rightarrow\overline{K}$, $\kappa(k)=(k,s_K(k))$, be a splitting
of $\overline{K}$ over $K$ (Lemma~\ref{lem:facts}(b)), and let
\begin{align*}
{\cal{H}}(\overline{G(F)},K)_{gen}
&=\{f\in C_c^\infty (\overline{G(F)})_{gen}
\ |\ f(\kappa(k)x\kappa(k'))=f(x)\ \text{for all } k,k'\in K,\
x\in\overline{G(F)}\}
\end{align*}
denote the Hecke algebra of genuine spherical functions; it is an
algebra under convolution because its elements have compact support.
We also write ${\cal H}(G(F),K)=C_c^\infty(G(F)//K)$.

\begin{example}
The unit element of ${\cal{H}}(\overline{G(F)},K)_{gen}$ (for the
measure giving $\overline{K}$ volume $1$, as in the next subsection) is
the function $e_K$ with $e_K(k,\varsigma)=\varsigma\, s_K(k)$ for $k\in K$ and
$e_K=0$ off $\overline{K}$. Thus $e_K|_{G(F)}=s_K\cdot 1_K$, which is
the characteristic function of $K$ only when the splitting $\kappa$ is
$k\mapsto(k,1)$. This is why bi-invariance is imposed with respect to
$\kappa(K)$, not by requiring $f|_{G(F)}$ to be bi-$K$-invariant.
\end{example}

\subsection{Local measures}

Let $F$ be a local field and let $\overline{F}$ denote a separable
algebraic closure.
Let $dk$ denote the Haar measure on $K=G({\cal O})$
normalized so that $K$ has
volume 1. Fix an additive non-trivial character $\psi$ of
$F$ and let $dn$ denote a Haar
measure on $N(F)$ which is induced by a self-dual measure on
the additive group $F$ (with respect to $\psi$).
Let $da$ be a Haar measure on $A(F)$. Every Cartan subgroup $T$ of
$G$ defined over $F$ is either $G(F)$-conjugate to $A$, or anisotropic
over $F$; in the second case $T(F)$ is compact, and we let $dt$ be the
Haar measure on $T(F)$ of total volume $1$. (Draft~1 defined $dt$ by
pulling back $da$ along an isomorphism $T\rightarrow A$ defined over
$\overline{F}$; such an isomorphism does not carry $T(F)$ to $A(F)$,
so it does not by itself transport Haar measures between these groups.
Any fixed normalization will do for the local comparisons below, provided
the same one is used on $T(F)$ and on $\overline{T(F)}$.)
Let $dg$ be the Haar measure
on $G(F)$ such that $dg=\delta _B(a)\,dk\,da\,dn$, $g=kan$, where
\begin{equation}
\delta _B(a)=|t|^2,\;\;a=\left(
\begin{array}{cc}
t & 0 \\
0 & t^{-1}
\end{array}
\right)
\end{equation}
is the modulus of $B$. (Draft~1 had $\delta_B(a)^{1/2}$ here.)

We extend these measures to the 2-fold cover by imposing the discrete
measure $d\varsigma $ on $\mu _2$ which gives $\mu _2$ total volume 1. For
example, the measure on $\overline{G(F)}$ will be denoted
$d\overline{g}=dg\cdot d\varsigma $.

\section{Global aspects}\label{sec:global}

Let $k$ denote a number field. Let $\Sigma _k$ denote the set of places of
$k$, $\Sigma _k^\infty $ denote the set of archimedean places of
$k$ (which we may identify with the set of inequivalent
complex embeddings $k\hookrightarrow {\mathbb{C}}$), and let
$\Sigma _k^f$ denote the set of non-archimedean places of
$k$ (which we may identify with the set of non-zero prime ideals of
${\cal{O}}_k$). Let ${\mathbb{A}}={\mathbb{A}}_k$ denote the adele
ring of $k$ and, for $v\in \Sigma _k$, let $k_v$ denote
the completion of $k$ with respect to $v$.

Let $G=SL(2)$.

The global cover is constructed from the local ones as follows
\cite{K2}. For each odd finite place $v$ let
$\kappa_v:K_{0,v}\rightarrow\overline{SL_2(k_v)}$ be the splitting of
Lemma~\ref{lem:facts}(b); for almost all $v$ we may take $\kappa_v(k)=(k,s_v(k))$, with
$s_v$ as in (\ref{eq:sv}) below. Let
$\prod'_v\overline{SL_2(k_v)}$ denote the restricted direct product of the
local covers with respect to the subgroups $\kappa_v(K_{0,v})$, and let
\[
Z_0=\Bigl\{(\varsigma_v)_v\in\bigoplus_{v\in\Sigma_k}\mu_2\ \Bigm|\
\prod_v\varsigma_v=1\Bigr\},
\]
a discrete central subgroup (only finitely many $\varsigma_v$ are
$\neq 1$). Then
\[
\overline{G({\mathbb{A}})}=\Bigl({\prod_v}'\,\overline{SL_2(k_v)}\Bigr)\Big/Z_0 ,
\]
and the product map $\mu_2^{(\Sigma_k)}\rightarrow\mu_2$ identifies the
image of $\bigoplus_v\mu_2$ with a central subgroup $\mu_2$, giving the
exact sequence
\begin{equation}
1\rightarrow \mu _2\rightarrow
\overline{G({\mathbb{A}})}\stackrel{p}{\rightarrow }G({\mathbb{A}})\rightarrow 1.
\end{equation}
Thus $\overline{G({\mathbb{A}})}$ is a \emph{quotient} of the restricted
product of the local covers, not a subgroup of it.

Equivalently, $\overline{G({\mathbb{A}})}$ is the set $G({\mathbb{A}})\times\mu_2$
with the multiplication (\ref{eq:mult}) for the 2-cocycle
\[
\beta_{\mathbb{A}}(g_1,g_2)=\prod_{v\in \Sigma_k}\beta'_v(g_{1,v},g_{2,v}),\qquad
\beta'_v(g_1,g_2)=\beta_v(g_1,g_2)\,s_v(g_1)\,s_v(g_2)\,s_v(g_1g_2)^{-1}
\]
(with $s_v$ as in (\ref{eq:sv}); $s_v=1$ at the archimedean places). The cocycle $\beta'_v$
is cohomologous to $\beta_v$ and is trivial on $K_{0,v}\times K_{0,v}$ for
almost all $v$ \cite{K2}, so for $g_1,g_2\in G({\mathbb{A}})$ the product
has only finitely many factors different from $1$. (Draft~1 wrote the
product of the unmodified $\beta_v$, which is not a finite product in
general.) The rules of multiplication, inverses, and conjugation are as in
(\ref{eq:mult})--(\ref{eq:conj}), with $\beta_{\mathbb{A}}$ in place of
$\beta$. If $H<G({\mathbb{A}})$ is a subgroup then we define
$\overline{H}=p^{-1}(H)<\overline{G({\mathbb{A}})}$.

\begin{example}
Let $v_1\neq v_2$ be places of $k$. The element of the restricted
product with $\varsigma_{v_1}=\varsigma_{v_2}=-1$ and all other
components trivial lies in $Z_0$, so it is the identity of
$\overline{G({\mathbb{A}})}$; the element with $-1$ at $v_1$ only is the
nontrivial central element of $\mu_2\subset\overline{G({\mathbb{A}})}$.
\end{example}

\begin{lemma}\label{lem:Gk-split}
$\overline{G(k)}\cong G(k)\times \mu _2$,
where the splitting isomorphism in $\cong$ will be
described below.
\end{lemma}

\begin{proof} This is well-known \cite{K2}.
\end{proof}

From \cite{K2}, we see that
the splitting of the above exact sequence indicated in the above
lemma is not given by $\gamma \longmapsto (\gamma ,1)$ but by $\gamma
\longmapsto (\gamma ,s_{{\mathbb{A}}}(\gamma )^{-1})$, where
\[
s_{{\mathbb{A}}}(\gamma )=
\prod_{v\in \Sigma _k}
s_v(\gamma),
\]
and
\begin{equation}\label{eq:sv}
s_v(\gamma)=\left\{
\begin{array}{c}
(c,d)_{v,2}^{-1},
\;\;\;cd\neq 0{\rm \ and\ }
2\;\mathrm{does\;not\;divide}\;v(c)
{\rm \ and\ }v\notin \Sigma_k^\infty, \\
1,\;\;\;\;\;\;\;\;\;\;\;\;\;\;\;\;
\mathrm{\;otherwise}.\;\;\;\;\;\;\;\;
\end{array}
\right.
\end{equation}
Here $\gamma =\left[
\begin{array}{cc}
a & b \\
c & d
\end{array}
\right] \in G(k)$.

{\em From now on, we identify $G(k)$ with its image
in $\overline{G({\mathbb{A}})}$ under this splitting (and the diagonal
embedding).} We identify $N({\mathbb{A}})$ with the
subgroup $N({\mathbb{A}})\times 1\subset \overline{G({\mathbb{A}})}$.

Let $K_0=\Pi_{v\in \Sigma_k}K_{0,v}$
denote the ``usual'' maximal compact subgroup of $G({\mathbb{A}})$.
Here, if $v\in \Sigma_k^f $ then $K_{0,v}=G({\cal O}_v)$ and
otherwise
$K_{0,\infty}=\prod_{\substack{ v|\infty\\ v\in\Sigma_k}}
K_v$, where $K_v=SO(2,{\mathbb{R}})$ if $v$ is real and
$K_v=SU(2,{\mathbb{C}})$ if $v$ is complex.

\subsection{Global algebras of functions}\label{sec:globalalg}

A function $f:\overline{G({\mathbb{A}})}\rightarrow {\mathbb{C}}$
such that (\ref{eq:genuine})
holds will be called \textbf{genuine}. Let
${\cal{H}}(\overline{G({\mathbb{A}})})$ denote the algebra of
functions generated by the elements of the
restricted tensor product
$\bigotimes_{v\in \Sigma_k }
C_c^\infty (\overline{G(k_v)})_{gen}$,
where almost all $f_v$ in the product
$f=\otimes _{v\in \Sigma_k}f_v\in
\bigotimes_{v\in \Sigma_k }C_c^\infty (\overline{G(k_v)})_{gen}$ are
spherical in the sense of \S\ref{sec:localalg}.

\subsection{Global measures}

Let $dk$ denote the Haar measure on $K=\prod_{v\in \Sigma_k }K_v$
defined by $dk=\prod_{v\in \Sigma_k }dk_v$.
Let $dg$ be the Tamagawa measure on
$G({\mathbb{A}})$ \cite{W2}. For any Cartan subgroup
$T$ of $G$ defined over $k$, let $dt$
denote the global Tamagawa measure on
$T({\mathbb{A}})$. If $T=A$ then this will
be denoted by $da$. Let $dn$ denote the product measure on $N({\mathbb{A}})$.

We extend these measures to the 2-fold cover by imposing the discrete
measure $d\varsigma $ on $\mu _2$ which gives $\mu _2$ total volume 1. For
example, the measure on $\overline{G({\mathbb{A}})}$ will be denoted
$d\overline{g}=dg\cdot d\varsigma $.
%% ===== end aabook1.tex =====



%% ===== begin aabook3.tex =====
\chapter{Representations of $\overline{SL(2,F)}$}\label{ch:reps}

This chapter presents a survey of certain aspects of the
representation theory of central covering groups of p-adic reductive groups
in the case of $SL(2)$.

\section{Characters}

Characters occur naturally in the trace formula, so we recall
some of their properties.

\subsection{The character as a function}

The existence of the character as a function has been used many times in
the literature but no proof, as far as I know, has been published.
The argument in \cite{S} extends to the metaplectic case with only
minor modification.

Let $\overline{G(F)}^{\prime }=\{(g,\varsigma )\in \overline{G(F)}%
\;|\;g\;{\rm has\;distinct\;eigenvalues}\},$ called the {\bf regular set} of
$\overline{G(F)}.$ In particular, every element in the regular set is
semi-simple (i.e., diagonalizable). The following result is a corollary of a
result of Howe and Harish-Chandra \cite{S}:

\begin{lemma}
There exists a genuine locally constant function denoted $\Theta _\pi $ on
the regular set $\overline{G(F)}^{\prime }$ which represents the trace of $%
\pi :$%
\[
\tr(\pi (f))=\int_{\overline{G(F)}}\Theta _\pi (x)f(x)dx,
\]
for all $f\in C_c^\infty (\overline{G(F)})_{gen}.$
\end{lemma}

This trace, which we sometimes denote simply by $\Theta _\pi (f)$, is called
the {\bf Fourier transform} of $f$ at $\pi $. The function $\Theta _\pi (x)$
is called the {\bf character} of $\pi $.

\subsection{Principal series representations}\label{sec:ps-char}

Let ${\cal T}$ denote a complete set of representatives of conjugacy classes of
Cartan subgroups of $G(F)$. We may and do assume that $A\in {\cal T}$. For each
$T\in {\cal T}$, let $W_T=N_{G(F)}(T)/T$ denote the Weyl
group of $T$. If $t\in T$ is regular then
\[
T\subset G(F)_t=Cent(t,G(F))
\]
is finite index.

\begin{lemma}
Assume $\chi\in \Pi (\overline{A(F)})$ and suppose
that $\chi \neq \chi ^w$. The character of the
irreducible representation $\pi =I_{A}(\chi )$\ is given on the regular
set by
\[
\Theta _\pi (x)=\left\{
\begin{array}{c}
\frac{\chi (a)+\chi (a^w)}{|a_1-a_2|},
\;\;x=a^g,\;{\rm for\;some}\;g\in \overline{G(F)},
\;a\in \overline{A(F)}, \\
0,\;\;\;\;\;\;\;\;\;\;\;\;\;\;\;\;\;\;\;\;\;\;\;\;\;\;\;\;\;
{\rm \;if\;}\;\;\;
x\in \overline{G(F)}^{\prime }-\overline{A(F)}^{\overline{G(F)}},
\end{array}
\right.
\]
where
\[
\overline{A(F)}^{\overline{G(F)}}=\{a^g\;|\;a\in \overline{A(F)},\;g\in
\overline{G(F)}\}.\
\]
\end{lemma}

The proof, which is omitted, follows the proof of \cite{Ge1}, Theorem 5.19.

\begin{remark} As noted above, on the regular set,
this character vanishes off the conjugacy class of $\overline{A(F)}$.
\end{remark}

Now we describe the character of a principal series representation
as the Fourier transform of the function $a\longmapsto J_G(a,f)$.

Let $\chi \in \Pi (\overline{A(F)})$,
$\pi =I_{A}(\chi )\in \Pi (\overline{G(F)})$.
This lemma and the Weyl integration formula
imply that, for $f\in C_c^\infty
(\overline{G(F)})$,
\begin{equation}
\begin{array}{c}
\Theta _\pi (f)=\int_{\overline{G(F)}}\Theta _\pi (x)f(x)dx \\
=\frac 12\int_{\overline{A(F)}}(\chi (a)+
\chi(a^w))F_f^{\overline{A(F)}}(a)da \\
=(F_f^{\overline{A(F)}})^{\wedge }(\chi ) ,
\end{array}
\end{equation}
where the Fourier transform is defined by
\begin{equation}
\begin{array}{c}
\varphi ^{\wedge }(\chi )=\int_{\overline{A(F)}}\chi (x)\varphi
(x)da \\
\ \ \ \ \ \ \chi \in \Pi (\overline{A(F)}).
\end{array}
\end{equation}
The last equality in the first display holds because $F_f^{\overline{A(F)}}$ is
invariant under $a\mapsto a^w$, since $a$ and $a^w$ are conjugate in
$\overline{G(F)}$. The same formulas, without bars, hold for $G(F)$.

\begin{lemma}
If $F(\chi)=\tr I_A(\chi)(f)$ then
\[
F^\vee(\gamma)
=|a_1-a_2|J_G(\gamma,f),
\]
for $\gamma =(a,\zeta)\in \overline{A(F)}$, $a\not= \pm 1$.
\end{lemma}

\begin{proof} This follows from our results above.
\end{proof}

\section{Description of the dual spaces}

The dual of $\overline{G(F)}$ is classified in \cite{Wa2}.
We shall recall some details from that paper and the
papers of Moen and Ariturk.

\subsection{Principal series, reducible principal series}

Let $\chi \in \Pi (\overline{A(F)})$. The induced representation
$I_{A}(\chi )$ is in general not irreducible. However, we do have the
following

\begin{theorem}
(Moen \cite{Mo2}) If $p>2$ then $I_{A}(\chi )$ is
irreducible for all $\chi \in \Pi _u(\overline{A(F)})$.
\end{theorem}

In fact, C. Moen \cite{Mo1}
explicitly computes the intertwining operators as
matrices using the Kirillov model.

\subsection{Complementary series}

Let $\mu \in \Pi (\overline{C})$,
$\chi =Ind_{\overline{C}}^{\overline{A(F)}}\mu \in
\Pi (\overline{A(F)})$. If $\mu (x)=\mu _0(x)|x|^s$, for some
character $\mu _0$ of finite order and some $s\in {\mathbb C}$ then
we write $s=s(\mu )=s(\chi )$.

\begin{lemma} If $0<s<1/2$ and $\mu (x)=|x|^s$ then $I_A(\mu )$ is unitary.
\end{lemma}

This follows from \cite{GS}, Proposition 1 following
the argument of \cite{Ar} sketched above. See also
\cite{Mo1} if $p>2$ and \cite{Ge1}, \S 5.4.

\subsection{Special representation, non-tempered representation}

The following result is due to Gelbart and Sally \cite{GS}:

\begin{lemma}
If $s=1/2$ and $\mu (x)=|x|^s$ then there is an irreducible
constituent of $I_A(\mu)$, ``the special
representation'' $\pi _{sp}$, which is elliptic (in the sense that its
character has support on the elliptic set) and square-integrable.
Its complement in $I_A(\mu)$ is an infinite-dimensional,
elliptic, non-tempered representation $\pi _{nt}$.
\end{lemma}

Furthermore, in this case

(1) the representation $\pi _{sp}$ is the only square-integrable
representation in $\Pi _u(\overline{G(F)})$ which is not cuspidal.

(2) $\pi _{sp}$ is tempered.

(3) The representation $\pi _{nt}$ and the complementary series comprise the
only non-tempered representations in $\Pi _u(\overline{G(F)})$.

\subsection{Cuspidal representations}

Thanks to ideas originally introduced by Weil and Shale, much is known about
the cuspidal representations in $\Pi _u(\overline{G(F)})$ in the case
where $p>2$. See for example \cite{M1}, \cite{M2}, \cite{M3}.
The information we shall need will be recalled in Chapter~\ref{ch:sc}.
%% ===== end aabook3.tex =====



%% ===== begin aabook4.tex =====
\chapter{On a supercuspidal correspondence}\label{ch:sc}

\section{Introduction}
In this chapter, we obtain a character relation between
``oddly corresponding'' supercuspidal
representations of $G(F)=SL(2,F)$ and its 2-fold metaplectic cover
$\overline{SL(2,F)}$, where $F$ is a p-adic field with $p>2$.
Following a suggestion of P. Kutzko, we use the Frobenius formula
to obtain character relations between ``corresponding'' supercuspidal
representations of $G(F)$ consistent with the global
correspondence in Chapter~\ref{ch:global}.
The parameterization of the supercuspidals for the correspondence
which we use is that of D. Manderscheid \cite{M1}-\cite{M2}, recalled
below (see also \cite{M3}).

\section{Filtrations and classification}\label{sec:filtrations}

Let
$F$ denote a p-adic field, let
$\varpi =\varpi_F$ be a fixed uniformizer of $F$, let $P:=P_F:=\varpi_F{\cal{O}}_F$,
and ${\mathbb{F}} :={\cal{O}}_F/P_F$.
Let ${\cal{B}}$ denote the tree for $G(F)$, $x_0$ the vertex fixed by
$K_0:=G({\cal{O}}_F)$ and $x_1$ the vertex fixed by $
K_1:=w^{-1}G({\cal{O}}_F)w$, where
$w:=
\left(\begin{array}{cc}
0&\varpi \\ 1& 0\end{array}\right)$.
We may find a metric $d$ on the ambient affine space of ${\cal{B}}$
for which the subgroup of $K_0$ fixing the ball
\begin{equation}
{\cal{B}}_m(x_0):=\{\text{ vertices }x\ |\ d(x,x_0)\leq m\}
\end{equation}
pointwise, $Fix\ {\cal{B}}_m(x_0)$, is the standard filtration subgroup:
\begin{equation}
Fix\ {\cal{B}}_m(x_0)=K_{0,m}:=
\{\left(\begin{array}{cc} a&b \\ c& d\end{array}\right)\in K_0\ |\
a-1\equiv d-1\equiv b\equiv c\equiv 0\  ({\rm mod\ } P^m)\}.
\end{equation}
Similarly,
\begin{equation}
Fix\ {\cal{B}}_m(x_1)=K_{1,m}=w^{-1}K_{0,m}w.
\end{equation}
For $i=0,1$, we call an irreducible representation $\sigma$ of
$K_i/K_{i,m}\cong G({\cal{O}}_F/\varpi^m{\cal{O}}_F)$
{\bf strongly cuspidal of level $m>1$} if
$\sigma$ does not contain the trivial character
of the subgroup $(N\cap K_{i,m-1-i})/K_{i,m}$. In case $m=1$ we call a
{\bf strongly cuspidal representation} of
$K_i/K_{i,1}\cong G({\mathbb{F}})$
any irreducible representation which is cuspidal as a representation of
$G({\mathbb{F}})$. Since $p> 2$, the metaplectic group
$\overline{SL(2,F)}$ splits over $K_0$ and $K_1$ (though the splittings
are different), so a similar definition may be given for
{\bf strongly cuspidal representation} of
$\overline{K_i}/\overline{K_{i,m}}$. The papers \cite{M1},
\cite{M2} show that each irreducible strongly cuspidal representation
of $K_i/K_{i,m}$ ($i=0,1;m\geq 1$) compactly induces to a irreducible
supercuspidal representation of $G(F)$, and similarly for its metaplectic
covers. Moreover, all irreducible supercuspidal representations
of $G(F)$ (or of $\overline{G(F)}$) arise in this way.

Let $pr$ denote the natural projection $pr:\overline{G(F)}\rightarrow G(F)$.

\begin{claim}[1995 draft; needs correction, see Remark~\ref{rem:sc-caveat}] Fix a strongly cuspidal of level $m\geq 1$
representation $\sigma$ of $K_i$, $i=0,1$. Let
$\overline{\sigma}:=\sigma \otimes \varepsilon$ on
$\overline{K_i}\cong K_i\times \mu_2$, where $\varepsilon$ is the
identity character of $\mu_2$, and let
$$
\begin{array}{c}
\overline{\pi}
:=\text{Ind }_{\overline{K_i}}^{\overline{G(F)}}\overline{\sigma},\\
\pi
:=\text{Ind }_{K_i}^{G(F)}\sigma.
\end{array}
$$
We have
\begin{equation}
\chi_{\overline{\pi}}(x)=\zeta\cdot\chi_{\pi}(pr(x)),\label{eq:sc-relation}
\end{equation}
if $x=(pr(x),\zeta )\in \overline{G(F)}$ with $pr(x)$
elliptic regular.
\label{prop:164}
\end{claim}

\begin{claim}[1995 draft; depends on Claim~\ref{prop:164}] Suppose that $p>2$. The character relation above
determines a bijective correspondence
$\overline{\pi}\longmapsto \pi$
between the equivalence classes of irreducible genuine supercuspidal
representations of $\overline{SL(2,F)}$, and the equivalence
classes of irreducible supercuspidal representations of $SL(2,F)$.
\end{claim}

\begin{remark}\label{rem:sc-caveat}
(Added 2026.) Two points in Claim~\ref{prop:164} and its proof need
attention. First, $\overline{\sigma}$ must be genuine; in the 1995 draft
it was written $\sigma\otimes 1$, which is not genuine, and we have
replaced it by $\sigma\otimes\varepsilon$, where $\varepsilon$ is the
identity character of $\mu_2$. Second, the isomorphism
$\overline{K_i}\cong K_i\times\mu_2$ is given by a splitting
$\kappa_i(k)=(k,s_i(k))$ which, by Lemma~\ref{lem:facts}(b), need not be
$k\mapsto(k,1)$. With respect to it,
$\overline{\sigma}((k,\zeta))=\zeta\, s_i(k)\,\sigma(k)$, whereas the proof
below evaluates $\chi_{\overline{\sigma}}^\circ$ as if $s_i$ were
trivial. Relatedly, the function $(x,\zeta)\mapsto\zeta\,\chi_\pi(x)$ in
(\ref{eq:sc-relation}) is not obviously invariant under conjugation in
$\overline{G(F)}$. A corrected statement presumably involves the factor
$s_i$ (compare the function $s_v$ of \S\ref{sec:global}, which is
nontrivial on some elliptic elements of $K_0$ when the lower left entry
has odd valuation). Until this is checked, the two claims are not used
in Chapter~\ref{ch:global}. Note also that over a p-adic field the
regular elliptic elements lie in several conjugacy classes of anisotropic
Cartan subgroups (\S\ref{sec:elliptic}); a corrected relation must be
checked on each of them.
\end{remark}

The 1995 argument for Claim~\ref{prop:164} is given in the next subsection,
with the gaps noted in Remark~\ref{rem:sc-caveat}.

\section{The Frobenius formula}

First we prove the following metaplectic version of the Frobenius formula.

\begin{proposition} If $\overline{\pi} \in \Pi_{gen}(\overline{G(F)})$
is induced from a finite dimensional representation $\overline{\sigma}$
of $\overline{K_i}$ then
$$
\chi_{\overline{\pi}} (x)=\sum_{\substack{ z \text{ in }\\
\overline{K_i}\setminus\overline{G(F)}/\overline{K_i}}}
[\sum_{\substack{ g \text{ in }\\
\overline{K_i}\setminus\overline{K_i}z\overline{K_i}}}
\chi_{\overline{\sigma}}^\circ (gxg^{-1})]
$$
converges for each $x\in \overline{G(F)}^{reg}$. (The convergence is not
uniform.)
\label{prop:166}
\end{proposition}

Let me first prove a preliminary lemma.
\begin{lemma} Let $x\in \overline{G(F)}^{reg}$ and let
$$
\chi_{\overline{\sigma}}^z(x):=\sum_{\substack{ g \text{ in }\\
\overline{K_i}\setminus\overline{K_i}z\overline{K_i}}}
\chi_{\overline{\sigma}}^\circ (gxg^{-1}).
$$
There exists a finite subset $Z$ of representatives $z$ of
$\overline{K_i}\setminus\overline{G(F)}/\overline{K_i}$
and an open neighborhood $U$ of $x$ such that
$\chi_{\overline{\sigma}}^z(y)=0$ unless $y\in U$ and $z\in Z$.
\end{lemma}
\begin{proof} Note $\chi_{\overline{\sigma}}^\circ$ is supported
on $\overline{K_i}$ and $K_i$ is the stabilizer in $G(F)$ of the vertex
$x_i$ of the tree ${\cal{B}}$.
It follows that $\chi_{\overline{\sigma}}^\circ (gxg^{-1})\not= 0$
if and only if $gxg^{-1}\in \overline{K_i}$ if and only if
$p(x)$ fixes the point $p(g)^{-1}x_i$ in the tree for $G(F)$.
It follows that
$\chi_{\overline{\sigma}}^z (x)= 0$ unless $p(x)$ is elliptic regular
in $G(F)$. Let $\Delta (x)=q^{-m}$ for $m\geq 0$. It is known that
(\cite{L4}, \S 5)
\begin{equation}
|\text{Fix}(x)|=
\left\{ \begin{array}{c}
\frac{q^{m+1}+q^m-2}{q-1},\ \ \ \ \text{if $F[x]/F$ is unramified},\\
\frac{2(q^m-1)}{q-1},\ \ \ \text{if $F[x]/F$ is ramified}.
\end{array} \right.
\end{equation}
From this formula, it follows that there is a neighborhood $U'$
of $p(x)$ in $G(F)^{reg}$ and a number $N=N(U')<\infty$
such that, for all $y\in U'$, we have $|\text{Fix}(y)|\leq N$.
In other words, the fixed point formula above implies that the
number of fixed points is locally uniformly bounded.

From this we obtain a finite set $Z$ as desired.
\end{proof}

\begin{proof}[Proof of Proposition~\ref{prop:166}]
The proposition follows immediately from
the lemma exactly as in \cite[p.~213]{Kut}.
\end{proof}

\begin{proof}[1995 argument for Claim~\ref{prop:164}]
The idea is as follows:
We compare the character of the supercuspidal representations
$\overline{\pi}\in \Pi_{gen}(\overline{G(F)})$ and
$\pi\in\Pi(G(F))$ corresponding to a given finite dimensional representation
$\sigma$ of $K_i$, $i=0,1$, using the Frobenius formula.
As the number of terms in the sum for $\chi_\pi$ and in
the sum for $\chi_{\overline{\pi}}$ are the same, by
Proposition \ref{prop:166} it suffices to compare the
characters $\chi_\sigma^\circ$
and $\chi_{\overline{\sigma}}^\circ$ term-by-term.
For $x$ in $\overline{K_0}$,
we want to show that two analogous terms are equal by
conjugating the elliptic element $x$ around by an element in the cover of
$K_0$ or $A(F)\cup \{t,t^{-1}\}$ or $N(F)$ (in that order), where
$t={\rm diag}(1,\varpi)$ is a diagonal matrix in $GL(2,F)$.
Note we may pass between $K_0$ and $K_1$ by conjugating by either
$t$ or its inverse.

If $T:=$ Cent$(x,\overline{G(F)})$ then $T$ is elliptic and we may embed $T$ into a maximal
compact subgroup $\overline{K}$ of $\overline{G(F)}$.
To effect the term-by-term comparison, we recall that $\overline{K}_0$
splits, so we must first worry about conjugation by
an element $(r,\zeta)$ with $r$ in $A(F)\cup \{t,t^{-1}\}$.
We first calculate
$$
(r,\zeta)\cdot (k,\zeta')\cdot ( r,\zeta)^{-1}
=(r,\zeta)\cdot (k,\zeta')\cdot (r^{-1},(r_1,r_2)_2\zeta^{-1} ),
$$
for $k\in K_0$ elliptic
and $r=\left(\begin{array}{cc} r_1&0\\0&r_2\end{array}\right)\in
GL(2,F)$. The details of this straightforward calculation are
included for the
reader's convenience. Writing
$k=\left(\begin{array}{cc} a&b\\c&d\end{array}\right)$ and
recalling the definition of $p:GL(2)\rightarrow SL(2)$ from \S\ref{sec:local}, we have
\begin{equation}
\begin{array}{c}
rkr^{-1}=\left(\begin{array}{cc} r_1a&r_1b\\r_2c&r_2d\end{array}\right)\cdot
\left(\begin{array}{cc} r_1^{-1}&0\\0&r_2^{-1}\end{array}\right)
=\left(\begin{array}{cc} a&r_1r_2^{-1}b\\r_1^{-1}r_2c&d\end{array}\right)\\
p(r)k=
\left(\begin{array}{cc} r_1a&r_1b\\r_1^{-1}c&r_1^{-1}d\end{array}\right)\\
p(rk)=
\left(\begin{array}{cc} r_1a&r_1b\\r_1^{-1}c&r_1^{-1}d\end{array}\right)\\
p(rk)^{d(1,r_1^{-1}r_2^{-1})}=
\left(\begin{array}{cc} r_1a&r_2^{-1}b\\r_2c&r_1^{-1}d\end{array}\right)
\end{array}
\end{equation}
If $k$ is elliptic then $c\not= 0$. We have, by (\ref{eq:mult}), 
$$
(r,\zeta)\cdot (k,\zeta')
=(rk,\beta (r,k)\zeta\zeta' )=(rk,(r_1,c)_2\zeta\zeta' ),
$$
so
$$
(r,\zeta)\cdot (k,\zeta')\cdot ( r,\zeta)^{-1}
=(rkr^{-1}, (r_1,c)_2(r_1,r_2)_2(r_1,r_1)_2^{-2}\beta (rk, r^{-1})\zeta' ).
$$
Since
$\beta (rk, r^{-1})=(r_1^{-1},r_1^{-2}r_2c)_2(r_1^{-1}r_2^{-1},r_1)_2$,
we find that
$$
(r,\zeta)\cdot (k,\zeta')\cdot ( r,\zeta)^{-1}
=(rkr^{-1}, (r_1,r_1r_2)_2\zeta' ).
$$
Since $r\in A(F)\cup \{t,t^{-1}\}$, we obtain the desired comparison
$$
\chi_{\overline{\sigma}}^\circ ((r,\zeta)
\cdot (k,\zeta')\cdot ( r,\zeta)^{-1})
=\zeta'\cdot \chi_{\sigma}^\circ (rkr^{-1}),
$$
in this case.
\par
Next we must conjugate our elliptic regular $x=(k,\zeta')$ by
$(n,\zeta)$, where $n\in N(F)$. We obtain, as the reader may verify,
$(n,\zeta)\cdot (k,\zeta')\cdot (n,\zeta)^{-1}=(nkn^{-1},\zeta')$.
\end{proof}
%% ===== end aabook4.tex =====



%% ===== begin aabook2a.tex =====
\chapter{Local comparisons on the groups}\label{ch:local}

\section{Introduction}
We prove the identities necessary to compare the
invariant trace formula to its metaplectic analog in Chapter~\ref{ch:global}.

We shall describe, in the real case, several
correspondences between functions in
$C_c^\infty(G({\mathbb{R}}))$ and functions in
$C_c^\infty(\overline{G({\mathbb{R}})})$, each of which is associated
to a matching of characters. See \S\ref{sec:realA} and
\S\ref{sec:realds}.

We shall describe, in the p-adic case with $p\not= 2$,
five correspondences between functions in
$C_c^\infty(G(F))$ and functions in
$C_c^\infty(\overline{G(F)})$, each of which is associated
to a matching of characters. The case $p=2$ is discussed in
\S\ref{sec:dyadic}, where no matching is established in this version.

\section{Orbital integrals - basics}

Let $G(F)_\gamma$ denote the connected component of the centralizer of
$\gamma$ in $G(F)$.

\subsection{Orbital integrals: definitions}\label{sec:weighted-defs}

Certain weighted orbital integrals occur naturally in the invariant trace
formula. We will need some of their properties to put the trace formula in
stable form as well as to compare trace formulas.

If $M\in \{A,G\}$, $f\in C_c^\infty (\overline{G(F)})_{gen}$, and $g\in
\overline{M(F)}$ is regular, let
\[
J_M(g,f)=|D(g)|^{1/2}
\int_{\overline{G(F)}_g\backslash \overline{G(F)}}
f(x^{-1}gx)v_M(x)dx.
\]
If $F$ is a local field then
\[
v_M(x)=\left\{
\begin{array}{ll}
1, & M=G, \\
\log (1+|n|^2)^{1/2}, & M=A,\ F\ {\rm archimedean},\\
\log \max(1,|n|_F), & M=A,\ F\ {\rm non\hbox{-}archimedean},
\end{array}
\right.
\]
where $x=a\left(
\begin{array}{cc}
1 & n \\
0 & 1
\end{array}
\right) k$.

Note that the centralizer in $\overline{G(F)}$ of a regular element
in $a\in A(F)$ is $\overline{A(F)}$. We may therefore write,
for a regular $a\in \overline{A(F)}$,
\begin{equation}
J_M(a,f)=
|D(g)|^{1/2}\int_{\overline{A(F)}\backslash \overline{G(F)}}
f(x^{-1}gx)v_M(x)dx.
\end{equation}

\subsection{Elliptic orbital integrals}

We now discuss the elliptic orbital integrals $J_G(\gamma)$,
$\gamma\in \overline{T}_{reg}$, where $\overline{T}_{reg}$
denotes the subset of elements of the form $(x,\zeta)$
where $x\in T$ has distinct eigenvalues and $T$ is an elliptic
Cartan of $SL(2,F)$.

By Lemma~\ref{lem:facts}(g), we have
\[
\overline{G(F)}_\gamma
=\overline{T},\ \ \ \ \ \
\gamma\in \overline{T}_{reg}.
\]
We have
\begin{equation}
\int_{\overline{G(F)}_\gamma\backslash \overline{G(F)}}
f(x^{-1}\gamma x)\frac{dx}{dt}
=\int_{\overline{T}\backslash \overline{G(F)}}
f(x^{-1}\gamma x)\frac{dx}{dt},
\end{equation}
where $\gamma =(g,1)\in \overline{T}_{reg}$.

\subsection{Hyperbolic orbital integrals}

By Lemma~\ref{lem:facts}(f), we have
\[
\overline{G(F)}_\gamma
=\overline{A(F)},\ \ \ \ \ \
\gamma\in \overline{A(F)}_{reg}.
\]
We have
\begin{equation}
\int_{\overline{G(F)}_\gamma\backslash \overline{G(F)}}
f(x^{-1}\gamma x)\frac{dx}{da}
=\int_{\overline{A(F)}\backslash \overline{G(F)}}
f(x^{-1}\gamma x)\frac{dx}{da},
\end{equation}
where $\gamma =(g,1)\in \overline{A(F)}_{reg}$.

\subsection{Unipotent orbital integrals}

Using (\ref{eq:conj}) and Lemma~\ref{lem:facts}(h), a straightforward calculation involving only the definitions
yields the following result:
\[
\overline{G(F)}_\gamma
=\overline{N(F)},\ \ \ \ \ \
\gamma=
(
\pm \left(
\begin{array}{cc}
1 & 1\\
0 & 1
\end{array}
\right)
,1)\in \overline{N(F)}.
\]
We have
\begin{equation}
\int_{\overline{G(F)}_\gamma\backslash \overline{G(F)}}
f(x^{-1}\gamma x)\frac{dx}{dn}
=\int_{\overline{N(F)}\backslash \overline{G(F)}}
f(x^{-1}\gamma x)\frac{dx}{dn},
\end{equation}
where $\gamma $ is as above.

\subsection{Weighted orbital integrals}

If $\gamma =(g,1)\in \overline{A(F)}_{reg}$ and if
\[
g=
\pm \left(
\begin{array}{cc}
a & 0\\
0 & a^{-1}
\end{array}
\right)
\in A(F),
\]
then
\begin{equation}
\int_{\overline{G(F)}_\gamma\backslash \overline{G(F)}}
f(x^{-1}\gamma x)v_A(x)\frac{dx}{da}
=\int_{\overline{A(F)}\backslash \overline{G(F)}}
f(x^{-1}\gamma x)v_A(x)\frac{dx}{da}.
\end{equation}

If
\[
\gamma=
(
\pm \left(
\begin{array}{cc}
1 & 0\\
0 & 1
\end{array}
\right)
,1)\in Z_{\overline{G}}(F),
\]
we have
\begin{equation}
J_A(\gamma,f)=
\lim_{\substack{ a\rightarrow 1\\ a\ {\rm regular}}}
\int_{\overline{A(F)}\backslash \overline{G(F)}}
f(x^{-1}\gamma a x)v_A(x)\frac{dx}{dt}.
\end{equation}

\section{Corresponding functions}

We define the notion of a ``smooth matching''. We shall use separate
definitions for $p=\infty$ and $p>2$; for $p=2$ see
\S\ref{sec:dyadic}.
In these cases we shall base our ``matching'' of functions (or, more
properly, of orbital integrals) on our knowledge of the characters of the
representations. For the purpose of avoiding confusion in later
chapters, we shall assume $p>2$ in some sections even though
that hypothesis will not be used until later.

\subsection{Corresponding functions on $A$}

We concentrate first on the set $A(F)_{reg}$.

\subsubsection{Characters of $\overline{A(F)}$}\label{sec:charsA}

We assume $p>2$ throughout this subsection. (The case $p=2$ is
not covered; see \S\ref{sec:dyadic}.) We write $a(x)={\rm diag}(x,x^{-1})$
and, for $a\in A(F)$, $x_a$ for the element with $a=a(x_a)$. We write
$(x,y)$ for the Hilbert symbol $(x,y)_2$; its values are $\pm1$, so
$(x,y)^{-1}=(x,y)=(y,x)$.

\begin{lemma}\label{lem:Acocycle}
(a) For $x,y\in F^\times$, $\beta(a(x),a(y))=(x,y)$. Hence
\[
(a(x),\zeta)(a(y),\zeta')=(a(xy),(x,y)\zeta\zeta').
\]

(b) Let $w=\left(\begin{smallmatrix}0&-1\\1&0\end{smallmatrix}\right)$.
For every $a\in A(F)$, $(w,1)(a,1)(w,1)^{-1}=(a^{-1},1)$.
(This holds for every local field $F$.)
\end{lemma}

\begin{proof}
(a) Since $c=0$, $x(a(x))=x^{-1}$, so by (\ref{eq:beta})
$\beta(a(x),a(y))=((xy)^{-1}/x^{-1},(xy)^{-1}/y^{-1})=(y^{-1},x^{-1})=(x,y)$.

(b) Let $a=a(x)$. We have $x(w)=1$, $x(w^{-1})=-1$, $x(a)=x^{-1}$,
$wa=\left(\begin{smallmatrix}0&-x^{-1}\\x&0\end{smallmatrix}\right)$ so
$x(wa)=x$, and $waw^{-1}=a(x^{-1})$ so $x(waw^{-1})=x$. Hence
$\beta(w,a)=(x,x^2)=1$, $\beta(w,w^{-1})=(1,-1)=1$ and
$\beta(wa,w^{-1})=(1,-x)=1$. By (\ref{eq:mult}) and (\ref{eq:inv}),
$(w,1)^{-1}=(w^{-1},1)$ and
$(w,1)(a,1)(w^{-1},1)=(wa,1)(w^{-1},1)=(a^{-1},1)$.
\end{proof}

Let
\[
A_2(F)=\{a(x)\ |\ \nu (x)\ {\rm \ even}\}.
\]
If $\nu(x)$ and $\nu(y)$ are even then $(x,y)=1$ (as $p$ is odd), so by
Lemma~\ref{lem:Acocycle}(a), $\overline{A_2(F)}=A_2(F)\times \mu_2$
(direct product), and every genuine character of $\overline{A_2(F)}$
has the form $(a(x),\zeta)\mapsto\zeta\,\chi_1(x)$ for a character
$\chi_1$ of $A_2(F)$ (which we regard as a character of the group of
$x$ with $\nu(x)$ even). Every element of $A(F)-A_2(F)$ is $a(x)$ with
$x/\varpi$ of even valuation.

\begin{lemma}\label{lem:genuine-chars}
Let $\chi_1\in\Pi(A_2(F))$ and let $\epsilon\in{\mathbb C}$ satisfy
\begin{equation}\label{eq:eps}
\epsilon^2=(\varpi,\varpi)^{-1}\chi_1(\varpi^2).
\end{equation}
Then
\begin{equation}\label{eq:chi-eps}
\overline{\chi}_\epsilon (a(x),\zeta)=
\left\{
\begin{array}{ll}
\zeta^{-1}\chi_1(x), & \nu(x)\ \text{even},\\
\epsilon\,\zeta^{-1}(x,\varpi)^{-1}\chi_1(x/\varpi),
& \nu(x)\ \text{odd},
\end{array}
\right.
\end{equation}
is a genuine character of $\overline{A(F)}$, and every genuine character
of $\overline{A(F)}$ is of this form for exactly one pair
$(\chi_1,\epsilon)$. Thus for each $\chi_1$ there are exactly two genuine
characters extending $\zeta\chi_1$.
\end{lemma}

\begin{proof}
By Lemma~\ref{lem:Acocycle}(a), a genuine function
$(a(x),\zeta)\mapsto\zeta^{-1}\psi(x)$ is a character if and only if
$\psi(xy)=(x,y)\psi(x)\psi(y)$ for all $x,y$. If $\psi$ is such a
function, put $\chi_1=\psi|_{A_2}$ and
$\epsilon=(\varpi,\varpi)\psi(\varpi)$. For $\nu(x)$ odd write
$x=\varpi x_0$; then
$\psi(x)=(\varpi,x_0)\psi(\varpi)\chi_1(x_0)=\epsilon(x,\varpi)\chi_1(x/\varpi)$,
because $(x,\varpi)=(\varpi,\varpi)(x_0,\varpi)$; and
$\chi_1(\varpi^2)=\psi(\varpi^2)=(\varpi,\varpi)\psi(\varpi)^2
=(\varpi,\varpi)\epsilon^2$, which is (\ref{eq:eps}). Conversely, let
$\psi$ be defined by (\ref{eq:chi-eps}) with $\zeta=1$. If $\nu(x)$,
$\nu(y)$ are both even the relation is clear. If $\nu(x)$ is odd and
$\nu(y)$ even, it reads $(xy,\varpi)=(x,y)(x,\varpi)$, i.e.\
$(y,x\varpi)=1$, which holds because $y$ and $x\varpi$ have even
valuation. If both are odd, it reads
$\chi_1(xy)=(x,y)\epsilon^2(x,\varpi)(y,\varpi)\chi_1(xy/\varpi^2)$; by
(\ref{eq:eps}) the right side is
$(x,y)(x,\varpi)(y,\varpi)(\varpi,\varpi)\chi_1(xy)
=(x\varpi,y\varpi)\chi_1(xy)=\chi_1(xy)$, since $x\varpi$ and $y\varpi$
have even valuation.
\end{proof}

Since $A(F)/A_2(F)$ has order 2, the characters of $A(F)$ are
also in 2-1 correspondence with the characters of $A_2(F)$:
given $\chi_1\in \Pi(A_2(F))$, its two extensions to $A(F)$ are
\begin{equation}\label{eq:chi-theta}
\chi_\vartheta (a(x))=
\left\{
\begin{array}{ll}
\chi_1(x), & \nu(x)\ \text{even},\\
\vartheta\, \chi_1(x/\varpi), & \nu(x)\ \text{odd},
\end{array}
\right.
\qquad \vartheta^2 = \chi_1(\varpi^2).
\end{equation}

Comparing (\ref{eq:eps}) with $\vartheta^2=\chi_1(\varpi^2)$, we see that
$\epsilon=\vartheta$ is impossible when $(\varpi,\varpi)=-1$. Since
$(\varpi,\varpi)=(\varpi,-1)=(-1)^{(q-1)/2}$, this happens exactly when
$q\equiv 3\pmod 4$. Fix once and for all $\eta_\varpi\in{\mathbb C}$ with
\[
\eta_\varpi^2=(\varpi,\varpi)^{-1},
\]
and take $\eta_\varpi=1$ when $(\varpi,\varpi)=1$.

\begin{definition} Let $\chi_1\in \Pi(A_2(F))$ be given.
We say that
$\overline{\chi}_\epsilon\in \Pi(\overline{A(F)})$ {\bf corresponds}
to $\chi_\vartheta \in \Pi(A(F))$ if $\epsilon =\eta_\varpi\vartheta$.
If $\chi \in \Pi(A(F))$ is given we let
$\overline{\chi}\in \Pi(\overline{A(F)})$ denote
its corresponding genuine character.
\label{def:corrchar}
\end{definition}

If $\vartheta^2=\chi_1(\varpi^2)$ then $\epsilon=\eta_\varpi\vartheta$
satisfies (\ref{eq:eps}), so Definition~\ref{def:corrchar} gives a
bijection $\chi\mapsto\overline{\chi}$ between $\Pi(A(F))$ and the genuine
characters of $\overline{A(F)}$. It depends on the choices of $\varpi$
and $\eta_\varpi$. (The 1995 draft required $\epsilon=\vartheta$, which is
inconsistent when $q\equiv3\pmod 4$.)

\begin{example}
Let $F={\mathbb Q}_3$, $\varpi=3$, and $\chi_1=1$. Then
$(3,3)_3=(3,-1)_3=-1$, so $\vartheta=\pm1$ while $\epsilon=\pm i$. With
$\eta_3=i$, the trivial character $\chi_{1}$ of $A({\mathbb Q}_3)$
corresponds to $\overline{\chi}_i$, and the character $\chi_{-1}$
($a(x)\mapsto(-1)^{\nu(x)}$) corresponds to $\overline{\chi}_{-i}$. For
$F={\mathbb Q}_5$, $(5,5)_5=1$, $\eta_5=1$, and the definition reduces to
that of the 1995 draft.
\end{example}

\subsubsection{Corresponding functions on $\overline{A(F)}_{reg}$}

We now determine the conditions on
$\phi \in C_c^\infty(A(F))$,
$\overline{\phi}
\in C_c^\infty(\overline{A(F)})_{gen}$
under which
$\phi^{\wedge}(\chi)=
\overline{\phi}^{\wedge}(\overline{\chi})$ for all $\chi \in \Pi(A(F))$.
Let $\chi=\chi_\vartheta$ extend $\chi_1$. Then
\begin{equation}\label{eq:FT-lin}
\phi^{\wedge}(\chi)=
\int_{A(F)}\phi(a)\chi(a)\, da
=\int_{A_2(F)}\phi(a)\chi_1(x_a)\, da +
\vartheta \int_{A(F)-A_2(F)}\phi(a)\chi_1(x_a/\varpi)\, da .
\end{equation}
Since $\overline{\phi}\,\overline{\chi}$ is invariant under
$\mu_2$ and $d\varsigma$ gives $\mu_2$ volume $1$, and since
$\overline{\chi}=\overline{\chi}_{\epsilon}$ with
$\epsilon=\eta_\varpi\vartheta$,
\begin{equation}\label{eq:FT-gen}
\begin{aligned}
\overline{\phi}^{\wedge}(\overline{\chi})&=
\int_{A(F)}\overline{\phi}(a,1)\overline{\chi}(a,1)\, da\\
&=\int_{A_2(F)}\overline{\phi}(a,1)\chi_1(x_a)\, da +
\eta_\varpi\vartheta \int_{A(F)-A_2(F)}\overline{\phi}(a,1)\,(x_a,\varpi)^{-1}
\chi_1(x_a/\varpi)\, da.
\end{aligned}
\end{equation}
(Draft~1 omitted the factor $(x_a,\varpi)^{-1}$ in (\ref{eq:FT-gen}) and
wrote $\epsilon$ in place of $\vartheta$ in (\ref{eq:FT-lin}).)

\begin{definition}\label{def:corrA} We say that
$\phi \in C_c^\infty(A(F))$ {\bf corresponds} to
$\overline{\phi}
\in C_c^\infty(\overline{A(F)})_{gen}$
if
\[
\overline{\phi}(a,1)=
\left\{
\begin{array}{ll}
\phi(a), & a\in A_2(F),\\
\eta_\varpi^{-1}(x_a,\varpi)\,\phi(a), & a\in A(F)-A_2(F).
\end{array}
\right.
\]
\end{definition}

\begin{lemma}\label{lem:FT-match}
(a) If $\phi \in C_c^\infty(A(F))$ corresponds to
$\overline{\phi}\in C_c^\infty(\overline{A(F)})_{gen}$
then
$\phi^{\wedge}(\chi)=
\overline{\phi}^{\wedge}(\overline{\chi})$
for all $\chi \in \Pi(A(F))$. Conversely, if this holds for all $\chi$,
then $\phi$ corresponds to $\overline{\phi}$.

(b) If $\phi$ corresponds to $\overline{\phi}$ and $\phi(a^{-1})=\phi(a)$
for all $a$, then $\overline{\phi}((a,1)^w)=\overline{\phi}(a,1)$ for all
$a$.
\end{lemma}

\begin{proof}
(a) The first terms of (\ref{eq:FT-lin}) and (\ref{eq:FT-gen}) agree
when $\overline{\phi}(a,1)=\phi(a)$ on $A_2(F)$. In the second term of
(\ref{eq:FT-gen}), substituting
$\overline{\phi}(a,1)=\eta_\varpi^{-1}(x_a,\varpi)\phi(a)$ gives
$\vartheta\int\phi(a)(x_a,\varpi)^2\chi_1(x_a/\varpi)\,da$, which is the
second term of (\ref{eq:FT-lin}). Conversely, as $\chi_1$ ranges over
$\Pi(A_2(F))$ and $\vartheta$ over both square roots of
$\chi_1(\varpi^2)$, Fourier inversion on $A_2(F)$ and on its coset
$A(F)-A_2(F)$ forces both pointwise identities.

(b) By Lemma~\ref{lem:Acocycle}(b), $(a,1)^w=(a^{-1},1)$, and
$(x_a^{-1},\varpi)=(x_a,\varpi)$.
\end{proof}

\begin{definition} We say that
$f \in C_c^\infty(G(F))$ {\bf corresponds oddly
on the hyperbolic set} to
$\overline{f}\in C_c^\infty(\overline{G(F)})_{gen}$
if the normalized orbital integral $F_f$ on $A(F)_{reg}$ corresponds to
the normalized orbital integral $F_{\overline{f}}$ on
$\overline{A(F)}_{reg}$ in the sense of Definition~\ref{def:corrA}.
\end{definition}

By Lemma~\ref{lem:FT-match}(b), this definition is compatible with the
invariance of $F_f$ and $F_{\overline f}$ under the Weyl group.

\begin{definition}
We say $f,\overline{f}$ {\bf correspond evenly on the hyperbolic set}
if $\Delta(g)O_{\overline{f}}(g)=\Delta(h)O_f(h)$, when
$g=t(h)$ is regular, and $\Delta(g)O_{\overline{f}}(g,1)=0$,
if $g\notin {\rm image}(t)$. Here $O_f$ is the orbital integral
\[
O_f(x)=\int_{G_x(F)\backslash G(F)} f(g^{-1}xg)dg,
\]
and $\Delta$, $t$ are as in \cite{Fl}. (This is the definition of the
1995 draft; the corresponding matching of representations, as in
\cite{Fl}, is not verified here.)
\end{definition}

\begin{claim}[1995 draft; false as stated]\label{lem:weighted-small}
For $f$ of sufficiently small support, if
$f \in C_c^\infty(G(F))$ corresponds oddly on $A(F)_{reg}$ to
$\overline{f}\in C_c^\infty(\overline{G(F)})_{gen}$
then the weighted orbital integral of
$f$ corresponds to the weighted orbital integral of
$\overline{f}$ in the sense above (as functions on $A(F)$).
\end{claim}

\begin{remark}\label{rem:weighted-noninv}
(Added in Draft~2.) The 1995 proof was ``a direct calculation''. The
claim cannot hold for \emph{all} corresponding pairs. Indeed, if
$(f,\overline{f})$ correspond, so does $(f,\overline{f}+\overline{h})$
for any genuine $\overline{h}$ whose orbital integrals vanish, for
example $\overline{h}=\overline{g}-\overline{g}^{\,y}$ with
$\overline{g}^{\,y}(x)=\overline{g}(yxy^{-1})$ and $y\in\overline{G(F)}$.
The claim would then force $J_A(a,\overline{h})=0$, i.e.,
\[
|D(a)|^{1/2}\int_{\overline{A(F)}\backslash\overline{G(F)}}
\overline{g}(x^{-1}ax)\,\bigl(v_A(x)-v_A(xy)\bigr)\,dx=0
\]
for all small-support $\overline{g}$ and all $y$, which fails because
weighted orbital integrals are not invariant distributions
\cite{A7}. A correct statement must either restrict to pairs constructed
in a specific way from one another, or replace the weighted orbital
integrals $J_A$ by the invariant distributions $I_A$ of the invariant
trace formula \cite{A6}. The same objection applies to the remark on
real weighted orbital integrals in \S\ref{sec:realA}. Neither statement
is used in Chapter~\ref{ch:global}.
\end{remark}

\begin{proposition}\label{prop:hyp-ps}
Suppose that $\overline{f}\in C^\infty_c(\overline{G(F)})_{gen}$
corresponds oddly on $A(F)_{reg}$ to
$f\in C_c^\infty(G(F))$. Let $\chi\in\Pi(A(F))$ and let $\overline{\chi}$
be the corresponding genuine character
(Definition~\ref{def:corrchar}). Then
\[
\tr I_A(\overline{\chi})(\overline{f})=\tr I_A(\chi)(f).
\]
\end{proposition}

\begin{proof}
By \S\ref{sec:ps-char},
\[
\tr I_A(\chi)(f)=(F_f)^\wedge(\chi),\qquad
\tr I_A(\overline{\chi})(\overline{f})=(F_{\overline{f}})^\wedge(\overline{\chi}).
\]
These agree by Lemma~\ref{lem:FT-match}(a). (The metaplectic character
formula of \S\ref{sec:ps-char} is quoted there with its proof omitted,
so the proposition depends on it.)
\end{proof}

\begin{example}
Let $F={\mathbb Q}_3$, $\varpi=3$, $\eta_3=i$, and let $\phi$ be the
characteristic function of $\{a(x)\ |\ x\in 3{\mathbb Z}_3^\times\}$.
Then the corresponding $\overline\phi$ has
$\overline{\phi}(a(3u),1)=-i\,(3u,3)_3=-i\,(3,3)_3(u,3)_3=i\left(\frac{u}{3}\right)$
for $u\in{\mathbb Z}_3^\times$, where $(\frac{\cdot}{3})$ is the Legendre
symbol, and $\overline\phi(a(x),1)=0$ otherwise. For the trivial character
$\chi=\chi_1$ (with $\chi_1=1$, $\vartheta=1$), $\phi^\wedge(\chi)={\rm vol}(3{\mathbb Z}_3^\times)$,
and $\overline{\phi}^\wedge(\overline\chi)
=i\int i\,(\tfrac{u}{3})(3u,3)_3\,da=-\int(\tfrac{u}{3})\cdot(-(\tfrac{u}{3}))\,da
={\rm vol}(3{\mathbb Z}_3^\times)$, as it should be.
\end{example}

A similar result holds in the real case, which is discussed in the
next subsection.

\subsection{Corresponding functions on
$\overline{A({\mathbb{R}})}$}\label{sec:realA}

Let $F={\mathbb{R}}$ in this subsection, and let $a(x)$ denote the
diagonal matrix with entries $x,x^{-1}$. The real Hilbert symbol is
\[
(x,y)_2=
\left\{
\begin{array}{ll}
-1, & x<0\ \text{and}\ y<0,\\
1, & \text{otherwise}.
\end{array}
\right.
\]
(Draft~1 gave a formula depending only on the sign of $x$.) By
Lemma~\ref{lem:Acocycle}(a), which holds for every local field,
$(a(x),\zeta)(a(y),\zeta')=(a(xy),(x,y)_2\zeta\zeta')$.

Let $\chi_1$ be a character of ${\mathbb{R}}_+^\times$. Define
$\chi$ on $A({\mathbb{R}})$ by
\[
\chi(a(x))=\chi_\epsilon(a(x))
=
\left\{
\begin{array}{ll}
\chi_1(x), & x>0,\\
\epsilon \chi_1(-x), & x<0,
\end{array}
\right.
\]
where $\epsilon =\pm 1$. Define
$\overline{\chi}$ on $\overline{A({\mathbb{R}})}$ by
\[
\overline{\chi}(a(x),\zeta)=\overline{\chi}_\eta(a(x),\zeta)
=
\left\{
\begin{array}{ll}
\zeta^{-1}\chi_1(x), & x>0,\\
\zeta^{-1}\eta \chi_1(-x), & x<0,
\end{array}
\right.
\]
where $\eta=\pm i$. This is a genuine character: the only nontrivial
condition, at $x=y=-1$, reads $1=(-1,-1)_2\,\eta^2=-\eta^2$. Every
genuine character of $\overline{A({\mathbb R})}$ is of this form (compare
\cite{GS}). We say that
$\chi_\epsilon$ {\bf corresponds} to $\overline{\chi}_\eta$ if
$\eta =\epsilon i$; this is a bijection
$\chi\mapsto\overline{\chi}$ from $\Pi(A({\mathbb R}))$ onto the genuine
characters of $\overline{A({\mathbb R})}$.

\begin{definition}\label{def:corrA-real}
We say that $\phi \in C_c^\infty(A({\mathbb R}))$ {\bf corresponds} to
$\overline{\phi}\in C_c^\infty(\overline{A({\mathbb R})})_{gen}$
if
$\phi^{\wedge}(\chi)=
\overline{\phi}^{\wedge}(\overline{\chi})$
for all $\chi \in \Pi(A({\mathbb R}))$.
\end{definition}

\begin{lemma}\label{lem:realA-match}
$\phi$ corresponds to $\overline{\phi}$ if and only if
\begin{equation}\label{eq:realA-match}
\overline{\phi}(a(x),1)=
\left\{
\begin{array}{ll}
\phi(a(x)), & x>0,\\
-i\,\phi(a(x)), & x<0.
\end{array}
\right.
\end{equation}
\end{lemma}

\begin{proof}
As in the p-adic case,
$\phi^\wedge(\chi_\epsilon)=\int_{x>0}\phi(a(x))\chi_1(x)\,d^\times x
+\epsilon\int_{x<0}\phi(a(x))\chi_1(-x)\,d^\times x$ and
$\overline{\phi}^\wedge(\overline{\chi}_{\epsilon i})
=\int_{x>0}\overline{\phi}(a(x),1)\chi_1(x)\,d^\times x
+\epsilon i\int_{x<0}\overline{\phi}(a(x),1)\chi_1(-x)\,d^\times x$.
These agree for all $\chi_1$ and both values of $\epsilon$ if and only if
$\overline{\phi}=\phi$ on $x>0$ and $i\,\overline{\phi}(a(x),1)=\phi(a(x))$
on $x<0$, i.e.\ $\overline{\phi}(a(x),1)=-i\,\phi(a(x))$.
\end{proof}

(Draft~1 stated (\ref{eq:realA-match}) with the factor $\epsilon$ in
place of $-i$; that is not compatible with $\eta=\epsilon i$, and since
$\epsilon$ depends on the character it cannot appear in a condition on
functions.) By Lemma~\ref{lem:Acocycle}(b), $(a,1)^w=(a^{-1},1)$, so the
condition (\ref{eq:realA-match}) is compatible with Weyl group
invariance.

\begin{definition} We say that
$f \in C_c^\infty(G({\mathbb R}))$ {\bf corresponds oddly on the
hyperbolic set} to
$\overline{f}\in C_c^\infty(\overline{G({\mathbb R})})_{gen}$
if $F_f$ corresponds to
$F_{\overline{f}}$
in the sense of Definition~\ref{def:corrA-real}.
\end{definition}

With this definition, the proof of Proposition~\ref{prop:hyp-ps} gives
$\tr I_A(\overline{\chi})(\overline{f})=\tr I_A(\chi)(f)$ for
$f$, $\overline{f}$ corresponding oddly on the hyperbolic set, given the
character formula of \S\ref{sec:ps-char} for ${\mathbb R}$.

\begin{remark}
The 1995 draft added here that, by a direct calculation, weighted
orbital integrals of functions corresponding on
$A({\mathbb{R}})_{reg}$ also correspond. For arbitrary corresponding
pairs this cannot hold, for the reason given in
Remark~\ref{rem:weighted-noninv}.
\end{remark}

\subsection{Corresponding functions on the elliptic set}\label{sec:elliptic}

Finally, we concentrate on the elliptic set. Let ${\cal T}_{ell}$ be a
set of representatives of the $G(F)$-conjugacy classes of anisotropic
(elliptic) Cartan subgroups of $G$ over $F$. If $F={\mathbb R}$, there is
one class, represented by $SO(2)$. If $F$ is p-adic, there are several:
each anisotropic torus is isomorphic to the norm-one torus of a
quadratic extension $E/F$, and for $p$ odd there are three such $E$ and
two $G(F)$-conjugacy classes of tori for each $E$. (Draft~1 used only one
elliptic Cartan in the p-adic case, which is not enough.)

For $T$ a Cartan subgroup, $t\in T(F)$ regular, and
$f\in C_c^\infty(G(F))$, let $\Phi_f(t)=|D(t)|\,O_f(t)$, where
$O_f(t)=\int_{T(F)\backslash G(F)}f(x^{-1}tx)\,dx/dt$, and define
$\Phi_{\overline{f}}$ on $\overline{T(F)}_{reg}$ in the same way for
$\overline{f}\in C_c^\infty(\overline{G(F)})_{gen}$. With this
normalization (and the measures of \S\ref{sec:local}), the Weyl
integration formula reads
\[
\Theta _{\pi}(f)=\sum_{T\in{\cal T}}\frac 1{|W_T|}\int_{T(F)}\Theta _{\pi}(t)
\Phi_f(t)\,dt,\;\;\;f\in C_c^\infty (G(F)),
\]
and
\[
\Theta _{\overline{\pi } }(\overline{f})
=\sum_{T\in{\cal T}}\frac 1{|W_T|}\int_{\overline{T(F)}}
\Theta _{\overline{\pi } }(t)
\Phi_{\overline{f}}(t)\,dt\,d\varsigma,
\;\;\;\overline{f}\in C_c^\infty (\overline{G(F)})_{gen},
\]
where ${\cal T}$ consists of $A$ and the elements of ${\cal T}_{ell}$,
and $W_T=N_{G(F)}(T)/T(F)$. Since $\Theta_{\overline{\pi}}$ and
$\Phi_{\overline{f}}$ are both genuine, the integrand on
$\overline{T(F)}$ is invariant under $\mu_2$, and the integral over
$\overline{T(F)}$ equals $\int_{T(F)}\Theta_{\overline\pi}(t,1)\Phi_{\overline f}(t,1)\,dt$.

\begin{definition}\label{def:ell-corr}
We say that $f$ and $\overline{f}$ {\bf have vanishing hyperbolic
orbital integrals} if $\Phi_f=0$ on $A(F)_{reg}$ and
$\Phi_{\overline f}=0$ on $\overline{A(F)}_{reg}$.

(a) If $f,\overline{f}$ have vanishing hyperbolic orbital integrals and
satisfy
\[
\Phi_f(t)=\Phi_{\overline{f}}(t,1),\ \ \ \ t\in T(F)_{reg},\ \text{for every } T\in{\cal T}_{ell},
\]
then we say that they \textbf{correspond oddly on the elliptic set}.

(b) If $f,\overline{f}$ have vanishing hyperbolic orbital integrals and
satisfy, for every $T\in{\cal T}_{ell}$,
\[
\Phi_{\overline{f}}(g)
=
\left\{
\begin{array}{ll}
\Phi_f(h), & g=\tau(h),\ h\in T(F)_{reg},\\
0, & g\notin \tau(T(F)),
\end{array}
\right.
\]
then we say that they \textbf{correspond evenly on the elliptic set}.
Here $\tau:h\longmapsto (h,1)^2$ is the map denoted
$h\longmapsto \tilde{h}$ in \cite{Fl}.
\end{definition}

For instance, if $f$ and $\overline{f}|_{G(F)}$ are supported in the
(open) set of regular elliptic elements, they have vanishing hyperbolic
orbital integrals.

\begin{lemma}\label{lem:ell-odd}
Let $f,\overline{f}$ correspond oddly on the elliptic set. Let $\pi$ and
$\overline{\pi}$ be an admissible representation of $G(F)$ and a genuine
admissible representation of $\overline{G(F)}$ whose characters satisfy
\begin{equation}\label{eq:ell-char-rel}
\Theta_{\overline{\pi}}(t,\zeta)=\zeta\,\Theta_\pi(t),\qquad
t\in T(F)_{reg},\ \zeta\in\mu_2,\ \text{for every }T\in{\cal T}_{ell}.
\end{equation}
Then $\Theta _{\pi }(f)=\Theta _{\overline{\pi }}(\overline{f})$.
\end{lemma}

\begin{proof}
In both Weyl integration formulas the term for $T=A$ vanishes, because
the hyperbolic orbital integrals vanish. For $T\in{\cal T}_{ell}$, the
term for the cover equals
$|W_T|^{-1}\int_{T(F)}\Theta_{\overline\pi}(t,1)\Phi_{\overline f}(t,1)\,dt
=|W_T|^{-1}\int_{T(F)}\Theta_{\pi}(t)\Phi_{f}(t)\,dt$ by
(\ref{eq:ell-char-rel}) and Definition~\ref{def:ell-corr}(a), which is
the term for $G$. Summing over ${\cal T}_{ell}$ gives the result.
\end{proof}

The character relation (\ref{eq:ell-char-rel}) for supercuspidal
representations is Claim~\ref{prop:164}, which needs correction
(Remark~\ref{rem:sc-caveat}); Lemma~\ref{lem:ell-odd} is therefore
stated with (\ref{eq:ell-char-rel}) as a hypothesis.

\subsection{Corresponding spherical functions, $p> 2$}\label{sec:spherical}

We sketch the Satake correspondence for $\overline{G(F)}$
following \cite{J}, \S 1.3 and using its notation.

Recall that the genuine characters of
$\overline{A(F)}$ are in 1-1 correspondence with
the characters of $A(F)$. Moreover, from results recalled in
Chapter~\ref{ch:reps} we deduce that all
the unramified genuine representations of
$\overline{G(F)}$ belong to the unramified principal series.
Replacing $GL(r)$ by
$SL(2)$ in \cite{J}, \S 1.3, we first observe that the
set $\Pi_{unr.p.s.}(\overline{G(F)})_{gen}$ is
parametrized by the $W$-orbit of unramified
genuine characters of $\overline{A(F)}$. Therefore, we have the following

\begin{lemma} The set $\Pi_{unr.p.s.}(\overline{G(F)})_{gen}$ is
in natural 1-1 correspondence with
$\Pi_{unr.p.s.}(G(F))$.
\end{lemma}

Following now \cite{J}, \S 1.4, we obtain a transform
\begin{equation}
\begin{array}{c}
{\cal{H}}(\overline{G(F)},K)_{gen}\rightarrow
{\cal{A}}_{cl.fcn}(PGL(2,{\mathbb{C}}))\\
f\longmapsto S_0 f(\chi)=\tr I_A(\chi)(f).
\end{array}
\end{equation}

This is essentially the inverse Fourier transform of the Satake
transform, which we define next. The Satake transform of a
genuine spherical function is defined as in the non-metaplectic Case:
\begin{equation}
Sf(a)=\delta(a)^{1/2}\int_{N(F)}f(an)dn.
\end{equation}
The relationship between this expression and an orbital integral
is well-known - see for example \cite{C}.

This defines an algebra isomorphism
\begin{equation}
\begin{array}{c}
{\cal{H}}(\overline{G(F)},K)_{gen}\rightarrow
{\cal{H}}(\overline{A(F)},K\cap A(F))^W_{gen}\\
f\longmapsto Sf(a).
\end{array}
\end{equation}

By Lemma~\ref{lem:FT-match}, the correspondence $\phi\mapsto\overline\phi$
of Definition~\ref{def:corrA} preserves Fourier transforms, hence
convolution; it maps $W$-invariant functions to $W$-invariant functions;
and, since the splitting $\kappa$ is $a(u)\mapsto(a(u),1)$ on $K\cap A(F)$
(by (\ref{eq:sv}), as $c=0$), it maps functions invariant under
$K\cap A(F)$ to genuine functions invariant under $\kappa(K\cap A(F))$
(the factor $(x_a,\varpi)$ transforms correctly under $a\mapsto a\,a(u)$
by Lemma~\ref{lem:Acocycle}(a)). It therefore gives an algebra isomorphism
${\cal{H}}(A(F),K\cap A(F))^W\cong
{\cal{H}}(\overline{A(F)},\kappa(K\cap A(F)))^W_{gen}$, and so
\[
{\cal{H}}(G(F),K)\cong{\cal{H}}(\overline{G(F)},K)_{gen}.
\]

\begin{definition}\label{def:satake-corr}
We say that $f\in{\cal H}(G(F),K)$ and
$\overline f\in{\cal H}(\overline{G(F)},K)_{gen}$ {\bf correspond under
the Satake isomorphism} if $Sf$ corresponds to $S\overline{f}$ in the
sense of Definition~\ref{def:corrA}.
\end{definition}

If $f$ and $\overline f$ correspond under the Satake isomorphism and $\chi$ is
an unramified character of $A(F)$, then
$\tr I_A(\overline\chi)(\overline f)=(S\overline f)^\wedge(\overline\chi)
=(Sf)^\wedge(\chi)=\tr I_A(\chi)(f)$, by Lemma~\ref{lem:FT-match}(a).
This is the \emph{spectral} matching supplied by the Satake isomorphism.

\begin{remark}\label{rem:satake-oi}
The 1995 draft concluded: ``This implies the matching of orbital
integrals for spherical functions.'' That does not follow from the
existence of an algebra isomorphism alone. What does follow is the
following. For $f\in{\cal H}(G(F),K)$ and regular $a\in A(F)$, the
normalized hyperbolic orbital integral $F_f(a)$ equals the Satake
transform $Sf(a)$ \cite{C}; granted the analogous identity on the cover
(the same computation, using the splitting $\kappa$ and
Lemma~\ref{lem:facts}(h); the cocycle bookkeeping has not been written out
here), functions corresponding under the Satake isomorphism correspond
oddly on the hyperbolic set. The \emph{elliptic} orbital integrals of
corresponding spherical functions are not determined by their Satake
transforms in any formal way, and their matching (a fundamental lemma
for this comparison) is not established here. It is one of the
ingredients assumed in Hypothesis~\ref{hyp:GH}.
\end{remark}

\begin{example}
For the unit elements $f=1_K$ and $\overline f=e_K$ of the two Hecke
algebras, $Sf=1_{K\cap A(F)}$ and $S\overline f$ is supported on
$\overline{A_2(F)}$, where it equals $(a,\zeta)\mapsto\zeta\,1_{K\cap A(F)}(a)$;
these correspond (the factor $\eta_\varpi^{-1}(x_a,\varpi)$ only enters off
$A_2(F)$). Hence $\tr I_A(\overline\chi)(e_K)=\tr I_A(\chi)(1_K)=1$ for
unramified $\chi$.
\end{example}

\begin{remark}\label{rem:weighted-units}
On the matching of
weighted orbital integrals of spherical functions:
every global smooth genuine function
$\overline{f}=\otimes_v\overline{f}_v$ is a restricted tensor product
of local smooth functions, almost all of which are units of the
local Hecke algebra of genuine spherical functions on
$\overline{G(k_v)}$. In order to match weighted orbital integrals
of such global functions, one must match them locally.
The 1995 draft states that a direct calculation of the weighted orbital
integrals shows that the weighted orbital integral of the unit element
$\overline{f}$ of the genuine spherical Hecke algebra on $\overline{A(F)}$
matches that of the unit element $f$ of the spherical Hecke algebra of
$G(F)$ on $A(F)$. That calculation is not reproduced here. (This is a
statement about a specific pair, so the objection of
Remark~\ref{rem:weighted-noninv} does not apply to it.)
\end{remark}

\begin{lemma}
Finally, if $f,\overline{f}$ are the respective unit elements
of the respective Hecke algebras, $p\not= 2$, then
$\Theta_\pi(f)=\Theta_{\overline{\pi}}(\overline{f})$,
where
$\pi \leftrightarrow \overline{\pi}$ are corresponding
unramified principal series as above.
\end{lemma}

\section{The dyadic places}\label{sec:dyadic}

In this section $F$ is a dyadic field, i.e., $p=2$, and we write
$G=SL(2,F)$ and $\overline{G}=\overline{SL(2,F)}$.

In the 1995 draft, this section attempted to construct, following
Howe and Moy \cite[ch.~3]{HM}, an injection of a metaplectic generalized
Hecke algebra $\overline{\cal H}_k$ (supported on a ``maximal split
subset'' $\overline{G_s}$ of $\overline{G}$) into the generalized Hecke
algebra ${\cal H}(G//B_k)$ of a deep filtration subgroup $B_k$ of the
Iwahori subgroup, and to deduce from it a matching of orbital integrals
and of characters at dyadic places. That argument is withdrawn in this
version, for the following reasons.

\begin{enumerate}
\item The presentation of ${\cal H}(G//B_k)$ included the braid relation
$f_{s_1}*f_{s_2}*f_{s_1}=f_{s_2}*f_{s_1}*f_{s_2}$, where
$s_1=\left(\begin{smallmatrix}0&-1\\1&0\end{smallmatrix}\right)$ and
$s_2=\left(\begin{smallmatrix}0&-\varpi^{-1}\\\varpi&0\end{smallmatrix}\right)$.
Since $s_1s_2={\rm diag}(-\varpi,-\varpi^{-1})$ has infinite order, the affine
Weyl group of $SL(2)$ is infinite dihedral and
$s_1s_2s_1\neq s_2s_1s_2$, so this relation does not hold. (It is the
analog of a relation valid for $GL(n)$, $n\geq 3$.)

\item The set
$X=\{g\in G\ |\ c\in F^{\times 2},\ d=0\ {\rm or}\ d\in F^{\times 2}\}$
was claimed to lie in a maximal subset $G_s$ of $G$ which is closed under
inversion and on which the cocycle $\beta$ is trivial. However, the
lower left entry of $g^{-1}$ is $-c$, so $X$ is not closed under inversion
when $-1\notin F^{\times 2}$ (as for $F={\mathbb Q}_2$); and the proof
computed $\beta(g_1,g_2^{-1})$ rather than $\beta(g_1,g_2)$.

\item The map
$N\left(\begin{smallmatrix}a&b\\c&d\end{smallmatrix}\right)=
\left(\begin{smallmatrix}a\epsilon_d^{-1}&b\epsilon_c^{-1}\\
c\epsilon_c&d\epsilon_d\end{smallmatrix}\right)$,
where $\epsilon_x$ denotes a fixed representative in ${\cal O}$ of the
square class of $x$, is not defined when $c=0$ or $d=0$, does not
map the Iwahori subgroup into itself when $c$ has odd valuation (since
then $b\epsilon_c^{-1}\notin{\cal O}$ in general), and is not
multiplicative.

\item The restrictions of genuine functions to $\overline{G_s}$ do not
form a convolution algebra, since $G_s$ is not a subgroup; the
computation of $\overline{f}_{g_1}*\overline{f}_{g_2}$ dropped cocycle
factors at points $x$ of $B_kg_2B_k$ outside $G_s$; and the
definition of the functions $\overline{f}_g$ ignored the splitting of
$\overline{B_k}$, which need not be $b\mapsto (b,1)$.

\item An injection of algebras does not by itself induce a map on
irreducible modules, so the claimed surjection
${\cal H}(G//B_k)^\wedge\rightarrow{\cal H}(\overline{G}//B_k)^\wedge$
with matching traces does not follow.

\item The orbital integral identity at a dyadic place was derived from
a global argument which assumed that the spectral sides of the two trace
formulas agree. That agreement is what Chapter~\ref{ch:global} needs to
establish, so the argument is circular.
\end{enumerate}

The generalized Hecke algebra material is revised separately in the
author's paper \cite{J3}.

In this version no local matching of functions or of characters is
established at the dyadic places. Instead, the dyadic places are always
put in the exceptional set $S$ of Chapter~\ref{ch:global}, and what is
needed there is assumed in Hypotheses~\ref{hyp:GH} and~\ref{hyp:LH}.

\begin{remark}
(1) The construction of \S\ref{sec:charsA} does not extend to $p=2$
without change: there $\overline{A_2(F)}$ need not be the direct product
$A_2(F)\times\mu_2$. For example, if $F={\mathbb Q}_2$, then
$(-I,1)^2=(I,(-1,-1)_2)=(I,-1)$, so the lifts of $-I\in A_2(F)$ have
order $4$.

(2) At Iwahori level, there are Hecke algebra correspondences at
$p=2$ for the metaplectic group (Takeda and Wood \cite{TW}, extending
work of Gan and Savin for $p\neq 2$), but they relate genuine
representations of $\overline{SL(2,F)}$ to representations of
$SO(3)\cong PGL(2)$ and its inner form, not of $SL(2)$. Whether such
results can be used to verify Hypothesis~\ref{hyp:LH} at dyadic places
is not addressed here.
\end{remark}

%% (aabook2b.tex withdrawn in Draft 1; see revision_log.md)



%% ===== begin aabook2c.tex =====
\section{On the correspondence of limits of d.s. for
$\overline{SL(2,\mathbb{R})}$}\label{sec:realds}

Let $T$ denote a Cartan of $G=SL(2)$.
%$F$ be a field and
%$\overline{F}$ a separable algebraic closure. By the density
%theorem (see for example, [Hum], \S 22.2), the conjugate $T^G$ contains
%a $\overline{F}$-Zariski-open dense subset of $G$.
%The $F$-rational
%points $T^G(F)$ remain $F$-Zariski open in $G(F)$ (but not necessarily dense,
%see for example, Theorem 1 of ch II, \S 4, in [Mum]). If $F$ is a
%complete local field then the points $T^G(F)$ are open (in the induced
%$F$-manifold topology) in $G(F)$. However, in general,
%$T(F)^{G(F)}\not= T^G(F)$.
If $F={\mathbb{R}}$ and $T$ is an
anisotropic Cartan of $G$ defined over $F$ then recall $T(F)^{G(F)}$ is the
elliptic set of $G$.

We shall describe 3 different correspondences
$c_i:\overline{\pi}\longmapsto \pi$, $i=1,2,3$, between representations of
$\overline{SL(2,\mathbb{R})}$ and those of $SL(2,\mathbb{R})$. They are
consistent with a ``matching'' $c_i^*:\overline{f}\longmapsto f$,
$i=1,2,3$, of functions whose hyperbolic orbital integrals vanish, in the
sense that the characters match:
$\Theta_{\overline{\pi}}(\overline{f})=\Theta_{\pi}(f)$, where
$\pi=c_i(\overline{\pi})$, $f=c_i^*(\overline{f})$, $i=1,2,3$
(Lemma~\ref{lem:realds}).
Furthermore, if $\pi$ is any element of the discrete series or limit of
discrete series of $G(F)$ then $\pi$ occurs in the image of at least one
of the $c_i$.

\subsection{Discrete series for $\overline{SL(2,\mathbb{R})}$}

Let
\begin{align*}
\overline{B} &=\overline{M}\cdot A(\mathbb{R})^o\cdot N(\mathbb{R}), \\
A(\mathbb{R})^o &=\{a(x)=\left(
\begin{array}{cc}
x & 0 \\
0 & x^{-1}
\end{array}
\right) \;|\;x>0\}, \\
N(\mathbb{R}) &=\{n=\left(
\begin{array}{cc}
1 & x \\
0 & 1
\end{array}
\right) \;|\;x\in \mathbb{R}\}, \\
\overline{M} &=<\gamma >,\;\gamma =(-\left(
\begin{array}{cc}
1 & 0 \\
0 & 1
\end{array}
\right) ,1),
\end{align*}
and let
\[
\overline{r}(\theta )=\left(
\begin{array}{cc}
\cos (\theta ) & -\sin (\theta ) \\
\sin (\theta ) & \cos (\theta )
\end{array}
\right) ,\;\;0\leq \theta \leq 4\pi ,
\]
which we use to denote the element of $\overline{K}$ lying over the
displayed rotation, chosen so that $\theta\mapsto\overline r(\theta)$ is a
continuous homomorphism ${\mathbb R}\rightarrow\overline K$ with
$\overline r(0)$ the identity. Then $\overline r(2\pi)=(I,-1)$, since
$\overline K$ is connected (Lemma~\ref{lem:facts}(d)).
Let $\overline{\chi }=\overline{\chi }_{\lambda ,\tau }$ denote the
character of $\overline{B}$ given by
\[
\overline{\chi }(\gamma ^ja(x)n)=e^{ij\tau \pi }x^\lambda ,\;\;\tau \in
\{0,1,1/2,-1/2\},\ \ \ \ \ \ x>0.
\]
This character is genuine if and only if $\tau \in \{1/2,-1/2\}$. Every
genuine principal series is given by
\[
\rho (\lambda ,\tau )=Ind(\overline{\chi },\overline{B},\overline{SL(2,\mathbb{R%
})}).
\]

Let $\Pi (\overline{SL(2,\mathbb{R})})_{ds}$ denote the set of equivalence
classes of genuine unitary discrete series representations of $\overline{%
SL(2,\mathbb{R})}$.

\subsection{Gelbart's parameterization \cite[p.~75]{Ge2}}

Let $k>1$ be an odd integer and let $\tau \in \{1/2,-1/2\}$. Let
\[
\phi _m(na(x)\overline{r}(\theta ))=x^\lambda e^{im\theta }=x^\lambda
e^{i(2s+\tau )\theta },\;\;\;\lambda \in \mathbb{C},\;\;m=2s+\tau ,\;\;s\in
\mathbb{Z}.
\]

These form a collection of functions which can be used to describe a basis
for the discrete series. The discrete series are tabulated as follows:
\[
\begin{array}{cccc}
\mathrm{d.s.\;repn} & k & \tau & \lambda \\
B_{-}^{1/2}(k/2) & k\equiv 1(4) & 1/2 & \frac k2-1 \\
B_{+}^{-1/2}(-k/2) & k\equiv 1(4) & -1/2 & \frac k2-1 \\
B_{+}^{1/2}(-k/2) & k\equiv 3(4) & 1/2 & \frac k2-1 \\
B_{-}^{-1/2}(k/2) & k\equiv 3(4) & -1/2 & \frac k2-1
\end{array}
\]
These representations exhaust $\Pi (\overline{SL(2,\mathbb{R})})_{ds}$. Another
parameterization is as follows:
\[
\begin{array}{cccc}
\mathrm{d.s.\;repn} & k & \tau & \lambda \\
B_{-}^{1/2}(-k/2+2) & k\equiv 3(4) & 1/2 & -\frac k2+1 \\
B_{+}^{-1/2}(k/2-2) & k\equiv 3(4) & -1/2 & -\frac k2+1 \\
B_{+}^{1/2}(k/2-2) & k\equiv 1(4) & 1/2 & -\frac k2+1 \\
B_{-}^{-1/2}(-k/2+2) & k\equiv 1(4) & -1/2 & -\frac k2+1
\end{array}
\]

\subsection{Discrete series for $SL(2,\mathbb{R})$}

Let
\begin{align*}
B &=M\cdot A(\mathbb{R})^o\cdot N(\mathbb{R}), \\
M &=\{\pm I_2\},\;I_2=\left(
\begin{array}{cc}
1 & 0 \\
0 & 1
\end{array}
\right) ,
\end{align*}
and let
\[
r(\theta )=\left(
\begin{array}{cc}
\cos (\theta ) & -\sin (\theta ) \\
\sin (\theta ) & \cos (\theta )
\end{array}
\right) ,\;\;0\leq \theta \leq 2\pi .
\]

For each non-zero $m\in \mathbb{Z}$ there is a $\pi _m\in \Pi (SL(2,\mathbb{R}%
))_{ds}$ whose character $\Theta _{\pi _m}$ satisfies
\[
\Theta _{\pi _m}(r(\theta ))=-sgn(m)
\frac{e^{im\theta }}{e^{i\theta }-e^{-i\theta }},
\;\;\;0<\theta <2\pi ,\;\theta \neq \pi ,
\]
and
\[
\Theta _{\pi _m}((-I_2)^j a(e^t ))=(-1)^{(m-1)j}
\frac{e^{-|mt| }}{|e^t-e^{-t}|}.
\]

For $m$ odd, $\pi _m$ has trivial central character. For $m$ even, $m\neq 0$
($m=0$ corresponds to the two irreducible components of the unitary
reducible principal series), $\pi _m$ has non-trivial central character.
This may be found in \cite{HC2}, page 660.

We denote the two irreducible components of the unitary reducible principal
series by $\pi _0^{+}$ and $\pi _0^{-}$, labeled in such a way that their
character $\Theta _{\pi _0^{\pm }}$ satisfies
\begin{align*}
\Theta _{\pi _0^\varepsilon }(r(\theta )) &=-\varepsilon \frac 1{e^{i\theta
}-e^{-i\theta }},\;\;\;0<\theta <2\pi ,\;\theta \neq \pi , \\
\Theta _{\pi _0^\varepsilon }((-I_2)^ja(x)) &=(-1)^j\frac
1{|x-x^{-1}|},\;\;\;\varepsilon =\pm .
\end{align*}

\subsection{A correspondence between the discrete series}

\emph{Labels and characters.} For $k\geq 3$ odd, let
$\overline{\pi}{}^{+}_{k/2}$ denote the genuine holomorphic discrete series
representation of $\overline{SL(2,{\mathbb R})}$ with lowest
$\overline{K}$-type $\overline{r}(\theta)\mapsto e^{ik\theta/2}$ (its
$\overline K$-types are $e^{i(k/2+2j)\theta}$, $j\geq0$), and
$\overline{\pi}{}^{-}_{k/2}$ the antiholomorphic one, with
$\overline K$-types $e^{-i(k/2+2j)\theta}$. On the regular elliptic set
their characters are given by the $\overline K$-characters, summed as
geometric series (Harish-Chandra; compare \cite{Ad2}):
\begin{equation}\label{eq:mds-char}
\Theta_{\overline{\pi}{}^{+}_{k/2}}(\overline r(\theta))
=-\frac{e^{i(k/2-1)\theta}}{e^{i\theta}-e^{-i\theta}},\qquad
\Theta_{\overline{\pi}{}^{-}_{k/2}}(\overline r(\theta))
=\frac{e^{-i(k/2-1)\theta}}{e^{i\theta}-e^{-i\theta}},\qquad
\theta\notin\pi{\mathbb Z}.
\end{equation}
The same computation for $SL(2,{\mathbb R})$ gives the formulas for
$\Theta_{\pi_m}$ and $\Theta_{\pi_0^\pm}$ above; in particular $\pi_m$,
$m>0$, is holomorphic with lowest $K$-type $m+1$, and $\pi_0^+$ is the
holomorphic limit of discrete series. (Here $\varepsilon=+$ always means
holomorphic. The 1995 draft used Gelbart's labels, and its statement of
Lemma~\ref{lem:realds} differs from the one below by $\varepsilon\mapsto-\varepsilon$
together with an interchange of $+$odd and $-$odd, which is the effect of
reversing the orientation $\theta\mapsto-\theta$; we have not checked the
labels of \cite{Ge2} against (\ref{eq:mds-char}).)

\emph{Weyl integration formula.} With $\Phi$ normalized as in
\S\ref{sec:elliptic}, $dk=d\theta/2\pi$ on $K=\{r(\theta)\}$ and
$d\overline k=d\theta/4\pi$ on $\overline K=\{\overline r(\theta)\}$,
the elliptic terms of the Weyl integration formulas for $\Theta_{\pi}(f)$
and $\Theta_{\overline\pi}(\overline f)$ are
\begin{gather*}
\int_0^{2\pi }\Theta _{\pi}(r(\theta ))\Phi_f(r(\theta ))\,\frac{d\theta}{2\pi}
\qquad\text{and}\\
\int_0^{4\pi }
\Theta _{\overline{\pi }}(\overline{r}(\varphi ))
\Phi_{\overline{f}}(\overline{r}(\varphi ))\,\frac{d\varphi}{4\pi}
=\int_0^{2\pi }
\Theta _{\overline{\pi }}(\overline{r}(2\theta ))
\Phi_{\overline{f}}(\overline{r}(2\theta ))\,\frac{d\theta}{2\pi}.
\end{gather*}
Here the Weyl group of $SO(2)$ in $SL(2,{\mathbb R})$ is trivial. Since
\[
e^{2i\theta}-e^{-2i\theta}=(e^{i\theta}-e^{-i\theta})(e^{i\theta}+e^{-i\theta}),
\]
the substitution $\varphi=2\theta$ in (\ref{eq:mds-char}) produces the
factor $e^{i\theta}+e^{-i\theta}=2\cos\theta$ in the denominator. This
factor vanishes at $\theta=\pi/2,3\pi/2$, where $r(\theta)$ is regular
but $\overline r(2\theta)$ lies over the central element $-I$. We
therefore state the matching conditions without division.

\begin{definition}\label{def:realds-corr}
Let $f,\overline{f}$ have vanishing hyperbolic orbital integrals, and let
$E=\{\pi/2,\pi,3\pi/2\}$. We say that $f,\overline f$
\textbf{correspond evenly}, \textbf{$+$oddly}, or \textbf{$-$oddly}
\textbf{on the elliptic set} if, for all $\theta\in(0,2\pi)-E$,
\begin{align*}
\Phi_{\overline{f}}(\overline{r}(2\theta ))
&=(e^{i\theta }+e^{-i\theta })\,\Phi_f(r(\theta )) &&(\mathrm{even}),\\
\Phi_{\overline{f}}(\overline{r}(2\theta ))
&=e^{-i\theta}(e^{i\theta }+e^{-i\theta })\,\Phi_f(r(\theta )) &&(+\mathrm{odd}),\\
\Phi_{\overline{f}}(\overline{r}(2\theta ))
&=e^{i\theta}(e^{i\theta }+e^{-i\theta })\,\Phi_f(r(\theta )) &&(-\mathrm{odd}),
\end{align*}
respectively.
\end{definition}

\begin{remark}\label{rem:realds-consistency}
(1) As $\theta$ runs over $(0,2\pi)-E$, $\overline r(2\theta)$ runs over
all regular elements of $\overline K$, so each condition determines
$\Phi_{\overline f}$ on the regular elliptic set from $\Phi_f$. Since
$\Phi_f$ is smooth near $r(\pi/2)$ and $r(3\pi/2)$, which are regular,
each condition forces $\Phi_{\overline f}(\overline r(\varphi))\rightarrow0$
as $\varphi\rightarrow\pi$ or $3\pi$. (In Draft~1 the conditions were
written with the factor $(e^{i\theta}+e^{-i\theta})^{-1}$, which is
undefined at these points.)

(2) Since $\overline r(2\pi)=(I,-1)$ and $\Phi_{\overline f}$ is genuine,
replacing $\theta$ by $\theta+\pi$ shows that the even condition can hold
only if $\Phi_f(-t)=\Phi_f(t)$, and the odd conditions only if
$\Phi_f(-t)=-\Phi_f(t)$, for regular elliptic $t$. This is consistent
with Lemma~\ref{lem:realds}: the even matching sees the $\pi_m$ with
trivial central character ($m$ odd), the odd matchings those with
nontrivial central character ($m$ even).
\end{remark}

\begin{example}
Let $\Phi_f(r(\theta))=\sin(\theta)\,g(\theta)$ with $g$ smooth,
$2\pi$-periodic and supported in a small neighborhood of $\theta=\pi/2$
and of $\theta=3\pi/2=\pi/2+\pi$ (with $g(\theta+\pi)=-g(\theta)$, so that
$\Phi_f(-t)=\Phi_f(t)$). The even condition requires
$\Phi_{\overline f}(\overline r(\varphi))=2\cos(\varphi/2)\sin(\varphi/2)g(\varphi/2)
=\sin(\varphi)g(\varphi/2)$, which is smooth and vanishes at
$\varphi=\pi,3\pi$, as Remark~\ref{rem:realds-consistency}(1) requires.
\end{example}

\begin{lemma}\label{lem:realds}
Let $k\geq3$ be odd and $\varepsilon\in\{+,-\}$.

(a) If $f,\overline{f}$ correspond evenly, then
$\Theta _{\pi _m}(f)
=\Theta _{\overline{\pi }{}_{k/2}^\varepsilon }(\overline{f})$ for
$m=\varepsilon (k-2)$.

(b) If $f,\overline{f}$ correspond $+$oddly, then
$\Theta _{\pi _m}(f)
=\Theta _{\overline{\pi }{}_{k/2}^\varepsilon }(\overline{f})$ for
$m=\varepsilon (k-2)-1$, unless $k=3$ and
$\varepsilon =+$; and
$\Theta _{\pi _{0}^{+}}(f)=\Theta _{\overline{\pi }{}_{3/2}^{+}}(\overline{f})$.

(c) If $f,\overline{f}$ correspond $-$oddly, then
$\Theta _{\pi _m}(f)
=\Theta _{\overline{\pi }{}_{k/2}^\varepsilon }(\overline{f})$ for
$m=\varepsilon (k-2)+1$, unless $k=3$ and
$\varepsilon =-$; and
$\Theta _{\pi _{0}^{-}}(f)=\Theta _{\overline{\pi }{}_{3/2}^{-}}(\overline{f})$.
\end{lemma}

\begin{proof}
Since the hyperbolic orbital integrals vanish, only the elliptic terms
of the Weyl integration formulas contribute. By (\ref{eq:mds-char}) and
the factorization above, for $\theta\notin E$,
\[
\Theta_{\overline{\pi}{}^{+}_{k/2}}(\overline r(2\theta))
=-\frac{e^{i(k-2)\theta}}{(e^{i\theta}-e^{-i\theta})(e^{i\theta}+e^{-i\theta})},
\qquad
\Theta_{\overline{\pi}{}^{-}_{k/2}}(\overline r(2\theta))
=\frac{e^{-i(k-2)\theta}}{(e^{i\theta}-e^{-i\theta})(e^{i\theta}+e^{-i\theta})}.
\]
Multiplying by $\Phi_{\overline f}(\overline r(2\theta))$ as given by
Definition~\ref{def:realds-corr}, the factor $e^{i\theta}+e^{-i\theta}$
cancels, and the elliptic integrand of $\Theta_{\overline\pi}(\overline f)$
becomes $\Phi_f(r(\theta))$ times:
\[
\renewcommand{\arraystretch}{1.6}
\begin{array}{c|cc}
 & \overline{\pi}{}^{+}_{k/2} & \overline{\pi}{}^{-}_{k/2}\\ \hline
\text{even} & -\dfrac{e^{i(k-2)\theta}}{e^{i\theta}-e^{-i\theta}}
 & \dfrac{e^{-i(k-2)\theta}}{e^{i\theta}-e^{-i\theta}}\\
+\text{odd} & -\dfrac{e^{i(k-3)\theta}}{e^{i\theta}-e^{-i\theta}}
 & \dfrac{e^{-i(k-1)\theta}}{e^{i\theta}-e^{-i\theta}}\\
-\text{odd} & -\dfrac{e^{i(k-1)\theta}}{e^{i\theta}-e^{-i\theta}}
 & \dfrac{e^{-i(k-3)\theta}}{e^{i\theta}-e^{-i\theta}}
\end{array}
\]
Comparing with
$\Theta_{\pi_m}(r(\theta))=-{\rm sgn}(m)e^{im\theta}/(e^{i\theta}-e^{-i\theta})$
and $\Theta_{\pi_0^{\pm}}(r(\theta))=\mp1/(e^{i\theta}-e^{-i\theta})$, the
entries are $\Theta_{\pi_m}(r(\theta))$ with $m=\pm(k-2)$ (even),
$m=k-3$ or $-(k-1)$ ($+$odd), $m=k-1$ or $-(k-3)$ ($-$odd); when $k=3$
the entries $k-3$ and $-(k-3)$ are replaced by $\pi_0^+$ and $\pi_0^-$.
These agree with $\Theta_{\pi_m}(r(\theta))\Phi_f(r(\theta))$ for almost
all $\theta$, so the integrals agree.
\end{proof}

In every case the infinitesimal character of the representation of
$SL(2,{\mathbb R})$ is $|m|\in\{2\nu-1,2\nu,2\nu+1\}$, where
$\nu=k/2-1$ is the infinitesimal character of $\overline\pi{}^\pm_{k/2}$
(both normalized so that the trivial representation of $SL(2,{\mathbb R})$
has infinitesimal character $1$). Every discrete series or limit of
discrete series representation of $SL(2,{\mathbb R})$ occurs in at least
one of (a)--(c).

\begin{remark}
First, we recall from \cite{HC2}, page 660, that
\[
\Theta _{\pi _m}((-I_2)^ja(x))
=(-1)^{(m-1)j}\frac{e^{-|m\ln(x)|}}{|x-x^{-1}|}%
,\;\;m\neq 0,\;x>0,
\]
and from \cite{Ad2}, (15.7)(b), that
\[
\Theta _{\overline{\pi }_{k/2}^\varepsilon }(a(x))=\frac{%
sgn(x)^{-\varepsilon k/2}e^{-|(\frac k2-1)\ln(x)|}}{|x-x^{-1}|}
\]
(with the labels of \cite{Ad2}, which have not been compared with those
of Lemma~\ref{lem:realds}).
By definition, if $f,\overline{f}$ are corresponding on the hyperbolic set
then $(F_f)^{\wedge }(\chi )
=(F_{\overline{f}})^{\wedge }(\overline{\chi })$,
where $\chi ,\overline{\chi }$ are corresponding characters and
$F_\phi$ is the normalized orbital integral. Combining these
together, we see that without the assumption that
the hyperbolic orbital integrals vanish, the result is
inconsistent with the matching of principal series.
\end{remark}

%\begin{center}
%\textbf{REFERENCES}
%\end{center}

%[A] J. Adams, Lifting of characters on orthogonal and metaplectic groups,
%preprint, 1995

%[Fl] Y. Flicker, Automorphic forms on covering groups of GL(2), Inv. Math.,
%1980

%[G] S. Gelbart, \underline{Weil's representation and the spectrum of the
%metaplectic group}, SLNM 530, 1976

%[HC] Harish-Chandra, \underline{Collected Works, vol III}, Springer-Verlag

%[Hum] J. Humpheries, \underline{Linear algebraic groups}, Springer-Verlag Grad Texts

%[Mum] D. Mumford, \underline{The red book}, Springer LNM, 1358, 1988

%[V] V. S. Varadarajan, An introduction to harmonic analysis on semi-simple
%Lie groups, Cambridge U. Press, 1989
%% ===== end aabook2c.tex =====


%% ===== begin aabook2d.tex (revised) =====
\section{On the correspondence of d. s. for
$\overline{SL(2,F)}$, $p>2$}\label{sec:padic-ds}

We have described several correspondences
$c_i:\overline{\pi}\longmapsto \pi$
between representations of
$\overline{SL(2,F)}$ and those of $SL(2,F)$. They are meant to be
consistent with a ``matching''
$c_i^*:\overline{f}\longmapsto f$
of functions, in the sense that the characters match:
$\Theta_{\overline{\pi}}(\overline{f})=\Theta_{\pi}(f)$, where
$\pi=c_i(\overline{\pi})$ and $f=c_i^*(\overline{f})$.
In the 1995 draft it was asserted that every element $\pi$ of the
unitary dual of $G(F)$ occurs in the image of at
least one of the $c_i$; in view of Remarks~\ref{rem:sc-caveat}
and~\ref{rem:savin-caveat}, this is not established here.

Let $\tau:G(F)\rightarrow \overline{G(F)}$ denote the map
$g\longmapsto \tilde{g}$ of \cite{Fl}.

\subsection{The Iwahori algebra}

Let $B$ denote the Iwahori subgroup of $G(F)$ and assume $p$ is odd.
Savin \cite{Sav} studied the genuine Iwahori--Hecke algebras of
metaplectic covers of split groups. The 1995 draft used his work in the
following form.

\begin{claim}[1995 draft, attributed to Savin \cite{Sav}; unverified, see
Remark~\ref{rem:savin-caveat}]\label{claim:savin} There is an isomorphism
of algebras
${\cal H}(\overline{G}//B)\rightarrow {\cal H}(G//B)$.
\end{claim}

\begin{claim}[1995 draft; depends on Claim~\ref{claim:savin}]
The isomorphism
${\cal H}(\overline{G}//B)\rightarrow {\cal H}(G//B)$
induces a set-theoretic
bijection ${\cal H}(\overline{G}//B)^\wedge
\rightarrow {\cal H}(G//B)^\wedge$ on the sets of equivalence
classes of irreducible (genuine, in the case of $\overline{G}$)
representations with a $B$-fixed vector.
Furthermore, if
\[
\begin{array}{ccc}
{\cal H}(G//B)&\rightarrow &{\cal H}(\overline{G}//B)\\
f & \longmapsto & \overline{f}
\end{array}
\]
and
\[
\begin{array}{ccc}
{\cal H}(\overline{G}//B)^\wedge &\rightarrow &{\cal H}(G//B)^\wedge\\
\overline{\pi}&\longmapsto &\pi
\end{array}
\]
then $\tr(\overline{\pi}(\overline{f}))=\tr(\pi(f))$.
Furthermore, if $\overline{\pi}$ is a discrete series representation
then so is $\pi$.
\end{claim}

\begin{remark}\label{rem:savin-caveat}
(Added 2026.) In the later literature on Hecke algebra correspondences
for metaplectic groups (see \cite{TW} and the references there), the
linear group paired with the metaplectic cover of $SL(2)=Sp(2)$ is
$SO(3)\cong PGL(2)$, not $SL(2)$, and the resulting bijection of
Iwahori-spherical representations is with representations of $PGL(2,F)$.
The two claims above, with $G=SL(2)$, should therefore be
checked against \cite{Sav} before being used. In particular, the 1995
draft's assertion that this isomorphism is associated with an identity
of orbital integrals rested on the argument of the 1995 version of
\S\ref{sec:dyadic}, which has been withdrawn.
Neither claim is used in Chapter~\ref{ch:global}.
\end{remark}

The 1995 draft intended this corollary to be used in the comparison of
the special representations at non-archimedean places $v$ not dividing
$2$.

%% ===== end aabook2d.tex =====


%% ===== begin aabook2e.tex =====
\section{On the character of the Weil representation}\label{sec:weil}

We formally compute the character of the ``even'' and ``odd'' parts
of the Weil representation on the ``big cell'' of the Bruhat decomposition
by integrating the kernel along the diagonal. (This integral does not
converge absolutely...). We then state some consequences for local matching
of representations. This section shall not be needed later.

Here $F$ may be either p-adic or real. If $F$ is p-adic we shall
sometimes impose the assumption that $p\not= 2$.

\subsection{A Bruhat decomposition}

Let $G=SL(2)$, let
\[
h(x)=\left(
\begin{array}{cc}
x & 0 \\
0 & x^{-1}
\end{array}
\right)
,\ \ \ \
n(x)=\left(
\begin{array}{cc}
1 & x \\
0 & 1
\end{array}
\right)
\]
and let
\[
w:=\left(
\begin{array}{cc}
0 & -1 \\
1 & 0
\end{array}
\right),\ \ \
r(\theta )=\left(
\begin{array}{cc}
cos(\theta) & -\sin(\theta) \\
\sin(\theta) & \cos(\theta)
\end{array}
\right) .
\]

Let $F$ be either ${\mathbb{R}}$ or a p-adic field.
If $F$ is p-adic, let $|\ |$ denote the norm which is normalized
so that $|\varpi|=q^{-1}$, where $q$ denotes the number of elements of the
residue field, and let ${\cal O}_F$ denote the ring of integers.
Let $(\ ,\ )$ denote
the quadratic Hilbert symbol and let $\psi$ denote a non-trivial
additive character of $F$. If $F$ is real, take $\psi(x)=\exp(\pi ix)$
and if $F$ is p-adic then assume $\psi$ has conductor ${\cal O}_F$.

Let $\beta: G(F)\times G(F)\rightarrow \mu_2=\{\pm 1\}$
denote the cocycle
\[
\beta(g,g')=({\frac{x(gg')}{x(g)}}, {\frac{x(gg')}{x(g')}} ),
\]
where $g,g'\in G(F)$ and
\[
x\left(
\begin{array}{cc}
a & b \\
c & d
\end{array}
\right)
=
\left\{
    \begin{array}{c}
    c,\ \ \ c\not= 0,\\
    d,\ \ \ {\rm otherwise}.
    \end{array}
\right.
\]
We define a group structure $\overline{G(F)}$ on
$G(F)\times \mu_2$ by
\[
(g,\zeta)\cdot (g',\zeta')=(gg',\beta(g,g')\zeta\zeta'),
\]
$g,g'\in G(F),\ \zeta,\zeta'\in \mu_2$.
The following simple fact has appeared in several places:

\begin{lemma} If
\[
g=\left(
\begin{array}{cc}
a & b \\
c & d
\end{array}
\right)
\]
has determinant 1 then
\[
(g,(-c,-1))=(n(a/c),1)\cdot (w,1)\cdot (h(-c),1)\cdot (n(d/c),1),
\]
provided $c\not= 0$.
\end{lemma}

\begin{proof} We have
\[
n(a/c)\cdot w =
\left(
\begin{array}{cc}
-a/c & 1 \\
-1 & 0
\end{array}
\right)
\]
so $\beta(n(a/c),w)
=({\frac{x(n(a/c)w)}{x(n(a/c))}},{\frac{x(n(a/c)w)}{x(w)}})
=(-1,1)=1$.
Also,
\[
n(a/c)\cdot w h(-c) =
\left(
\begin{array}{cc}
a & -1/c \\
c & 0
\end{array}
\right)
\]
so $\beta(n(a/c)w,h(-c))
=({\frac{x(n(a/c)wh(-c))}{x(n(a/c)w)}},{\frac{x(n(a/c)wh(-c))}{x(h(-c))}})
=(-c,-c^2)=(-c,-1)$.
Finally,
\[
n(a/c)\cdot w h(-c)n(d/c) =
\left(
\begin{array}{cc}
a & b \\
c & d
\end{array}
\right)=g
\]
so $\beta(n(a/c)wh(-c),n(d/c))
=({\frac{x(g)}{x(n(a/c)wh(-c))}},{\frac{x(g)}{x(n(d/c))}})
=(1,c)=1$.
\end{proof}

\subsection{The Weil representation}

Let $\omega$, $q(x)$, and $\gamma(q)$ be as in \cite{Fl}, \S 2.2. See also
\cite{Ge2}, page 36, and \cite{Rao}, Appendix.

For $\phi \in L^1(F)$, let
\[
\hat{\phi}_q(x)=\int_F \phi(y)\psi(2qxy)dy,
\]
where $dy$ is normalized as in \cite{Ge2}, ch 2, so that this extends to
an isometry of $L^2(F)$. This implies that $dy$ depends on $\psi, q$
in a way we will make explicit later.
(In the case where $F={\mathbb{R}}$, if $\psi(x)=\exp(\pi i x)$ then
$dy$ is $\sqrt{|q|}$ times Lebesgue measure.)

Let $\epsilon \in \{\pm\}$ and let
\[
L^2(F)_\epsilon=
\{\phi \ |\ \phi(-x)=\epsilon \phi(x),\ \ \forall x\in F\}.
\]
For $\phi\in L^2(F)_\epsilon$, we have
\[
\int_F \phi(x)\psi(2qxy)dx =
\int_{F/\{\pm \}}
\phi(x) (\psi(2qxy)+\epsilon \psi(-2qxy))\overline{dx},
\]
if $dx=2\overline{dx}$.

We define the Weil representation $r=r_{\psi,q}$ of
$\overline{G(F)}$ on $L^2(F)$ as in \cite{Fl}, \S 2.2, so

(i) $r(n(b),\zeta)\phi(x)=\zeta\cdot \psi(bq(x))\phi(x)$,

(ii) $r(w,\zeta)\phi(x)
=\zeta\cdot \gamma(q)\hat{\phi}_q(x)$,

(iii) $r(h(a),\zeta)\phi(x)=\zeta\cdot |a|^{1/2}\omega(a)\phi(ax)$.

We define the Weil representation $r_\epsilon=r_{\psi,q,\epsilon}$ of
$\overline{G(F)}$ on $L^2(F)_\epsilon$ by

(i) $r_\epsilon(n(b),\zeta)\phi(x)=\zeta\cdot \psi(bq(x))\phi(x)$,

(ii) $r_\epsilon(w,\zeta)\phi(x)
=\zeta\cdot {\frac{1}{2}}\gamma(q)(\hat{\phi}_q(x)
+\epsilon \hat{\phi}_{-q}(x))$,

(iii) $r_\epsilon(h(a),\zeta)\phi(x)=\zeta\cdot |a|^{1/2}\omega(a)\phi(ax)$.

We will assume that $r_\epsilon$ agrees with the restriction of
$r$ to $L^2(G)_\epsilon$.

We have
\[
r_\epsilon(g,1)=(-c,-1)
r_\epsilon(n(a/c),1)\cdot r(w,1)\cdot r_\epsilon(h(-c),1)
\cdot r_\epsilon(n(d/c),1),
\]
where $g$ is as in the lemma above. This implies
\[
r_\epsilon(g,1)\phi(x)=(-c,-1)|c|^{-1/2}\omega(-c)
\psi_q({\frac{a+d}{c}}q(x))
\gamma(q)(\hat{\phi}_q(-x/c)+
\epsilon \hat{\phi}_{-q}(-x/c)).
\]

Since $r_+$ is even and $r_-$ is odd, we point out that
the characters $\Theta_\epsilon$ of each of these
representations may be obtained from the
character of the sum (i.e., the character of $r$)
by virtue of the trick
\[
\Theta_+(-g,1)+\Theta_-(-g,1)=
\Theta_+(g,1)-\Theta_-(g,1).
\]
The character $\Theta_+(g)+\Theta_-(g)$ is determined in \cite{H}.
We shall take another approach.

\subsection{Integrating the kernel}

For $\phi\in L^2(F)_\epsilon$, we have
\[
r_\epsilon(g,1)\phi(x)=
\int_{F}
k_\epsilon (x,y)\phi(y)dy,
\]
where the kernel function (which also depends on $g$)
\[
\begin{array}{c}
k_\epsilon (x,y)=
(-c,-1)|{\frac{1}{c}}|^{1/2}
\omega(-c)\psi({\frac{a+d}{c}}q(x))
\gamma(q)(\psi_{-q}({\frac{2}{c}}xy)+\epsilon \psi_q({\frac{2}{c}}xy))\\
=(-c,-1)|{\frac{1}{c}}|^{1/2}
\omega(-c)\gamma(q)
(\psi_{-q}(-{\frac{a+d-2}{c}}xy)+\epsilon \psi_q({\frac{a+d+2}{c}}xy))
\end{array}
\]
Note that $k_\epsilon(x,x)$ may be regarded as a function on either
$F$ or $F/\{\pm 1\}$.
Also, note if $\phi\in L^2(F)_\epsilon$ then
the functions $x\longmapsto k_\epsilon(x,y)\phi(x)$,
$y\longmapsto k_\epsilon(x,y)\phi(y)$ are
functions on $F/\{\pm 1\}$.

The restriction $r_\epsilon$ of $r$ to $L^2(F)_\epsilon$
is an irreducible admissible representation (see, for example,
\cite{Fl}, \S 2.2, or \cite{Ge2}). For each genuine $f\in C_c^\infty(\overline{G(F)})$,
the trace of the operator $r_\epsilon(f)$ exists and there is an
anti-genuine function $\Theta_\epsilon(g)$ on $\overline{G(F)}$
such that
\[
\tr(r_\epsilon(f))=\int_{\overline{G(F)}} \Theta_\epsilon (g)f(g)dg.
\]

We shall assume the following

\begin{conjecture}
On the regular elements of the ``big cell''
$N(F)wN(F)\overline{A(F)}$, the character $\Theta_\epsilon$
is given by the principal
value integral of the kernel function along the diagonal:
\[
\Theta_\epsilon(g,1)=
\lim_{N\rightarrow \infty}
\int_{\{x\in F\ |\ |x|<N\}}
k_\epsilon(x,x)dx.
\]
Here $dx$ is the normalized Haar measure as above.
\end{conjecture}

\begin{remark} This conjecture implies
\[
\begin{array}{c}
\Theta_\epsilon(g,1)\\
=(-c,-1)|{\frac{1}{c}}|^{1/2}\omega(-c)\gamma(q)
\lim_{N\rightarrow \infty}
\int_{\{x\in F\ |\ |x|<N\}}
{\frac{\psi_{-q}(-{\frac{a+d-2}{c}}x^2)
+\epsilon \psi_q({\frac{a+d+2}{c}}x^2)}{2}}dx.
\end{array}
\]
\end{remark}

Let
\[
\int_{\mathbb{R}} \exp(aix^2)dx
=\lim_{N\rightarrow \infty}
\int_{\{x\in {\mathbb{R}}\ |\ |x|<N\}} \exp(aix^2)dx.
\]
The following result is elementary and well-known:

\begin{lemma}\label{lem:gauss} Let $d_Lx$ denote the usual Lebesgue measure.
If $F=\mathbb{R}$ and $a\not= 0$ is real then
\[
\int_{F} \exp(aix^2)d_Lx =\exp(\pi i sgn(a)/4)\sqrt{\frac{\pi}{|a|}}.
\]
\end{lemma}

\begin{remark}
If $<\ ,\ >$ is the usual inner product on
$F^n$ then
\[
\int_{F^n} \exp(ai<\vec{x},\vec{x}>)d_L x =({\frac{\pi}{|a|}})^{n/2}.
\]
Using simple generalizations like this it is hoped that it will
be possible to generalize some of the results below to
$\overline{Sp(2n,{\mathbb{R}})}$.
\end{remark}

\begin{proof} We may assume without loss of generality that $a>0$.
Let $C_R=C_{R,a}\cup C_{R,b}\cup C_{R,c}$
denote the sector of the circle of radius $R$,
consisting of $z=re^{i\theta}$ with (a) $\theta=0,\ 0\leq r\leq R$,
(b) $\theta=\pi/4,\ 0\leq r\leq R$, and (c) $0\leq \theta \leq \pi/4,\ r=R$.
We have
\[
\begin{array}{c}
\int_{C_{R,c}} e^{iaz^2}dz=
Ri\int_0^{\pi/4} e^{R^2(-\sin 2\theta +i\cos 2\theta )}e^{i\theta} d\theta\\
<< R\int_0^{\pi/4} e^{-R^2\sin 2\theta }d\theta\\
<< R\int_0^{\epsilon} e^{-2R^2x}dx +Re^{-\epsilon R^2}
<<ln(R)/R,
\end{array}
\]
picking $\epsilon =2ln(R)/R^2$.
This allows us to use a standard residue theorem argument
applied to $C_R$ to obtain
\[
\int_F e^{iax^2}dx=\exp(\pi i/4)\int_F e^{-ay^2}dy.
\]
The result now follows from the fact that, for $a>0$,
\[
\int_0^\infty e^{-ax^2}dx={\frac{1}{2}}\sqrt{\frac{\pi}{a}}.
\]
\end{proof}

Let $\tau (x)=\exp (\pi i q x)$, so $\tau=\psi_q$,
and let $d_\tau x$ denote
the Haar measure normalized as in \cite{Ge2}, ch 2. If
\[
\gamma(\tau)=
\int_{\mathbb{R}} \tau(x^2)d_\tau x
\]
then $\gamma(\tau) =\gamma(q)$ (in the notation of \cite{Fl}, \S 2.2) and
\[
\gamma(\tau)=\exp ({\frac{\pi i}{4}} sgn(q)),
\]
by \cite{Ge2}, page 36. Let $\tau_a(x)=\tau(ax)$. Then
\[
\int_{\mathbb{R}} \tau_a(x^2)d_\tau x
=|a|^{-1/2}\exp ({\frac{\pi i}{4}} sgn(q)).
\]
In particular,
\[
\begin{array}{c}
\gamma(q)
\int_{\mathbb{R}}
(\psi_{-q}(-{\frac{a+d-2}{c}}x^2)+
\epsilon\psi_q({\frac{a+d+2}{c}}x^2)d_\tau x\\
=\gamma(q)\gamma(-q)|{\frac{a+d-2}{c}}|^{-1/2}+
\epsilon \gamma(q)^2|{\frac{a+d+2}{c}}|^{-1/2}.
\end{array}
\]

In the p-adic case, let
\[
\kappa(\tau)=\lim_{N\rightarrow - \infty}
\int_{\varpi^N{\cal O}_F}
\tau(u^2) du.
\]
For a $fixed$ Haar measure $du$, we have
$\kappa(\tau_{a^2})=|a|^{-1}\kappa(\tau)$, for all
$a\in F^\times$. Now normalize the measure $du$ as in \cite{Mo3}
(so $du=d_\tau u$ depends on $\tau$) and
call the result $\kappa_{nor}(\tau)$. It is known
that (see for example \cite{Mo3})
\[
\kappa_{nor}(\tau)=
\left\{
    \begin{array}{c}
    1,\ \ \ 2|f(\tau),\\
    G(\tau),\ \ \ {\rm otherwise}.
    \end{array}
\right.
\]
where the conductor of $\tau$ is given by
$\varpi^{f(\tau)}{\cal O}_F$ and where $G(\tau)$ denotes the Gauss sum
normalized so as to be a $8^{th}$ root of unity as in \cite{Mo3}.
Now, the Haar measure $d_\tau u$ has been chosen to depend on
$\tau$ in such a way that $vol({\cal O}_F,d_\tau u)=q^{f(\tau)/2}$.
This is the measure for which $(\hat{\phi}_q)^\wedge_q(x)=\phi(-x)$,
by \cite{Shal}.
Note that $f(\tau_a)=f(\tau)-\nu(a)$, where
$\nu : F^\times \rightarrow {\mathbb{Z}}$ is the normalized
valuation. Putting these together gives
$\kappa_{nor}(\tau_{a^2})=\kappa_{nor}(\tau)$. It is also known that
$\kappa_{nor}(\tau_{a})=({\frac{a}{\varpi}})\kappa_{nor}(\tau)$, for
$a\in {\cal O}_F^\times$, where $\left(\frac{\cdot}{\varpi}\right)$ is the
Legendre symbol. Of course,
\[
vol({\cal O}_F,d_\tau u)^{-1}
\kappa_{nor}(\tau)
=vol({\cal O}_F,du)^{-1}\kappa(\tau),
\]
so it is easy to convert from one to the other.
We assume that $du$ is the measure in the conjecture above.

\begin{proposition} (Assume the conjecture above.)
If $c\not= 0$ then
\[
\Theta_\epsilon (g,1)=
(-c,-1)||{\frac{q}{c}}|^{1/2}\omega(-c)
{\frac{\gamma(q)\gamma(-q)\kappa(\psi_\alpha)+
\epsilon \gamma(q)^2\kappa(\psi_\beta)}{2}},
\]
where $\alpha={\frac{a+d-2}{c}}q$, $\beta={\frac{a+d+2}{c}}q$.
\end{proposition}

\begin{proof} This follows by putting the results above into the conjecture
and the remark following it.
\end{proof}

\begin{corollary} (Assume the conjecture above)
If $F={\mathbb{R}}$ and
$g=r(\theta)$ is elliptic regular then
\[
\begin{array}{c}
\Theta_\epsilon (g,\zeta_g)=
\zeta_g
{\frac{(-\sin(\theta),-1)\omega(-\sin(\theta))}{2}}
({\frac{\gamma(q)\gamma(-q)}{|\det(g-1)|^{1/2}}}
+\epsilon {\frac{\gamma(q)^2}{|\det(g+1)|^{1/2}}})\\
=\zeta_g
{\frac{(-\sin(\theta),-1)\omega(-\sin(\theta))}{2}}
({\frac{1}{2|\sin(\theta /2)|^{1/2}}}
+\epsilon {\frac{\gamma(q)^2}{2|\cos(\theta/2)|^{1/2}}}).
\end{array}
\]
for $\zeta_g=sgn(\cos(\theta))\in \mu_2$.
\end{corollary}

\begin{remark}
Note that
$(-\sin(\theta),-1)=-sgn(\sin(\theta))$. According to \cite{Ge2}, page
36, $\gamma(q)=\exp(-{\frac{i\pi}{4}}sgn(q))$ is the Weil index
associated to the quadratic character
$x\longmapsto \exp(\pi i qx^2)$ (here $\gamma(q)$ is as in \cite{Fl}, \S 2.2;
apparently, one should replace $q$ by $-q$ in \cite{Fl} to get formulas
consistent with \cite{Ge2}, ch 2). In the notation of \cite{Fl}, we have
\[
\omega(a)=\gamma(q)/\gamma(aq)
=
\left\{
\begin{array}{c}
1,\ \ \ \ \ \ \ \ a>0,\\
-i,\ \ \ \ a<0,\ \ q>0,\\
i,\ \ \ \ a<0,\ \ q<0.
\end{array}
\right.
\]
Using this, the corollary agrees with Adams' formula in the case
$0<\theta <\pi$. Since $r(-\theta)=wr(\theta)w^{-1}$,
we see that this also agrees with Adams' formula for $-\pi<\theta <0$.
\end{remark}

\begin{remark}
This should be compared with the formulas in \cite{Ad1}
(see especially Proposition 4.5, 4.7, and equation (5.5)),
where one finds that
\[
\theta(\psi_\epsilon)_{even}(g,\zeta_g)
={\frac{1}{4}} ({\frac{1}{\cos(\theta/2)}}
+\epsilon {\frac{i}{\sin(\theta/2)}}),
\]
and
\[
\theta(\psi_\epsilon)_{odd}(g,\zeta_g)
={\frac{1}{4}} (-{\frac{1}{\cos(\theta/2)}}
+\epsilon {\frac{i}{\sin(\theta/2)}}),
\]
for $0<\theta <4\pi$, $\theta\not= \pi,2\pi,3\pi$.
\end{remark}

\begin{proof} Since $a+d\pm 2=\det(g\pm 1)$, this follows from
the proposition above and Lemma~\ref{lem:gauss}.
\end{proof}

\begin{corollary} (Assume the conjecture above)
If $F$ is p-adic and
$g=\left(
\begin{array}{cc}
a & \eta b \\
b & a
\end{array}
\right)$ is elliptic regular ($\eta\in F$ is a non-square)
then
\[
\begin{array}{c}
\Theta_\epsilon (g,1)\\
=(-b,-1)\omega(-b)|2q{\frac{\sqrt{\eta}}{\lambda_1 -\lambda_2}}|^{1/2}
{\frac{\gamma(q)\gamma(-q)\kappa(\psi_\alpha)+
\epsilon \gamma(q)^2\kappa(\psi_\beta)}{2}}\\
=C(g,q,du)
({\frac{\gamma(q)\gamma(-q)\kappa_{nor}(\psi_\alpha)}{|\det(g-1)|^{1/2}}}+
\epsilon {\frac{\gamma(q)^2\kappa_{nor}(\psi_\beta)}{|\det(g+1)|^{1/2}}}),
\end{array}
\]
where $C(g,q,du)=vol({\cal O}_F,du)(-b,-1)\omega(-b)/2$,
where
\[
\alpha=2q{\frac{a-1}{b}}=
{\frac{2q\sqrt{\eta}}{\lambda_1-\lambda_2}}\det(g-1),
\ \ \ \
\beta=2q{\frac{a+1}{b}}
={\frac{2q\sqrt{\eta}}{\lambda_1-\lambda_2}}\det(g+1),
\]
and where $\lambda_1,\lambda_2$ are the eigenvalues of $g$.
\end{corollary}

\begin{remark}
This should be compared with the formulas in \cite{H}.
\end{remark}

\subsection{Matchings for non-tempered representations}

Let $r_{even}$ denote the even part of the Weil representation
on $\overline{GL(2,F)}$ in \cite{Fl}, \S 2.2 (where it is denoted by
$r_\mu$, $\mu$ an even character which we take here to be the
trivial character).

\subsubsection{The real case}

Let $F={\mathbb{R}}$. Let $\overline{\pi}_{+,nt},\overline{\pi}_{-,nt}$
denote the two irreducible, genuine, non-tempered, unitary representations
of $\overline{G(F)}$. In the notation of \cite{Ad1}, \cite{Ge2}, we have
\[
\overline{\pi}_{+,nt}
=\overline{\pi}^+_{1/2}=\omega(\psi_+)_{even},
\ \ \ \ \
\overline{\pi}_{-,nt}
=\overline{\pi}^-_{1/2}=\omega(\psi_-)_{even}.
\]
On the Grothendieck group of the unitary dual of $\overline{G(F)}$,
we have
\[
r_{even}|_{\overline{G(F)}}=
\overline{\pi}_{+,nt}+\overline{\pi}_{-,nt}.
\]
Let $\pi_{nt}$ denote the (non-tempered) trivial representation of $G(F)$.

Using the notation above and that
of \cite{Ad1}-\cite{Ad2}, we have the following additional
identities in the Grothendieck group
of the unitary dual of $\overline{G(F)}$ or $G(F)$
(whichever is appropriate to the context):
\begin{equation}\label{eq:groth}
\begin{array}{c}
\pi'_{PS}(1/2,1/2)=\overline{\pi}_{+,nt}+\pi'_{DS}(-,1/2),\\
\pi'_{PS}(-1/2,1/2)=\overline{\pi}_{-,nt}+\pi'_{DS}(+,1/2),\\
\pi_{PS}(0,1)=\pi_{nt}+\pi_{DS}(+,1)+\pi_{DS}(-,1),\\
\pi_{PS}(0,-1)=\pi_{nt}+\pi_{DS}(+,1)+\pi_{DS}(-,1).
\end{array}
\end{equation}
Here $\pi'_{DS}(\pm,1/2)$ are the genuine discrete series with
$\overline K$-types $\pm(3/2+2j)$, $j\ge0$, i.e.\
$\overline\pi{}^{\pm}_{3/2}$ in the notation of \S\ref{sec:realds}, and
$\pi_{DS}(\pm,1)=\pi_{\pm1}$. (For instance, the even oscillator
representation $\overline\pi_{+,nt}$ has $\overline K$-types
$1/2+2j$, $j\geq0$, and the remaining $\overline K$-types $-3/2-2j$ of the
principal series containing it are those of $\overline\pi{}^-_{3/2}$.)

\begin{lemma}\label{lem:ds-3/2} If
$f,\overline{f}$ correspond evenly on the elliptic set in
the sense of Definition~\ref{def:realds-corr}, then
$\Theta_{\pi'_{DS}(\pm,1/2)}(\overline{f})
=\Theta_{\pi_{DS}(\pm,1)}(f)$.
\end{lemma}

\begin{proof}
This is Lemma~\ref{lem:realds}(a) with $k=3$.
\end{proof}

\begin{lemma} If $f,\overline{f}$ correspond evenly on the elliptic set
in the sense of Definition~\ref{def:realds-corr}, then
\[
\Theta_{\overline{\pi}_{+,nt}}(\overline{f})
+\Theta_{\overline{\pi}_{-,nt}}(\overline{f})
=\Theta_{\pi_{nt}}(f).
\]
\end{lemma}

\begin{proof}
(Added in Draft~2.) The characters of principal series representations
vanish on the regular elliptic set, and the hyperbolic orbital
integrals of $f$ and $\overline f$ vanish. Hence, by (\ref{eq:groth}) and
the Weyl integration formula,
$\Theta_{\overline{\pi}_{\pm,nt}}(\overline f)=-\Theta_{\pi'_{DS}(\mp,1/2)}(\overline f)$
and
$\Theta_{\pi_{nt}}(f)=-\Theta_{\pi_{DS}(+,1)}(f)-\Theta_{\pi_{DS}(-,1)}(f)$.
The result follows from Lemma~\ref{lem:ds-3/2}. (The identities
(\ref{eq:groth}) are quoted from \cite{Ad1}, \cite{Ge2}, and have not
been rechecked here.)
\end{proof}

\begin{remark}
Indeed, it seems likely that \cite{Fl} implies that
if $f,\overline{f}$ correspond evenly on the hyperbolic set in
the sense of \S\ref{sec:charsA}, then
\[
\Theta_{\overline{\pi}_{+,nt}}(\overline{f})
+\Theta_{\overline{\pi}_{-,nt}}(\overline{f})
=\Theta_{\pi_{nt}}(f).
\]
\end{remark}

\subsubsection{The $p$-adic case}

Let $F$ be p-adic. Let $r_{odd}$ denote the odd part of some Weil
representation (associated to some odd character)
on $\overline{GL(2,F)}$ as in \cite{Fl}, \S 2.2. This is an
irreducible genuine supercuspidal representation of
$\overline{GL(2,F)}$. Its restriction to
$\overline{SL(2,F)}$ decomposes into $4=[F^\times:F^{\times 2}]$
irreducible genuine supercuspidal components $\overline{\pi}_{sc,i}$
\cite{M3}.

Let $r_{even}$ denote the even part of the Weil
representation associated to the trivial character on
$\overline{GL(2,F)}$ as in \cite{Fl}, \S 2.2. This is an
irreducible genuine representation of
$\overline{GL(2,F)}$, which is not supercuspidal. (Draft~1 called it
supercuspidal.) Its restriction to
$\overline{SL(2,F)}$ decomposes into $4=[F^\times:F^{\times 2}]$
irreducible genuine infinite dimensional
non-tempered components $\overline{\pi}_{nt,i}$.

Let $\sigma$ denote the special representation of $SL(2,F)$
and $\pi_{nt}$ the trivial representation.

The 1995 draft derived the following from \cite{Fl} and the fact that the
restriction of the special (Steinberg) representation of $GL(2,F)$
to $SL(2,F)$ is irreducible.

\begin{claim}[1995 draft; unverified] (a) If $f, \overline{f}$ correspond evenly on the elliptic set
then
\[
\sum_i \Theta_{\overline{\pi}_{sc,i}}(\overline{f})
=\Theta_\sigma(f).
\]
(b) If $f,\overline{f}$ correspond evenly on the elliptic set
then
\[
\sum_i \Theta_{\overline{\pi}_{nt,i}}(\overline{f})
=\Theta_{\pi_{nt}}(f).
\]
(c) Let $\overline{\sigma}_i$ denote the complement of
$\overline{\pi_{nt,i}}$ (in the associated full unitarily
induced representation). If $f,\overline{f}$ correspond
evenly on the elliptic set then
\[
\sum_i \Theta_{\overline{\sigma}_{i}}(\overline{f})
=-\Theta_{\pi_{nt}}(f).
\]
\end{claim}

\begin{remark} If $f,\overline{f}$ correspond oddly on the elliptic set
then we have the supercuspidal
correspondence described in Chapter~\ref{ch:sc} (see Remark~\ref{rem:sc-caveat}).
\end{remark}

%\begin{center}
%\textbf{REFERENCES}
%\end{center}

%[A] J. Adams, ``Character of the oscillator representations'',
%preprint, 1995 (to appear in Isr. J. Math.)

%[A'] ---, Notes on the Flicker correspondence

%[A''] ---, ``Lifting of characters on orthogonal and metaplectic groups'',
%preprint, 1995

%[Fl] Y. Flicker, ``Automorphic forms on covering groups of GL(2)'',
%Inv. Math., 1980

%[G] S. Gelbart, \underline{Weil's representation and the spectrum of the
%metaplectic group}, SLNM 530, 1976

%[H] R. Howe, unpublished

%[HC] Harish-Chandra, \underline{Collected Works, vol III}, Springer-Verlag

%[Hum] J. Humpheries, \underline{Linear algebraic groups},
%Springer-Verlag Grad Texts

%[Mum] D. Mumford, \underline{The red book}, Springer LNM, 1358, 1988

%[J] D. Joyner, book

%[J'] ---, Notes on the ds correspondence for $\overline{SL(2,F)}$,
%$F$ p-adic

%[J''] ---, Notes on the ds correspondence for
%$\overline{SL(2,{\mathbb{R}})}$

%[Man] D. Manderscheid, ``Supercuspidal representations and the
%theta correspondence'', preprint

%[M] C. Moen, ``The dual pair $(U(1),U(1))$ over a p-adic field'',
%Pac J Math, vol 158, 1993

%[R] R. Rao, ``On some explicit formulas in the theory of Weil
%representations'', Pac J Math \underbar{157}(1993)335-371

%[S] J. Shalika, ``Representations of the two by two unimodular group
%over local fields'', IAS

%[V] V. S. Varadarajan, An introduction to harmonic analysis on semi-simple
%Lie groups, Cambridge U. Press, 1989
%% ===== end aabook2e.tex =====



%% ===== begin aabook5.tex =====
\chapter{The trace formulas}\label{ch:tf}

We recall the trace formula in the case of $SL(2)$, as formulated by
Arthur. Then we state the metaplectic analog of the formula which
can be derived analogously. We shall be very brief since there are
already several expository works explaining the ideas needed here
(e.g., \cite{A15}, \cite{MW}).

\section{The trace formula for $SL(2)$}

We shall merely write down the terms, with some definitions, from the
work of Arthur \cite{A2}, \cite{A3}, \cite{A7}.

\subsection{The definition of $J_M(u,f)$}

Let $f\in C_c^\infty (G({\mathbb {A}}_S))$ and let
$\gamma \in M({\mathbb {A}}_S)$ have Jordan decomposition
$\gamma =\sigma u$. Suppose that
$M({\mathbb {A}}_S)_\gamma =G({\mathbb {A}}_S)_\gamma $.
Assume that $S$ is a ``sufficiently large'' finite set which contains
the archimedean places. In this case, define
\begin{equation}
J_M(\gamma ,f)=|D(\gamma )|^{1/2}
\int_{M({\mathbb {A}}_S)\backslash G({\mathbb {A}}_S)}
\int_{{\cal O}_\gamma (M({\mathbb {A}}_S))}f(x^{-1}\mu x)v_M(x)d\mu dx,
\end{equation}
where ${\cal O}_\gamma (M({\mathbb {A}}_S))
=\{m^{-1}\gamma m\;m\in M({\mathbb {A}}_S))$.
Here $v_M$ is the ``volume'' of the convex hull of the points
$\{H_B(x),H_{B^{-}}(x)\}$ in ${\mathfrak a}$. If $x=nk$ then
$H_B(x)=0$, so $v_M(nk)=|H_{B^{-}}(x)|=|\log\delta (wx)|$.
In any case, this integral can be written as
\[
J_M(\gamma ,f)=|D(\gamma )|^{1/2}
\int_{G({\mathbb {A}}_S)_\gamma \backslash G({\mathbb{A}}_S)}
f(x^{-1}\gamma x)v_M(x)dx.
\]
This is a function on $M({\mathbb {A}}_S)$.

Suppose now $\gamma \in M({\mathbb {A}}_S)$ is arbitrary and $a\in
A_M({\mathbb {A}}_S)$ is chosen so that
$M({\mathbb {A}}_S)_{a\gamma }=G({\mathbb {A}}_S)_{a\gamma }$.
Following \cite{A7}, (8.1)
we define $J_M(\gamma ,f)$ to be the limit at $a=1$
of $|D(a\gamma )|^{1/2}$ times the integral
\begin{equation}
\int_{G({\mathbb {A}}_S)_\sigma \backslash G({\mathbb {A}}_S)}
\int_{M({\mathbb {A}}_S)_\gamma \backslash G({\mathbb {A}}_S)_\sigma }
f(y^{-1}\sigma x^{-1}a u xy) \left( \sum_{L\in {\cal L}(M)}r_M^L(\gamma
,a)v_L(xy)\right) dxdy,
\end{equation}
where $\gamma\in M({\mathbb A}_S)$ is as above and $r_M^L$ is as in \cite{A7}.

\begin{example}
Let $G=SL(2),\;M=A,\;a=\left(
\begin{array}{cc}
t & 0 \\
0 & t^{-1}
\end{array}
\right) $, where $t\neq \pm 1$. In this case,
\[
r_A(\gamma ,a)=1,
\]
and
\[
v_A(x)=\sum_{v\in S}v_{A,v}(x_v),
\]
where $x=\prod_{v\in S}a_v\cdot \left[
\begin{array}{cc}
1 & n_v \\
0 & 1
\end{array}
\right] \cdot k_v,$ for
$a_v\in {A(k_v)}$, $k_v\in {K_v}$
where $v_{A,v}$ is the local weight of \S\ref{sec:weighted-defs}
(\cite{A6}, Intro., \cite{Sho}, chap. 6). Therefore,
$J_A(z,f)$, for $z\in Z(G(k))$,
is the limit as $a\rightarrow 1$ of
\begin{equation}
\ |D(a)|^{1/2}
\int_{A({\mathbb {A}}_S)\backslash G({\mathbb {A}}_S)}
f(y^{-1} az y)v_A(y)dy
\end{equation}
\cite{Sho}, section 6.1, p. 64.
\end{example}

\subsection{Definitions}

Let ${\cal X}$ denote the equivalence classes of pairs $(M,\rho)$ as in
\cite{A2}, \cite{A3}.
For each standard parabolic $P=MN$ of $G$ defined over $k$
and each $\chi \in {\cal X}$, we define
$L^2(M(k)\backslash M({\mathbb{A}})^1)_\chi $
as follows. If there is no $P_1\in {\cal P}_\chi $
with $P_1\subset P$ then we define
$L^2(M(k)\backslash M({\mathbb{A}})^1)_\chi $ to be $0$.
Otherwise, let
$\Psi (P_1,P,\chi,W )$ denote the space of
all complex-valued functions $\psi :N({\mathbb{A}})M_1(k)\backslash G(%
{\mathbb{A}})\rightarrow {\mathbb C}$, $P_1=N_1M_1\in {\cal P}_\chi $
with $P_1\subset P$, such that

\begin{enumerate}
\item  $\psi (amk)$, as a function of
$a\in A_1(k)\backslash A_1({\mathbb{A}})$ (fixed
$m\in M_1(k)\backslash M_1({\mathbb{A}})^1,\;k\in K$),
is compactly supported,

\item  $\psi (amk)$, as a function of $k\in K_\infty $, belongs to the
vector space defined by the span of the matrix coefficients of some
irreducible representation $W$ of $K_\infty $,

\item  $\psi (amk)$, as a function of $m\in M$, belongs to
$L^2(M(k)\backslash M({\mathbb{A}})^1)_{cusp,\chi }$.
\end{enumerate}

Given $\psi \in \Psi (P_1,P,\chi,W )$, $P_1\subset P$, let
\[
\psi ^{\wedge ,M}(m)=
\sum_{\delta \in P_1(k)\cap M(k)\backslash
M(k)}\psi (\delta m),\;\;\;\;\;m\in M(k)\backslash
M({\mathbb{A}})^1,
\]
and by \cite[Lemma 1]{L1} the function $\psi ^{\wedge ,M}$ lies in
$L^2(M(k)\backslash M({\mathbb{A}})^1)$. We call this a {\bf theta function}.
Let $L^2(M(k)\backslash M({\mathbb{A}})^1)_\chi $ denote
the $L^2$-closure of the span of all the $\psi ^{\wedge ,M}$ with $\psi \in
\Psi (P_1,P,\chi,W )$, some $W$ and $P_1\subset P$ as above.
Lemma 2 of \cite{L1} implies
\begin{equation}
L^2(M(k)\backslash M({\mathbb{A}})^1)=\bigoplus_{\chi \in
{\cal X}}L^2(M(k)\backslash M({\mathbb{A}})^1)_\chi .
\end{equation}

Now, define $I_P(\pi)$ induced representation as in \cite{A2},
\cite{A3}. If $f$ is a smooth function and $\chi\in {\cal X}$,
Define $I_P(\pi,f)_\chi$ denote the
$L^2(M(k)\backslash M({\mathbb{A}})^1)_\chi$-isotypic component
of $I_P(\pi,f)$.

\begin{definition}
Let $G=SL(2)$. Define $M^T_P$ as follows:

if $\chi =(G,\rho )$, let $M_G^T(\pi )_\chi =1$;

if $\chi =(G,\rho )$, let $M_B^T(\pi )_\chi =1$;

if $P=B,\chi =(A,\rho )$ unramified, let
\[
\begin{array}{c}
M_B^T(\pi )_\chi \\
=vol({\mathbb{R}}/{\mathbb{Z}})
\lim_{\varsigma \rightarrow 0}
\left[
\frac{e^{\varsigma (T)}
M(1,\pi )^{-1}M(1,\pi ,\varsigma )}{\varsigma (\alpha^{\vee })}
+\frac{e^{-\varsigma (T)}
M(s,\pi )^{-1}M(s,\pi ,\varsigma )}{-\varsigma (\alpha^{\vee })}
\right]\\
=M(s,\pi)^{-1}
\lim_{\varsigma \rightarrow 0}
{\frac{M(s,\pi)-e^{-2\varsigma (T)}M(s,\pi,\varsigma)}{\varsigma(\alpha^\vee)}} ,
\end{array}
\]
where $\alpha $ denotes the unique positive root.
\end{definition}

\begin{lemma}
For $G=SL(2)$, each $(G,\pi)\in\chi\in{\cal{X}}$ is
unramified. Moreover, $(A,\mu)\in\chi\in{\cal{X}}$ is unramified
if and only if $\mu^w\not=\mu$.
In the case case of the above definition, at a function $\phi$ on $K$,
\[
\begin{array}{c}
(M^T_B(\pi)_\chi \phi)(k)=
M(s,\pi)^{-1}
\lim_{\varsigma \rightarrow 0}
{\frac{M(s,\pi)\phi(k)-e^{-2\varsigma (T)}e^{\varsigma (H(k))}
M(s,\pi_\varsigma)\phi_\varsigma(k)}{\varsigma(\alpha^\vee)}} \\
=M(s,\pi)^{-1}
\lim_{\varsigma \rightarrow 0}
{\frac{M(s,\pi)\phi(k)-e^{-2\varsigma (T)}M(s,\pi_\varsigma)\phi_\varsigma (k)}{\varsigma(\alpha^\vee)}} \\
=M(s,\pi)^{-1}M'(s,\pi)\phi(k)-2 \phi(k) T,
\end{array}
\]
where the prime denotes differentiation with respect to
$\varsigma$. Thus
\[
M_B^T(\pi )_\chi =
M(s,\pi)^{-1}M'(s,\pi)-2\cdot I\cdot T,
\]
where $I$ denotes the identity.
\end{lemma}

\subsection{The trace formula}

The geometric side of the trace formula is the sum of the following terms:

(a) the sum over classes ${\cal{O}}\subset G(k)$
such that ${\cal{O}}\cap Z(G(k))\not= \emptyset$ of
\begin{equation}
J_{\cal{O}}(f)=
\sum_{\gamma\in {\cal O}}
a^G(S,\gamma)J_G(\gamma,f)
+{\frac{1}{2}}\sum_{\gamma\in {\cal O}\cap A({\mathbb A})}
a^A(S,\gamma)J_A(\gamma,f),
\end{equation}
where $J_M$ is defined above and
the constants $a^M(S,\gamma )$ satisfy, if $\gamma$ is semisimple,
the property
\[
a^M(S,\gamma )=\left\{
\begin{array}{c}
vol(M(k)_\gamma \backslash M({\mathbb {A}}_S)_\gamma ^1),
\;\;\;\;{\rm if\;}\gamma \;{\rm is}\;
k-{\rm elliptic\;in}\;M(k), \\
0,\;\;\;\;\;\;\;\;\;\;\;\;\;\;\;\;\;\;{\rm otherwise},\;\;\;\;\;
\end{array}
\right.
\]

(b) the sum over elliptic regular classes ${\cal{O}}\subset G(k)$
of
\begin{equation}
J_{\cal{O}}(f)=
\sum_{\gamma\in {\cal O}}
a^G(S,\gamma)J_G(\gamma,f),
\end{equation}
where $J_G$ (an orbital integral) and $a^G$ (a volume)
are as above,

(c) the sum over ``hyperbolic'' regular classes
${\cal{O}}\subset G(k)$
of
\begin{equation}
J_{\cal{O}}(f)=
{\frac{1}{2}}\sum_{\gamma\in {\cal O}\cap A({\mathbb A})}
a^A(S,\gamma)J_A(\gamma,f),
\end{equation}
where $J_A$ and $a^A$ are as above,

The spectral side of the trace formula is the sum of the following terms:

(d) the sum over classes $\chi\in {\cal{X}}(G)$ of
\begin{equation}
J_\chi(f)=
\tr(\pi(f)),\ \ \ (\pi,G)\in\chi,
\end{equation}

(e) the sum over classes $\chi\in {\cal{X}}-{\cal{X}}(G)$ of
\begin{equation}
J_\chi(f)=
{\frac{1}{2}}\int_{\Pi_u(A({\mathbb{A}}_S))}
\tr(M_A^T(\pi)_\chi
I_B(\pi,f)_\chi)d\pi,
\end{equation}
where $M_A^T(\pi)_\chi$ is (essentially a logarithmic derivative
of an intertwining operator) given by the definition and lemma above.

\section{The trace formula for $\overline{SL(2)}$}

The analogs of the definitions above for the 2-fold cover of $SL(2)$
are so similar that we do not repeat them (precise definitions
are all given in \cite{J2}).

The geometric side of the trace formula is the sum of the following terms:

(a) the sum over classes ${\cal{O}}\subset G(k)$
such that ${\cal{O}}\cap Z(G(k))\not= \emptyset$ of
\begin{equation}
J_{\cal{O}}(f)=
\sum_{\gamma\in {\cal O}}
a^G(S,\gamma)J_G(\gamma,f)+{\frac{1}{2}}
\sum_{\gamma\in {\cal O}\cap A({\mathbb A})}
a^A(S,\gamma)J_A(\gamma,f),
\end{equation}
where $J_M$ and $a^M$
are given by the ``obvious'' analogs of the formulas
in the previous section,

(b) the sum over elliptic regular classes ${\cal{O}}\subset G(k)$
of
\begin{equation}
J_{\cal{O}}(f)=
\sum_{\gamma\in {\cal O}}
a^G(S,\gamma)J_G(\gamma,f),
\end{equation}
where $J_G$ (an orbital integral) and $a^G$ (a volume)
are given by the ``obvious'' analogs of the formulas
in the previous section,

(c) the sum over ``hyperbolic'' regular classes
${\cal{O}}\subset G(k)$
of
\begin{equation}
J_{\cal{O}}(f)=
{\frac{1}{2}}\sum_{\gamma\in {\cal O}\cap A({\mathbb A})}
a^A(S,\gamma)J_A(\gamma,f),
\end{equation}
where $J_A$ and $a^A$
are given by the ``obvious'' analogs of the formulas
in the previous section,

The spectral side of the trace formula is the sum of the following terms:

(d) the sum over the cuspidal classes $\chi\in {\cal{X}}(G)$ of
\begin{equation}
J_\chi(f)=
\tr(\pi(f)),\ \ \ (\pi,G)\in\chi,
\end{equation}

(e) the sum over classes $\chi\in {\cal{X}}-{\cal{X}}(G)$ of
\begin{equation}
J_\chi(f)=
{\frac{1}{2}}\int_{\Pi_u(\overline{A({\mathbb{A}}_S)})}
\tr(M_A(\pi)_\chi I_B(\pi,f)_\chi)d\pi,
\end{equation}
where $M_A(\pi)_\chi =M_A^T(\pi)_\chi|_{T=0}$.
%% ===== end aabook5.tex =====


%% ===== begin aabook6.tex (rewritten in Draft 1) =====
\chapter{Global comparison and multiplicity one}\label{ch:global}

Throughout this chapter, $k$ is a number field, ${\mathbb{A}}={\mathbb{A}}_k$,
$G=SL(2)$, and $\overline{G({\mathbb{A}})}$ is the 2-fold cover of
\S\ref{sec:global}. Let ${\cal H}(G({\mathbb{A}}))$ and
${\cal H}(\overline{G({\mathbb{A}})})$ denote the spaces of global
functions of \S\ref{sec:globalalg}. We write $R_{cusp}$ for the
representation of $G({\mathbb{A}})$ on
$L^2_{cusp}(G(k)\backslash G({\mathbb{A}}))$ and $\overline{R}_{cusp}$ for
the representation of $\overline{G({\mathbb{A}})}$ on the genuine part of
$L^2_{cusp}(G(k)\backslash \overline{G({\mathbb{A}})})$, where $G(k)$ is
embedded in $\overline{G({\mathbb{A}})}$ by the splitting of
\S\ref{sec:global}. For a cuspidal automorphic representation $\pi$ of
$G({\mathbb{A}})$ (resp. genuine cuspidal automorphic representation
$\overline{\pi}$ of $\overline{G({\mathbb{A}})}$), $m(\pi)$ (resp.
$m(\overline{\pi})$) denotes its multiplicity in $R_{cusp}$ (resp.
$\overline{R}_{cusp}$). As usual, global test functions are assumed
$K_\infty$-finite.

For each place $v$ we write $K_{0,v}=G({\cal O}_v)$ if $v$ is finite,
and, if $v$ is finite and odd, $K_{1,v}=w^{-1}K_{0,v}w$ with $w$ as in
\S\ref{sec:filtrations}. These are representatives of the two conjugacy
classes of hyperspecial maximal compact subgroups of $G(k_v)$. Since $v$
is odd, $\overline{K_{i,v}}$ splits (Lemma~\ref{lem:facts}(b)), and
${\cal H}(\overline{G(k_v)},K_{i,v})$ denotes the Hecke algebra of
genuine functions that are bi-invariant under the image of the
splitting.

The argument is modeled on that of Gelbart and Jacquet \cite{GJ}, who
compare automorphic representations of $GL(2)$ and of the multiplicative
group of a quaternion algebra. Besides the groups being different, there
are two other differences: $SL(2)$ and its 2-fold cover do not satisfy
strong multiplicity one, and the local comparison of functions is more
complicated. The input that replaces strong multiplicity one is
Waldspurger's multiplicity one theorem for the cover, applied through
Lemma~\ref{lem:A00}.

\section{The hypotheses}\label{sec:hypotheses}

\begin{definition}\label{def:exceptional}
For a cuspidal automorphic representation $\pi=\otimes_v\pi_v$ of
$G({\mathbb{A}})$, let $S(\pi)$ denote the finite set of places $v$ of $k$
such that $v$ is archimedean, or $v$ is dyadic, or $v$ is odd and $\pi_v$
has no nonzero $K_{0,v}$-fixed vector.
\end{definition}

\begin{definition}\label{def:corresponding}
Let $S$ be a finite set of places of $k$ containing the archimedean and
dyadic places. For each $v\in S$, let ${\cal C}_v$ be a vector space of
pairs $(f_v,\overline{f}_v)$ with $f_v\in C_c^\infty(G(k_v))$ and
$\overline{f}_v\in C_c^\infty(\overline{G(k_v)})_{gen}$ ($K_v$-finite if
$v$ is archimedean). A pair $(f,\overline{f})$ with
$f=\otimes_v f_v\in{\cal H}(G({\mathbb{A}}))$ and
$\overline{f}=\otimes_v\overline{f}_v\in{\cal H}(\overline{G({\mathbb{A}})})$
is called {\bf $(S,{\cal C})$-corresponding} if
\begin{enumerate}
\item $(f_v,\overline{f}_v)\in{\cal C}_v$ for each $v\in S$;
\item for each $v\notin S$, $f_v\in{\cal H}(G(k_v),K_{0,v})$ and
$\overline{f}_v\in{\cal H}(\overline{G(k_v)},K_{0,v})_{gen}$ correspond
under the Satake isomorphism of \S\ref{sec:spherical};
\item for almost all $v$, $f_v$ and $\overline{f}_v$ are the unit
elements of the respective spherical Hecke algebras.
\end{enumerate}
\end{definition}

\begin{hypothesis}[local spectral transfer]\label{hyp:LH}
Let $v$ be a place of $k$ and $\pi_v$ an irreducible unitary
representation of $G(k_v)$. We say that {\bf Hypothesis LH holds for
$\pi_v$ with the class ${\cal C}_v$} if ${\cal C}_v$ is a vector space of
pairs as in Definition~\ref{def:corresponding} with the following
properties.
\begin{enumerate}
\item[(LH0)] ${\cal C}_v=\bigcup_L{\cal C}_v(L)$ is a directed union of
subspaces ${\cal C}_v(L)$, indexed by open compact subgroups $L$ of
$G(k_v)$ if $v$ is finite (with the pairs in ${\cal C}_v(L)$ bi-invariant
under $L$, respectively under its image in a splitting of $\overline{L}$
over a subgroup of finite index), and by finite sets $L$ of
$K_v$-types if $v$ is archimedean (with the pairs in ${\cal C}_v(L)$
having $K_v$-types in $L$, respectively in the corresponding genuine
$\overline{K_v}$-types).
\item[(LH1)] (Separation) For each $L$, and for each finite family
$\sigma_1,\dots,\sigma_r$ of pairwise inequivalent irreducible
admissible representations of $G(k_v)$ none of whose characters vanishes
identically on ${\cal C}_v(L)$, the linear functionals
$(f_v,\overline{f}_v)\longmapsto\tr \sigma_j(f_v)$ on
${\cal C}_v(L)$ are linearly independent.
\item[(LH2)] (Transfer) For each genuine irreducible unitary
representation $\overline{\tau}$ of $\overline{G(k_v)}$, either
$\tr \overline{\tau}(\overline{f}_v)=0$ for all
$(f_v,\overline{f}_v)\in{\cal C}_v$, or there is an irreducible
admissible representation $t_v(\overline{\tau})$ of $G(k_v)$ such that
\[
\tr \overline{\tau}(\overline{f}_v)=\tr t_v(\overline{\tau})(f_v)
\qquad\text{for all }(f_v,\overline{f}_v)\in{\cal C}_v .
\]
Moreover, $t_v$ is injective on the set of $\overline{\tau}$ for which it
is defined.
\item[(LH3)] (Detection) $\tr \pi_v(f_v)\neq 0$ for some
$(f_v,\overline{f}_v)\in{\cal C}_v$.
\item[(LH4)] (Infinitesimal characters; $v$ archimedean only) For each
infinitesimal character $\nu_v$ of $G(k_v)$, the genuine $\overline\tau$
for which $t_v(\overline\tau)$ is defined and has infinitesimal character
$\nu_v$ have only finitely many distinct infinitesimal characters.
\end{enumerate}
\end{hypothesis}

For example, for the elliptic matchings of \S\ref{sec:realds}, a genuine
discrete series representation of infinitesimal character $\overline\nu$
is matched with a representation of infinitesimal character
$2\overline\nu$ or $2\overline\nu\pm1$, so (LH4) holds for them.

\begin{definition}\label{def:infchar}
Let $\mathfrak{z}_\infty$ denote the center of the universal enveloping
algebra of the complexified Lie algebra of
$G_\infty=\prod_{v|\infty}G(k_v)$. Since $\overline{G_\infty}$ has the same
Lie algebra, $\mathfrak{z}_\infty$ acts by a character (the infinitesimal
character) on every irreducible admissible representation of $G_\infty$
or of $\overline{G_\infty}$. For a cuspidal automorphic representation
$\pi'$ of $G({\mathbb{A}})$, let $\nu_\infty(\pi')$ be the infinitesimal
character of $\pi'_\infty=\otimes_{v|\infty}\pi'_v$.

Suppose that, for each archimedean $v$, ${\cal C}_v$ satisfies (LH2), with
transfer $t_v$. A genuine cuspidal automorphic representation
$\overline\pi'$ of $\overline{G({\mathbb{A}})}$ {\bf lies over} the
character $\nu$ of $\mathfrak{z}_\infty$ if $t_v(\overline\pi'_v)$ is
defined for each archimedean $v$ and $\otimes_{v|\infty}t_v(\overline\pi'_v)$
has infinitesimal character $\nu$. For global functions $f$ and genuine
global functions $\overline f$, put
\[
\tr R_{cusp}^\nu(f)=\sum_{\pi':\ \nu_\infty(\pi')=\nu}m(\pi')\tr\pi'(f),
\qquad
\tr \overline{R}{}_{cusp}^{\,\nu}(\overline f)=
\sum_{\overline\pi'\text{ lying over }\nu}m(\overline\pi')\tr\overline\pi'(\overline f).
\]
These are subseries of the absolutely convergent series for
$\tr R_{cusp}(f)$ and $\tr\overline R_{cusp}(\overline f)$.
\end{definition}

\begin{example}
Let $k={\mathbb Q}$, and let $\pi'$ be a cuspidal automorphic
representation of $SL(2,{\mathbb A})$ occurring in the restriction of the
representation of $GL(2,{\mathbb A})$ generated by the cusp form $\Delta$
of weight $12$. Then $\pi'_\infty$ is one of the discrete series
$\pi_{\pm11}$ (lowest $K$-types $\pm12$), and $\nu_\infty(\pi')$ is the
infinitesimal character $11$. Take ${\cal C}_\infty$ to be one of the
elliptic matchings of \S\ref{sec:realds}. By Lemma~\ref{lem:realds}, a
genuine $\overline\pi'$ lies over $\nu=11$ only if
$\overline\pi'_\infty=\overline\pi{}^{\pm}_{k/2}$ with $k-2=11$ (even
matching) or with $k-1=11$ or $k-3=11$ (odd matchings). Since $k$ must
be odd, only the even matching contributes, and then
$\overline\pi'_\infty=\overline\pi{}^{\pm}_{13/2}$, of infinitesimal
character $11/2$: the finite set allowed by (LH4) is $\{11/2\}$.
\end{example}

\begin{hypothesis}[global trace identity]\label{hyp:GH}
For the given $S$ and $({\cal C}_v)_{v\in S}$, and for every character
$\nu$ of $\mathfrak{z}_\infty$,
\[
\tr R_{cusp}^\nu(f)=\tr \overline{R}{}_{cusp}^{\,\nu}(\overline{f})
\]
for all $(S,{\cal C})$-corresponding pairs $(f,\overline{f})$.
\end{hypothesis}

If $(f,\overline f)$ is $(S,{\cal C})$-corresponding and $\overline\pi'$
lies over no $\nu$, then some archimedean $\tr\overline\pi'_v$ vanishes on
${\cal C}_v$ by (LH2), so $\tr\overline\pi'(\overline f)=0$. Summing
Hypothesis~\ref{hyp:GH} over $\nu$ therefore gives
$\tr R_{cusp}(f)=\tr\overline R_{cusp}(\overline f)$, which is the form
of the hypothesis used in Draft~1.

\begin{remark}\label{rem:why-refine}
(Added in Draft~2.) (1) Draft~1 assumed only the unrefined identity
$\tr R_{cusp}(f)=\tr\overline R_{cusp}(\overline f)$, and reduced it to a
finite identity by a ``finiteness lemma'' asserting that only finitely
many cuspidal representations have a vector fixed by a given open compact
subgroup $L_f$ with $K_\infty$-types in a given finite set. That is
false: for $k={\mathbb Q}$, $L_f=SL(2,\widehat{\mathbb Z})$, and the
trivial $SO(2)$-type, there are infinitely many spherical Maass cusp forms
of level one, with Laplace eigenvalues tending to infinity (Weyl's law
for the cuspidal spectrum). What is true is the finiteness for a
\emph{fixed infinitesimal character} (Lemma~\ref{lem:finiteness}), and
Hypothesis~\ref{hyp:GH} is refined accordingly.

(2) One would hope to derive the refined identity from the unrefined one
by spectral multipliers at the archimedean places: for a $W$-invariant
Paley--Wiener function $\alpha$ there is an operation
$f_\infty\mapsto f_{\infty,\alpha}$ on $K_\infty$-finite functions in
$C_c^\infty(G_\infty)$ with $\tr\pi(f_{\infty,\alpha})=\alpha(\nu_\pi)\tr\pi(f_\infty)$
(Arthur's multipliers, \cite{APW}). If each archimedean class ${\cal C}_v$
were stable under such multipliers, applied compatibly to $f_v$ and
$\overline f_v$ (compatibility being governed by the relation between
infinitesimal characters in (LH4)), then the unrefined identity would hold
with $f$ replaced by $f_\alpha$ for all $\alpha$, and one could hope to
isolate a single $\nu$ using the absolute convergence of the spectral
expansions. Neither the stability of the classes nor the isolation
argument is carried out here, so we assume the refined form.
\end{remark}

\begin{remark}\label{rem:GH}
(1) Hypothesis~\ref{hyp:GH} is what a full comparison of the trace
formulas of Chapter~\ref{ch:tf} would give, refined as in
Remark~\ref{rem:why-refine}(2). Write the trace formula of
$\overline{G({\mathbb{A}})}$ as
\[
\tr \;\overline{R}_{cusp}(\overline{f})+
\sum_{\chi \in \overline{{\cal X}}-\overline{{\cal X}}(G)}
J_\chi (\overline{f}) =\sum_{{\cal O}\in {\cal C}}J_{{\cal O}}(\overline{f}),
\]
and the trace formula of $G({\mathbb{A}})$ as
\[
\tr \;R_{cusp}(f)+
\sum_{\chi \in {\cal X}-{\cal X}(G)}J_\chi(f)
=\sum_{{\cal O}\in {\cal C}}J_{{\cal O}}(f).
\]
If the geometric sides agree term by term for corresponding pairs, then
subtracting gives
\begin{equation}\label{eq:decomp-measures}
\tr \;\overline{R}_{cusp}(\overline{f})-\tr \;R_{cusp}(f)
=\sum_{\chi \in {\cal X}-{\cal X}(G)}J_\chi (f)
-\sum_{\chi \in \overline{{\cal X}}-\overline{{\cal X}}(G)}
J_\chi(\overline{f}).
\end{equation}
(The 1995 draft had a plus sign in front of the second sum.)
The unrefined identity then amounts to the vanishing of both sides
of (\ref{eq:decomp-measures}). Langlands' decomposition-of-measures
argument \cite{L4}, as adapted by Flicker \cite[\S\S 4.1--4.2]{Fl} (see
also \cite[\S\S 10.5, 14.5]{R}), is designed to prove such a vanishing.
The 1995 draft asserted that Flicker's argument applies almost verbatim
to the present situation; this has not been checked.

(2) A term-by-term comparison of the geometric sides requires matching
of orbital integrals and of weighted orbital integrals (or, in the
invariant trace formula, of the invariant distributions $I_M$) at every
place. Chapter~\ref{ch:local} provides this only in part. For spherical
functions outside $S$, the Satake isomorphism gives the matching of
hyperbolic orbital integrals but not of elliptic ones
(Remark~\ref{rem:satake-oi}); weighted orbital integrals of spherical
functions are matched only for the unit elements
(Remark~\ref{rem:weighted-units}), while Definition~\ref{def:corresponding}
lets $f_v$ range over the whole spherical Hecke algebra for $v\notin S$;
weighted orbital integrals cannot be matched for arbitrary corresponding
pairs (Remark~\ref{rem:weighted-noninv}); and nothing is matched at the
dyadic places. In particular, Hypothesis~\ref{hyp:GH} contains much of
the global comparison that the 1995 project hoped to establish.
\end{remark}

\section{The main theorem}\label{sec:main}

\begin{theorem}\label{thm:main}
Let $\pi=\otimes_v\pi_v$ be a cuspidal automorphic representation of
$G({\mathbb{A}})$ and let $S=S(\pi)$. Suppose that for each $v\in S$
Hypothesis~\ref{hyp:LH} holds for $\pi_v$ with a class ${\cal C}_v$, and
that Hypothesis~\ref{hyp:GH} holds for $S$ and $({\cal C}_v)_{v\in S}$.
Then:
\begin{enumerate}
\item[(a)] $m(\pi)=1$;
\item[(b)] there is a unique genuine cuspidal automorphic
representation $\overline{\pi}=\otimes_v\overline{\pi}_v$ of
$\overline{G({\mathbb{A}})}$ such that $\overline{\pi}_v$ corresponds to
$\pi_v$ under the Satake correspondence of \S\ref{sec:spherical} for
$v\notin S$, and $t_v(\overline{\pi}_v)\cong\pi_v$ for $v\in S$.
\end{enumerate}
\end{theorem}

We shall use the following facts.

\begin{lemma}[finiteness]\label{lem:finiteness}
Let $L_f$ be an open compact subgroup of $G({\mathbb{A}}_f)$, $E$ a
finite set of $K_\infty$-types, and $\nu$ a character of
$\mathfrak z_\infty$. Then only finitely many cuspidal
automorphic representations $\pi'$ of $G({\mathbb{A}})$, counted with
multiplicity, have $\nu_\infty(\pi')=\nu$ and a nonzero vector which is
$L_f$-fixed and whose $K_\infty$-types lie in $E$. The same holds for
genuine cuspidal automorphic representations of
$\overline{G({\mathbb{A}})}$, with $L_f$ replaced by an open compact
subgroup of $\overline{G({\mathbb{A}}_f)}$ meeting $\mu_2$ trivially and
$E$ by a finite set of genuine $\overline{K_\infty}$-types.
\end{lemma}

\begin{proof}
Distinct copies of such representations in the cuspidal spectrum
contribute linearly independent vectors to the space of cusp forms which
are $L_f$-invariant, have $K_\infty$-types in $E$, and are annihilated by
the kernel of $\nu$ in $\mathfrak z_\infty$. That space is
finite-dimensional by Harish-Chandra's finiteness theorem \cite{HC68}
(see also \cite{BJ}, \cite{MW}). For the cover, one reduces as usual to
finitely many quotients $\Gamma\backslash\overline{G_\infty}$ with
$\Gamma$ arithmetic, and Harish-Chandra's argument applies; we have not
written out the details. Without the condition on $\nu$ the statement
is false (Remark~\ref{rem:why-refine}(1)).
\end{proof}

\begin{lemma}[tensor products]\label{lem:tensor}
For $v$ in a finite set $S$, let $V_v$ be a complex vector space and
$\Lambda_v$ a finite linearly independent set of linear functionals on
$V_v$. Then the functionals $\otimes_{v\in S}\lambda_v$
($\lambda_v\in\Lambda_v$) on $\otimes_{v\in S}V_v$ are linearly
independent.
\end{lemma}

\begin{proof} Choose $x_{\mu}\in V_v$ ($\mu\in\Lambda_v$) with
$\lambda(x_\mu)=\delta_{\lambda\mu}$ for $\lambda,\mu\in\Lambda_v$, and
evaluate a linear relation at $\otimes_v x_{\mu_v}$.
\end{proof}

Let ${\cal A}_0$ denote the space of genuine cusp forms on
$\overline{G({\mathbb{A}})}$ (so that $\overline R_{cusp}$ acts on its
completion), let $\Theta_0\subset{\cal A}_0$ be the closed span of the
cuspidal elementary theta series, i.e.\ the theta series attached to the
Weil representations $r_{\psi,\chi}$ of one-dimensional quadratic forms
(\cite{GePS}, \cite{Ge2}), and let ${\cal A}_{00}$ be the orthogonal
complement of $\Theta_0$ in ${\cal A}_0$. Waldspurger's multiplicity one
theorem \cite{Wa1} (see also \cite{PS}) states that every irreducible
genuine representation occurs in ${\cal A}_{00}$ with multiplicity at
most one. It is a statement about ${\cal A}_{00}$, not ${\cal A}_0$; Draft~1
quoted it for all of ${\cal A}_0$.

\begin{lemma}\label{lem:A00}
Let $\pi$ be a cuspidal automorphic representation of $G({\mathbb{A}})$,
let $S\supset S(\pi)$ be finite, and let $\overline{\pi}'$ be a genuine
cuspidal automorphic representation of $\overline{G({\mathbb{A}})}$ such
that $\overline{\pi}'_v$ corresponds to $\pi_v$ under the Satake
correspondence for every $v\notin S$. Then no irreducible constituent of
$\Theta_0$ is isomorphic to $\overline\pi'$. Consequently every copy of
$\overline\pi'$ in the genuine cuspidal spectrum lies in ${\cal A}_{00}$,
and $m(\overline{\pi}')\leq 1$.
\end{lemma}

\begin{proof}
(Added in Draft~2.) We use three facts.

(i) Every irreducible constituent $\sigma$ of $\Theta_0$ is a quotient of
a global Weil representation $\otimes_v r_{\psi_v,\chi_v}$, and so
$\sigma_v$ is the even part $r^+_{\psi_v,\chi_v}$ of the local Weil
representation for almost all $v$ \cite{GePS}, \cite{Ge2}.

(ii) Let $v$ be odd, with $\psi_v$ of conductor ${\cal O}_v$ and $\chi_v$
unramified. By formula (iii) of \S\ref{sec:weil}, the functional
$\phi\mapsto\phi(0)$ identifies the Jacquet module of $r^+_{\psi_v,\chi_v}$
with the genuine character $(h(a),\zeta)\mapsto\zeta|a|^{1/2}\omega(a)$ of
$\overline{A(k_v)}$. After normalizing by $\delta_B^{-1/2}(h(a))=|a|^{-1}$,
Frobenius reciprocity embeds $r^+_{\psi_v,\chi_v}$ (which is irreducible
and has a vector fixed by $\kappa_v(K_{0,v})$, namely the characteristic
function of ${\cal O}_v$) in $I_A(\overline\chi)$ with
$|\overline\chi(a(x),1)|=|x|^{-1/2}$. Under the Satake correspondence of
Definition~\ref{def:satake-corr}, which does not change $\chi_1$ on
$A_2(k_v)$, it therefore corresponds to the spherical constituent of
$I_A(\chi)$ for a character $\chi$ with $|\chi(a(x))|=|x|^{\pm1/2}$.

(iii) By Labesse--Langlands \cite{LL}, $\pi$ is a constituent of the
restriction to $SL(2,{\mathbb{A}})$ of a cuspidal automorphic
representation $\Pi$ of $GL(2,{\mathbb{A}})$. At each place $v$ where $\Pi_v$
is unramified, $\Pi_v\cong I(\mu_1,\mu_2)$ with
$|\mu_i(x)|=|x|^{\pm s_v}$ and $|s_v|<1/4$ \cite{GJ78}. Hence $\pi_v$ is
the spherical constituent of $I_A(\chi)$ with $\chi(a(x))=\mu_1\mu_2^{-1}(x)$,
and $|\chi(a(x))|=|x|^{\pm 2s_v}$ with $|2s_v|<1/2$.

Now let $\sigma$ be an irreducible constituent of $\Theta_0$. Choose an odd
place $v\notin S$ at which $\Pi_v$, $\psi_v$, $\chi_v$ are unramified and
$\sigma_v\cong r^+_{\psi_v,\chi_v}$; all but finitely many places qualify.
If $\sigma\cong\overline\pi'$, then $r^+_{\psi_v,\chi_v}$ would correspond
to $\pi_v$ under the Satake correspondence, which by (ii) and (iii)
is impossible, because a spherical representation determines the
character $\chi$ of the principal series containing it up to
$\chi\mapsto\chi^{-1}$. Hence $\overline\pi'$ is isomorphic to no
constituent of $\Theta_0$. Since $\Theta_0$ is a closed invariant
subspace of the (discretely decomposable) genuine cuspidal spectrum, the
orthogonal projection to $\Theta_0$ of any copy of $\overline\pi'$ is
zero, so every copy lies in ${\cal A}_{00}$. Waldspurger's theorem gives
$m(\overline\pi')\leq1$.
\end{proof}

\begin{remark}
The proof of Lemma~\ref{lem:A00} uses the bound $|s_v|<1/4$ only through
the strict inequality $|2s_v|<1/2$; any bound toward the Ramanujan
conjecture for $GL(2)$ that is strictly better than the trivial one
$|s_v|<1/2$, and that holds at infinitely many places, would do. The
normalizations in (ii) (of the Weil representation and of the
correspondence of Definition~\ref{def:satake-corr}) should be rechecked
against \cite{Ge2} and \cite{GePS}.
\end{remark}

\medskip

\begin{proof}[Proof of Theorem~\ref{thm:main}]
Let $\nu=\nu_\infty(\pi)$.
For each $v\in S$ use (LH3) and (LH0) to fix an index $L_v$ such that
$\tr \pi_v$ does not vanish identically on ${\cal C}_v(L_v)$.
Restrict attention to $(S,{\cal C})$-corresponding pairs with
$(f_v,\overline{f}_v)\in{\cal C}_v(L_v)$ for $v\in S$. Let ${\cal H}^S$ be
the restricted tensor product of the spherical Hecke algebras
${\cal H}(G(k_v),K_{0,v})$ for $v\notin S$; for $h\in{\cal H}^S$, let
$\overline{h}$ denote the corresponding genuine function, and write
$f=f_S\otimes h$, $\overline{f}=\overline{f}_S\otimes\overline{h}$.

{\em Step 1 (finiteness).} If $\tr \pi'(f)\neq 0$ for some such $f$,
then $\pi'_v$ has a nonzero $K_{0,v}$-fixed vector for $v\notin S$, and
$\pi'_S=\otimes_{v\in S}\pi'_v$ has a nonzero vector fixed by
$\prod_{v\in S, v<\infty}L_v$ with $K_\infty$-types in $\prod_{v|\infty}L_v$.
By Lemma~\ref{lem:finiteness}, only a finite set ${\cal A}$ of cuspidal
automorphic representations $\pi'$ with $\nu_\infty(\pi')=\nu$ can
contribute to $\tr R^\nu_{cusp}(f)$. Similarly, a genuine $\overline\pi'$
contributing to $\tr\overline R{}^{\,\nu}_{cusp}(\overline f)$ has a nonzero
vector fixed by the image of a splitting of an open compact subgroup of
$\overline{G({\mathbb A}_f)}$ determined by the $L_v$ ($v\in S$ finite) and
the $\kappa_v(K_{0,v})$ ($v\notin S$), with genuine $\overline{K_\infty}$-types
in a finite set determined by the $L_v$ ($v|\infty$); and since it lies
over $\nu$, (LH4) allows only finitely many infinitesimal characters for
$\overline\pi'_\infty$. By Lemma~\ref{lem:finiteness} again, only a finite
set $\overline{{\cal A}}$ of such $\overline\pi'$ can contribute.
Hypothesis~\ref{hyp:GH} for this $\nu$
therefore reads
\begin{equation}\label{eq:finite-identity}
\sum_{\pi'\in{\cal A}}m(\pi')\,\omega_{\pi'}^S(h)
\prod_{v\in S}\tr \pi'_v(f_v)
=\sum_{\overline{\pi}'\in\overline{{\cal A}}}m(\overline{\pi}')\,
\omega_{\overline{\pi}'}^S(\overline{h})
\prod_{v\in S}\tr \overline{\pi}'_v(\overline{f}_v),
\end{equation}
where $\omega_{\pi'}^S$ (resp. $\omega_{\overline{\pi}'}^S$) is the
character of ${\cal H}^S$ (resp. of its genuine analog) by which it acts
on the line of spherical vectors of $\otimes_{v\notin S}\pi'_v$ (resp. of
$\otimes_{v\notin S}\overline{\pi}'_v$).

{\em Step 2 (isolating the components outside $S$).} By
\S\ref{sec:spherical}, $\omega_{\overline{\pi}'}^S(\overline{h})$ is a
character of ${\cal H}^S$ as a function of $h$, and it equals
$\omega_{\pi''}^S(h)$ for the unramified representation
$\otimes_{v\notin S}\pi''_v$ corresponding to
$\otimes_{v\notin S}\overline{\pi}'_v$ under the Satake correspondence.
Fix $f_S$ and regard both sides of (\ref{eq:finite-identity}) as
functions of $h$. Distinct characters of an algebra are linearly
independent (Artin), so, after grouping the finitely many terms
according to their characters, we may equate the coefficients of the
character $\omega_\pi^S$. Since a $K_{0,v}$-spherical irreducible
representation is determined by its Hecke character, and the Satake
correspondence is a bijection on spherical representations, this gives
\begin{equation}\label{eq:S-identity}
\sum_{\substack{\pi'\in{\cal A}\\ \pi'_v\cong\pi_v\ (v\notin S)}}
m(\pi')\prod_{v\in S}\tr \pi'_v(f_v)
=\sum_{\substack{\overline{\pi}'\in\overline{{\cal A}}\\
\overline{\pi}'_v\leftrightarrow\pi_v\ (v\notin S)}}
m(\overline{\pi}')\prod_{v\in S}\tr \overline{\pi}'_v(\overline{f}_v),
\end{equation}
where $\overline{\pi}'_v\leftrightarrow\pi_v$ means that they correspond
under the Satake correspondence. This holds for all
$(f_v,\overline{f}_v)\in{\cal C}_v(L_v)$, $v\in S$, and hence, by
linearity, on $\otimes_{v\in S}{\cal C}_v(L_v)$.

{\em Step 3 (transfer).} By (LH2), each term on the right of
(\ref{eq:S-identity}) either vanishes identically or equals
$m(\overline{\pi}')\prod_{v\in S}\tr t_v(\overline{\pi}'_v)(f_v)$.
The representations $\overline{\pi}'$ occurring on the right are
isomorphic to one another outside $S$ (the genuine Satake correspondence
is injective), so two distinct ones differ at some $v\in S$, and by the
injectivity of $t_v$ the corresponding tuples
$(t_v(\overline{\pi}'_v))_{v\in S}$ are distinct. On the left, distinct
$\pi'$ occurring in (\ref{eq:S-identity}) have distinct tuples
$(\pi'_v)_{v\in S}$, since they agree outside $S$.

{\em Step 4 (comparison of coefficients).} For each $v\in S$, let
$\Lambda_v$ be the set of distinct representations of $G(k_v)$ which
occur as $\pi'_v$ on the left or as $t_v(\overline{\pi}'_v)$ on the right
of (\ref{eq:S-identity}), and whose characters do not vanish identically
on ${\cal C}_v(L_v)$; it is finite by Step 1. Terms involving some other
representation vanish on $\otimes_{v\in S}{\cal C}_v(L_v)$. By (LH1), the
characters of the elements of $\Lambda_v$ are linearly independent on
${\cal C}_v(L_v)$, so by Lemma~\ref{lem:tensor} the product functionals
are linearly independent on $\otimes_{v\in S}{\cal C}_v(L_v)$. By the
choice of the $L_v$, $\pi_v\in\Lambda_v$ for every $v\in S$. Comparing
the coefficients of $\otimes_{v\in S}\tr \pi_v$ on the two sides of
(\ref{eq:S-identity}) gives
\[
m(\pi)=\sum_{\overline{\pi}'}m(\overline{\pi}'),
\]
the sum over the genuine cuspidal automorphic representations
$\overline{\pi}'$ in $\overline{\cal A}$ with
$\overline{\pi}'_v\leftrightarrow\pi_v$ for
$v\notin S$ and $t_v(\overline{\pi}'_v)\cong\pi_v$ for $v\in S$. (Every
such $\overline\pi'$ lies over $\nu$, so it does lie in $\overline{\cal A}$.)
By Step 3, there is at most one such $\overline{\pi}'$, and by
Lemma~\ref{lem:A00}, $m(\overline{\pi}')\leq 1$. Hence $m(\pi)\leq 1$.
Since $\pi$ is automorphic, $m(\pi)\geq 1$, which proves (a). Then the
sum on the right equals $1$, so it has exactly one term, which proves (b).
\end{proof}

\begin{remark}
Hypothesis~\ref{hyp:GH} is only used for $\nu=\nu_\infty(\pi)$ and for
pairs with $f_S$ in $\otimes_{v\in S}{\cal C}_v(L_v)$ and arbitrary
$h\in{\cal H}^S$.
\end{remark}

\section{The local hypothesis at hyperspecial places}

\begin{proposition}\label{prop:LH-spherical}
Let $v$ be an odd finite place, $i\in\{0,1\}$, and let $\pi_v$ be an
irreducible unitary representation of $G(k_v)$ with a nonzero
$K_{i,v}$-fixed vector. Let ${\cal C}_v$ be the space of pairs
$(f_v,\overline{f}_v)$ with $f_v\in{\cal H}(G(k_v),K_{i,v})$ and
$\overline{f}_v\in{\cal H}(\overline{G(k_v)},K_{i,v})_{gen}$
corresponding under the Satake isomorphism (Definition~\ref{def:satake-corr};
for $i=1$, transported from $i=0$ by conjugation by $w$ in the
2-fold cover of $GL(2,k_v)$ of \S\ref{sec:local}). Then
Hypothesis~\ref{hyp:LH} holds for $\pi_v$ with the class ${\cal C}_v$.
\end{proposition}

\begin{proof} (LH0) holds with the single index $L=K_{i,v}$. If $\sigma$ is an
irreducible admissible representation of $G(k_v)$ whose character does
not vanish on ${\cal H}(G(k_v),K_{i,v})$, then $\sigma$ has a
$K_{i,v}$-fixed line, on which ${\cal H}(G(k_v),K_{i,v})$ acts by an
algebra character $\omega_\sigma$, and $\tr \sigma(f)=\omega_\sigma(f)$
for $f\in{\cal H}(G(k_v),K_{i,v})$. Since $\sigma$ is determined by
$\omega_\sigma$, inequivalent such $\sigma$ give distinct characters,
which are linearly independent (Artin). This proves (LH1). Similarly,
if $\overline{\tau}$ is a genuine irreducible representation whose
character does not vanish on the genuine spherical Hecke algebra, then
$\tr \overline{\tau}(\overline{f})=\omega_{\overline{\tau}}(\overline{f})$
for an algebra character $\omega_{\overline{\tau}}$. Composing with the
Satake isomorphism of \S\ref{sec:spherical} gives an algebra character
of ${\cal H}(G(k_v),K_{i,v})$, which is $\omega_\sigma$ for a unique
spherical irreducible representation $\sigma=t_v(\overline{\tau})$.
Since $\overline{\tau}$ is determined by $\omega_{\overline{\tau}}$,
$t_v$ is injective. This proves (LH2). Finally, (LH3) holds with
$f_v$ the unit element, since $\tr \pi_v(f_v)=1$. Condition (LH4) concerns
archimedean places only.
\end{proof}

\begin{corollary}\label{cor:main}
Suppose that, for each archimedean or dyadic place $v$, and for each
irreducible unitary representation $\pi_v$ of $G(k_v)$ that occurs as a
local component of a cuspidal automorphic representation, there is a
class ${\cal C}_v$ for which Hypothesis~\ref{hyp:LH} holds for $\pi_v$,
and that Hypothesis~\ref{hyp:GH} holds for these classes together with
those of Proposition~\ref{prop:LH-spherical}. Then every cuspidal
automorphic representation $\pi$ of $G({\mathbb{A}})$ such that, at each
odd place $v$, $\pi_v$ has a nonzero vector fixed by $K_{0,v}$ or by
$K_{1,v}$, occurs with multiplicity one.
\end{corollary}

\begin{proof} Combine Theorem~\ref{thm:main} with
Proposition~\ref{prop:LH-spherical}.
\end{proof}

\section{What the local chapters provide toward Hypothesis LH}
\label{sec:LH-status}

We record what Chapters~\ref{ch:reps}--\ref{ch:local} contribute to
Hypothesis~\ref{hyp:LH} at the places of $S(\pi)$, and what is missing.

\begin{enumerate}
\item {\em Odd places, $\pi_v$ hyperspecial-spherical.} Hypothesis LH
holds (Proposition~\ref{prop:LH-spherical}). This covers the two
components of a reducible unramified unitary principal series. (The
1995 draft treated the non-$K_{0,v}$-spherical component by conjugating
by an element of the cover of $GL(2,k_v)$; this is the case $i=1$ of
Proposition~\ref{prop:LH-spherical}.)

\item {\em Odd places, $\pi_v$ a ramified principal series.}
Proposition~\ref{prop:hyp-ps}, with the corrected correspondence of
characters of \S\ref{sec:charsA}, gives the character identity
$\tr I_A(\overline{\chi})(\overline{f})=\tr I_A(\chi)(f)$ for
functions corresponding oddly on the hyperbolic set, when
$\chi\neq\chi^w$. For (LH2) one would also need to control the traces of
the genuine discrete series representations on that class, and for
(LH1) a separation statement; neither is established here.
When $\chi=\chi^w$ is ramified, $I_A(\chi)$ is reducible and its two
components are not separated by such functions.

\item {\em Odd places, $\pi_v$ supercuspidal.} Chapter~\ref{ch:sc} and
\S\ref{sec:elliptic} propose a correspondence through functions
corresponding oddly on the elliptic set (Lemma~\ref{lem:ell-odd}, which
requires the character relation on every anisotropic Cartan subgroup).
See Remark~\ref{rem:sc-caveat} for a correction needed in the character
relation of Chapter~\ref{ch:sc}. For (LH2) one would also need to
control the traces of the genuine special and non-tempered
representations on that class.

\item {\em Odd places, $\pi_v$ special.} See
Remark~\ref{rem:savin-caveat}.

\item {\em Archimedean places.} \S\ref{sec:realA} and
\S\ref{sec:realds} give character identities for principal series
(through the hyperbolic matching) and for discrete series and limits of
discrete series (through the three elliptic matchings of
Lemma~\ref{lem:realds}, for functions with vanishing hyperbolic orbital
integrals). Condition (LH4) holds for these matchings (see the remark
following Hypothesis~\ref{hyp:LH}), and for the hyperbolic matching,
since corresponding principal series have the same infinitesimal
character. A verification of (LH1)--(LH3) for these classes is not
carried out here. Complex places are not treated.

\item {\em Dyadic places.} Nothing is established; see
\S\ref{sec:dyadic}.
\end{enumerate}

The trivial representation never occurs as a local component of a
cuspidal automorphic representation of $G({\mathbb{A}})$ (a consequence of
the strong approximation theorem), so it need not be considered.

\section{Remarks on the 1995 argument}\label{sec:1995-argument}

The 1995 draft stated an unconditional theorem asserting, for each
``global transfer'' (a choice of local matchings at all places), a
correspondence between cuspidal automorphic representations of
$G({\mathbb{A}})$ and genuine ones of $\overline{G({\mathbb{A}})}$, and
deduced multiplicity one. Besides its dependence on the local results
discussed in \S\ref{sec:LH-status} and on the trace identity of
Remark~\ref{rem:GH}, the final step of that argument had two gaps.

\begin{enumerate}
\item It fixed a finite set $V$ of places containing the archimedean
places, derived an identity of the form (\ref{eq:S-identity}) with
$S$ replaced by $V$, and then took $V=\{v_0\}$ for a single finite
unramified place $v_0$. This is not permitted, since $V$ must contain the
archimedean places, and the character identities it then used were only
assumed at finite places.

\item To obtain a single genuine representation on the right side of
the identity, it used a count attributed to Waldspurger \cite{Wa2}
(via \cite{PS}) of the genuine cuspidal representations
$\overline{\pi}'$ with $\overline{\pi}'_v\cong\overline{\pi}_v$ for all
$v\notin V$: namely $1$ if $\Sigma\subset V$ and
$2^{|\Sigma|-|\Sigma\cap V|-1}$ otherwise, where $\Sigma$ is the set of
places where $\overline{\pi}_v$ is a discrete series representation and
$|\Sigma\cap V|$ is assumed even. But $V$ was chosen to contain at most
one place of $\Sigma$, so the count as stated is $2^{|\Sigma|-1}$, which
is greater than $1$ as soon as $|\Sigma|\geq 2$.
\end{enumerate}

The proof of Theorem~\ref{thm:main} avoids both points: the set $S$ is
never shrunk, and uniqueness on the metaplectic side comes from the
injectivity in (LH2) together with Waldspurger's multiplicity one
theorem, rather than from a count of near-equivalent representations.

%% ===== end aabook6.tex =====


%% ===== begin aabookrf.tex =====
\begin{thebibliography}{99}

\bibitem[Ad1]{Ad1} J. Adams, ``Character of the oscillator representation'',
preprint (to appear in Isr. J. Math.)

\bibitem[Ad2]{Ad2} -------,
``Lifting of characters on orthogonal and metaplectic groups'',
Duke Math. J. \underline{92} (1998) 129--178

%\bibitem[Ad3]{Ad3} -------, Notes on the Flicker correspondence, 1995

\bibitem[Ar]{Ar} Ariturk, ``On the composition series of principal series
representations of a three-fold covering group of $SL(2,K)$'', Nagoya Math
J. \underline{77}(1980)177-196

%\bibitem[A1]{A1} J. Arthur,
%``The Selberg trace formula for $F$-rank one groups'', Ann.
%Math. \underline{100}(1974)326-385

\bibitem[A1]{A2} ------,
``A trace formula for reductive groups, I : terms associated to
classes in $G(\mathbb{Q})$'', Duke Math. J. \underline{45} (1978)911-952

\bibitem[A2]{A3}  ------,
``A trace formula for reductive groups, II : applications of a
truncation operator'', Comp. Math. \underline{40} (1980)87-121

%\bibitem[A4]{A4} ------,``Eisenstein series and the trace formula'', in
%\underline{Automorphic forms, }
%\underline{Representations, } \underline{and
%L-functions, vol. I}, 1979 Corvallis Conference, ed. A. Borel, W. Casselman,
%Proc. Symp. Pure Math. vol 33, AMS, 1979

%\bibitem[A5]{A5} ------,
%``The trace formula for non-compact quotient'', Proc. Inter.
%Cong. Math., Warsaw, Poland, 1983

\bibitem[A3]{A6} ------, ``The trace formula in invariant form'', Ann. Math.
\underline{114}(1981)

\bibitem[A4]{A7} ------, ``The local behaviour of weighted orbital integrals'',
Duke Math. J. \underline{56}(1988)223-293

%\bibitem[A8]{A8} ------, ``A measure on the unipotent variety'', Can. J. Math.
%\underline{37}(1985)1237-1274

%\bibitem[A9]{A9} ------, ``The invariant trace formula, I. Local theory'',
%J.A.M.S.\underline{1}(1988)323-383

%\bibitem[A10]{A10} ------, ``The invariant trace formula, II. Global theory'',
%J.A.M.S.\underline{1}(1988)501-554

%\bibitem[A11]{A11} ------, ``Intertwining operators and residues, I: weighted
%characters'', J. Func. Anal.\underline{84}(1989)19-84

%\bibitem[A12]{A12} ------, ``A family of distributions obtained from
%orbits'', Can. J. Math. \underline{38}(1986)179-214

%\bibitem[A13]{A13} ------, ``A family of distributions obtained from
%Eisenstein series, I'', preprint

%\bibitem[A14]{A14} ------, ``A family of distributions obtained from
%Eisenstein series, II,'' preprint

\bibitem[A5]{A15} ------,``The trace formula for reductive groups,'' in
\underline{Journ\'ees Automorphes}, Publ. Math. Univ. Paris VII,
U.E.R. Math, 1983

%\bibitem[A16]{A16} ------,``On the Fourier transform of weighted
%orbital integrals,'' preprint

\bibitem[APW]{APW} J. Arthur, ``A Paley--Wiener theorem for real reductive
groups'', Acta Math. \underline{150} (1983) 1--89

\bibitem[AC]{AC} ----- and L. Clozel, \underline{Simple algebras,
base change, and the advanced} \underline{theory of the trace formula}
Annals of Math. Studies, \underbar{120}, Princeton Univ. Press

%\bibitem[As]{As} M. Assem,
%``The Fourier transform and some character formulae for p-adic
%$SL(\ell )$, $\ell $ a prime'', preprint

%\bibitem[Ba]{Ba} G. Bachman,
%\underline{Elements} \underline{of} \underline{abstract}
%\underline{harmonic} \underline{analysis}, Academic Press, New York, 1965

%\bibitem[B]{B} J. Bernstein, ``All reductive p-adic groups are tame'',
%Funct. Anal. Appl. \underline{8}(1974)91-93

%\bibitem[BD]{BD} ----, ``Le centre de Bernstein'', (notes by P. Deligne),
%in \underline{Repr\'esentations} \underline{des}
%\underline{groupes} \underline{r\'eductifs}
%\underline{sur} \underline{un} \underline{corp} \underline{local},
%Hermann, Paris, 1984

%\bibitem[BDK]{BDK} J. Bernstein, P. Deligne, D. Kazhdan, ``Trace
%Paley-Wiener theorem for
%reductive p-adic groups'', J. d'Anal. Math. \underline{47}(1986)180-192

%\bibitem[BZ1]{BZ1} J. Bernstein, A. Zelevinski, ``Representations of
%the group $GL(n,F)$,
%where $F$ is a non-archimedean local field'', Russian Math Surveys
%\underline{31}(1976)1-68

%\bibitem[BZ2]{BZ2} ----, ``Induced representations of reductive p-adic
%groups, I'', Ann. scient. \'Ec. norm. Sup. \underline{10}(1977)441-472

\bibitem[Blas]{Blas} D. Blasius, ``On multiplicities for $SL(n)$'',
Israel J. Math. \underbar{88}(1994)237-251

\bibitem[Bl]{Bl} C. Blondel, ``Les repr\'esentations supercuspidales
des groupes
m\'etaplectiques sur GL(2) et leurs caract\'eres'',
Bull de la Soc. Math. France, Memoir 18, 1985

%\bibitem[Bou]{Bou} A. Bouaziz, ``Int\'egrales orbitales sur les groupes
%de Lie r\'eductifs'', Ann. Sc. \'Ec. Norm. Sup. \underline{27}(1994)573-609

\bibitem[Bo]{Bo} A. Borel, et al, \underline{Seminar in algebraic groups and
related finite groups}, Springer LNM \underline{131}, 1970

\bibitem[BJ]{BJ} A. Borel and H. Jacquet, ``Automorphic forms and automorphic
representations'', in \underbar{Automorphic} \underbar{forms},
\underbar{representations}, \underbar{and} \underbar{L-functions},
vol I, ed. A. Borel, W. Casselman, Proc. Symp. Pure Math. 33, AMS,
Providence, RI, 1979

\bibitem[CZ]{CZ} J. Chai and Q. Zhang, ``A strong multiplicity one theorem
for $SL_2$'', Pacific J. Math. \underline{285} (2016) 345--374


\bibitem[Car]{C} P. Cartier, ``Representations of p-adic groups: a survey'',
in \underbar{Automorphic} \underbar{forms},
\underbar{representations}, \underbar{and} \underbar{L-functions},
vol I, ed. A. Borel, W. Casselman, AMS, Providence, RI, 1979

%\bibitem[Ca]{Ca} W. Casselman, ``Introduction to the theory of admissible
%representations on reductive p-adic groups'', preprint

%\bibitem[Cl]{Cl} L. Clozel, ``Invariant harmonic analysis
%on the Schwartz space of a
%reductive p-adic group'', in \underbar{Harmonic} \underbar{analysis} %
%\underbar{on} \underbar{reductive} \underbar{groups}, ed. W. Barker, P.
%Sally, Birkh\"auser, Boston, 1991

%\bibitem[CD]{CD} ------ and P. Delorme, ``Le th\'eor\`eme de Paley-Wiener
%invariant pour les groupes de Lie reductif, II'',
%Ann. Scient. Ec. Norm. Sup., \underline{23}(1990)193-228

%\bibitem[D]{D} J. Dixmier, \underbar{$C^{*}$-algebras}, North-Holland,
%Amsterdam, 1977

%\bibitem[Fe]{Fe} J. Fell, ``$C^{*}$-algebras with smooth dual'', Ill. J.
%Math. \underline{4}(1960)221-230

%\bibitem[Fla]{F} D. Flath, ``Decomposition of representations into tensor
%products'', in \underline{Automorphic forms, }
%\underline{Representations, } \underline{and
%L-functions}, 1979 Corvallis Conference, ed. A. Borel, W. Casselman,
%Proc. Symp. Pure Math. vol 33, AMS, 1979

\bibitem[Fl]{Fl} Y. Flicker, ``Automorphic forms on covering groups of
$GL(2)$'', Inv
Math \underline{57}(1980)119-182

%\bibitem[Fr]{Fr} S. Franklin, ``The reducible principal series of
%SL(2) over a p-adic field'', Thesis, Univ. Chicago, 1971

\bibitem[Ge1]{Ge1} S. Gelbart, \underline{Automorphic forms on adele groups},
Annals of
Math. Studies, Princeton Univ. Press, Princeton, 1975

\bibitem[Ge2]{Ge2} S. Gelbart, \underline{Weil's representation and the
spectrum of the metaplectic} \underline{group},
SLN Math \underline{530}, 1976

\bibitem[GJ]{GJ} S. Gelbart and H. Jacquet,
``Forms of $GL(2)$ from the analytic
point of view'', in \underline{Automorphic forms},
\underline{representations}, and \underline{L-functions},
(ed. A. Borel, W. Casselman), Proc. Symp. Pure Math, vol 33, AMS, 1979

%\bibitem[GeKn]{GeKn} ------ and T. Knapp, ``L-indistinguishability and R
%groups for the special linear group'', Adv. Math. 43(1982)101-121

\bibitem[GJ2]{GJ78} S. Gelbart and H. Jacquet, ``A relation between
automorphic representations of $GL(2)$ and $GL(3)$'', Ann. Sci. \'Ec.
Norm. Sup. \underline{11} (1978) 471--542

\bibitem[GePS]{GePS} S. Gelbart and I. Piatetski-Shapiro, ``Distinguished
representations and modular forms of half-integral weight'', Invent.
Math. \underline{59} (1980) 145--188

\bibitem[GPS]{GPS} S. Gelbart and I. Piatetski-Shapiro, ``Some remarks
on metaplectic cusp
forms and the correspondences of Shimura and Waldspurger'', Israel
J. Math.
\underline{44}(1983)97-126

\bibitem[GS]{GS} ------ and P. Sally, ``Intertwining operators and
automorphic forms on the metaplectic group'', Proc. Nat. Acad. Sci.
\underline{72}(1975)1406-1410

%\bibitem[GG]{GG} I. M. Gelfand, M. I. Graev,
%``Representations of a group of the second
%order with elements from a locally compact field'', Russian Math. Surveys,
%18(1963)29-100

%\bibitem[GGPS]{GGPS} ------ and I. Piatetski-Shapiro,
%\underline{Representation
%theory} \underline{and} \underline{automorphic functions}, W. B. Saunders,
%Philadelphia, 1969

%\bibitem[Go1]{Go1} R. Godement, ``Domaines fondamentaux des groupes
%arithm\'etiques'', Sem. Bourbaki, no. 257, 1962/63

%\bibitem[Go2]{Go2} ------, ``The spectral decomposition of cusp forms'', in
%\underline{Algebraic} \underline{groups}
%\underline{and} \underline{discontinuous}
%\underline{subgroups}, Proc Symp Pure Math, vol. 9, AMS, Providence, 1966

%\bibitem[Go3]{Go3} ------, ``Analyse spectrale des functions modulaires'',
%Sem Bourbaki, no 278, 1964/65

\bibitem[HC1]{HC} Harish-Chandra, \underline{Harmonic analysis on
reductive p-adic
groups}, (notes by G. van Dijk), SLN Math \underline{162}, 1970

\bibitem[HC2]{HC2} -------, \underbar{Collected Works of
Harish-Chandra, vol III}, (ed. V. Varadarajan et al), Springer-Verlag

\bibitem[HC3]{HC68} -------, \underline{Automorphic forms on semisimple
Lie groups}, (notes by J. G. M. Mars), SLN Math \underline{62}, 1968

%\bibitem[H1]{H1} L. H\"ormander,
%\underbar{An} \underbar{introduction} \underbar{to}
%\underbar{complex} \underbar{analysis} \underbar{in}
%\underbar{several} \underbar{complex} \underbar{variables},
%2nd edition, North-Holland, Amsterdam, 1973

%\bibitem[H2]{H2} -----, \underbar{The} \underbar{analysis}
%\underbar{of} \underbar{linear} \underbar{partial}
%\underbar{differential} \underbar{operators}, \underbar{I},
%Springer-Verlag, New York, 1983

%\bibitem[HT]{HT} R. Howe, E. C. Tan,
%\underline{Non-abelian harmonic analysis}, Springer-Verlag, 1992

\bibitem[H]{H} R. Howe, ``Character of the oscillator representation'',
(appendix to ``Elements of Oscillator Representations'',
unpublished notes)

\bibitem[HM]{HM} R. Howe (with the collaboration of A. Moy),
\underbar{Harish-Chandra}
\newline
\underbar{homomorphisms} \underbar{for}
\underbar{p-adic} \underbar{groups}, CBMS \underbar{59}, AMS (1985).

%\bibitem[Hu]{Hu} J. E. Humphreys, \underline{Arithmetic groups}, Springer
%Lecture Notes in Math., vol. \underline{789} (1980)

\bibitem[JL]{JL} H. Jacquet and R. Langlands,
\underline{Automorphic} \underline{forms}
\underline{on} \underline{$GL(2)$}, Springer Lecture Notes in Math.,
vol 114, 1970

\bibitem[J1]{J} D. Joyner, ``The metaplectic analog of of Kazhdan's
endoscopic lift'', Israel J. Math. \underline{66}(1988)113-154

\bibitem[J2]{J2} ------, ``An application of Arthur's invariant
trace formula to $SL(2)$'', preprint, Feb/1996

\bibitem[J3]{J3} ------, ``A correspondence for the generalized Hecke
algebra of the metaplectic cover $\overline{SL(2,F)}$, $F$ p-adic'',
New York J. Math. \underbar{4} (1998); revised version in preparation (2026)

%\bibitem[K]{K} D. Kazhdan,
%``Cuspidal geometry of p-adic groups'', J. d'Anal. Math.
%47(1986)1-36

%\bibitem[KP]{KP} ------ and  S. Patterson, ``Metaplectic forms'',
%Publ. Math. IHES \underline{59}(1984)35-142

%\bibitem[KP2]{KP2} ------, ``A Shimura correspondence'',
%Adv. Math. \underline{}(1986)

%\bibitem[Ko1]{Ko1} R. Kottwitz, ``Stable trace formula: elliptic singular
%terms'', Math. Ann. \underline{275}(1986)365-399

%\bibitem[Ko2]{Ko2} ------, ``Stable trace formula: cuspidal tempered terms'',
%Duke Math. J. \underline{51}(1984)611-650

\bibitem[K1]{K1} T. Kubota, ``Topological covering of $SL(2)$ over a
local field'', J.
Math. Soc. Japan \underline{19} (1967)114-121

\bibitem[K2]{K2} ------, \underline{Automorphic functions and the reciprocity
law in a number field}, Kyoto U., 1969

%\bibitem[K3]{K3} ------, \underline{Eisenstein series for $SL(2,\mathbb{R})$},
%John Wiley and Sons, 1973

\bibitem[Kut]{Kut} P. Kutzko,
``Character formulas for supercuspidal representations of
$GL(\ell)$, $\ell$ a prime'', Amer J Math \underline{109}(1987)201-222

\bibitem[Lab]{Lab} J.-P. Labesse, ``L-indistinguishable representations
and the trace formula for $SL(2)$'', in
\underbar{Lie groups and their representations}, John Wiley, 1975

\bibitem[LL]{LL} ------- and R. P. Langlands,
``L-indistinguishability for
$SL(2)$'', Can. J. Math. \underline{33}(1979)

%\bibitem[La]{La} S. Lang, \underline{$SL(2,{\mathbb{R}})$}, Addison-Wesley, 1975

\bibitem[L1]{L1} R. Langlands, ``Eisenstein series'', Proc Symp Pure Math,
vol. 9, AMS, Providence, 1966

\bibitem[L2]{L2} ------, \underline{On} \underline{the} \underline{functional}
\underline{equations} \underline{satisfied}
\underline{by} \underline{Eisenstein}
\underline{series}, Springer Lecture Notes, \underline{544} (1976)

%\bibitem[L3]{L3} ------, \underline{Les d\'ebuts d'une formule des traces
%stable}, Publ. Math. de L'Univ. Paris, VII, 1983

\bibitem[L3]{L4} ------, \underline{Base} \underline{change}
\underline{for $GL(2)$},
Annals of Math. Studies, Princeton

%\bibitem[L5]{L5} ------, ``On the notion of an automorphic
%representation'', in \underbar{Automorphic} \underbar{forms},
%\underbar{representations}, \underbar{and} \underbar{L-functions},
%vol I, ed. A. Borel, W. Casselman, AMS, Providence, RI, 1979

%\bibitem[Li]{Li}
%R. Lipsman, \underline{Group representations: A survey of some current
%topics}, Springer LNM \underline{388}, 1974

%\bibitem[Ma]{Ma} G. Mackey, ``Borel structure in groups and their duals'',
%Trans. AMS \underline{85}(1957)134-165

%\bibitem[Mac]{Mac} G. Mackey, ``Unitary representations of group extensions, I'',
%Acta Math. \underbar{99}(1958)265-311

\bibitem[M1]{M1} D. Manderscheid, ``On the supercuspidal representations of
$SL(2)$ and its 2-fold cover, I'', Math. Ann. \underbar{266}(1987)287-295

\bibitem[M2]{M2} ------, ``On the supercuspidal representations of
$SL(2)$ and its 2-fold cover, II'', Math. Ann.
\underbar{266}(1987)297-305

\bibitem[M3]{M3} ------, ``Supercuspidal representations and the
theta correspondence'', preprint

\bibitem[MW]{MW} C. Moeglin, J.-L. Waldspurger,
\underline{Spectral decomposition and Eisenstein series},
Cambridge Univ Press, 1995

\bibitem[Mo1]{Mo1} C. Moen, ``Intertwining operators for covering groups
of $SL(2)$ over
a local field'', preprint

\bibitem[Mo2]{Mo2} ------, ``Irreducibility of unitary principal series
representations
for covering groups of $SL(2)$'', Pac. J. Math. 135(1988)89-110

\bibitem[Mo3]{Mo3} -----, ``The dual pair $(U(1),U(1))$ over a p-adic field'',
Pac. J. Math. \underbar{158}(1993)

\bibitem[M]{M} C. Moore, ``Group extensions of p-adic linear groups'',
Publ. Math. IHES \underline{35} (1968)

%\bibitem[O]{O} T. Ono, ``On Tamagawa numbers'', in Proc. Symp. Pure Math.,
%vol. 9, AMS, Providence, 1966

\bibitem[PS]{PS} I. Piatetski-Shapiro, ``Work of Waldspurger'', in Lie groups
representations, II (ed. R. Herb, et al), Springer LNM
\underline{1041},1983

%\bibitem[PR]{PR} G. Prasad, M. Raganathan,
%``Topological central extensions of
%semi-simple groups over local fields'', Annals of Math.
%\underline{119}(1984)143-268

%\bibitem[RS]{RS} C. Rader, A. Silberger, ``Some consequences of
%Harish-Chandra's submersion principle'', Proc. AMS \underbar{118}
%(1993)1271-1279

\bibitem[Rao]{Rao} R. Rao, ``On some explicit formulas in the theory
of Weil representation'', Pac. J. Math. \underbar{157}(1993)335-371

%\bibitem[Ro]{Ro} A. Robert, \underline{Introduction} \underline{to}
%\underline{the} \underline{representation} \underline{theory} \underline{of}
%\underline{compact} \underline{and} \underline{locally}
%\underline{compact} \underline{groups},
%London Math. Series, Cambridge Univ. Press, Cambridge, 1983

\bibitem[Rog]{R} J. Rogawski,
\underline{Automorphic representations of unitary groups
in three} \underline{variables}
Annals of Math. Studies, \underline{123}, Princeton Univ. Press, 1990

%\bibitem[Ru]{Ru} W. Rudin, \underline{Functional} \underline{Analysis},
%McGraw-Hill, New York, 1973

\bibitem[Ra]{Ra} D. Ramakrishnan, ``Modularity of the Rankin-Selberg
$L$-series, and multiplicity one for $SL(2)$'', Ann. of Math.
\underbar{152} (2000) 45--111

\bibitem[Sa]{Sa} P. Sally, ``Analytic continuation of the irreducible unitary
representations of the covering group of $SL(2,{\mathbb{R}})$'', Memoirs AMS
\underline{69}, 1967

%\bibitem[SS1]{SS1} ------- and J. Shalika,
%``Characters of the discrete series of
%representations of SL(2) over a p-adic field'', Proc. Acad. Sci. USA
%61(1968)1231-1237

%\bibitem[SS2]{SS2} ------, ``The Fourier transform of orbital integrals
%on SL(2) over a p-adic field'', in \underbar{Lie} \underbar{Group}
%\underbar{Representations}, \underbar{II},
%ed., R. Herb et al, Lecture Notes in Math. vol 1041,
%Springer-Verlag, New York, 1984

\bibitem[Sav]{Sav} G. Savin, ``A local Shimura correspondence'',
Math Ann. \underbar{280}(1988)185-190
\bibitem[Sch]{Sch} L. Schwartz, \underbar{Theorie} \underbar{des}
\underbar{distributions}, \underbar{vol. I}, Hermann, Paris, 1950

%\bibitem[Sen]{Sen} S. Sen, ``Explicit reciprocity laws, II', Crelle's J.,
%\underline{323}(1981)

%\bibitem[Se]{Se} J.-P. Serre,  \underline{Representations of finite groups},
%Springer-Verlag

%\bibitem[Sh]{Sh} F. Shahidi, ``Local coefficients as Artin factors
%for real groups'', Duke Math J. \underbar{52}(1985)973-1007

\bibitem[Shal]{Shal} J. Shalika, ``Representations of the two by two
unimodular group over local fields'', IAS

%\bibitem[Sha]{Sha} S. Shatz, \underline{Profinite groups, arithmetic, and
%geometry}, Princeton Univ. Press, Princeton, 1972

%\bibitem[She]{She} D. Shelstad, ``L-indistinguishability for real groups'',
%Math Ann \underbar{259}(1982)385-430

\bibitem[Sho]{Sho} S. Shokranian, \underline{The Selberg-Arthur trace formula},
Springer Lecture Notes in Math., vol. \underline{1503} (1992)

\bibitem[S]{S} A. Silberger, \underbar{Introduction} \underbar{to}
\underbar{harmonic} \underbar{analysis} \underbar{on} \underbar{reductive}
\underbar{p-adic} \underbar{groups}, Princeton Univ. Press, Princeton, 1979

%\bibitem[SpSt]{SpSt} T. Springer and R. Steinberg, ``Conjugacy classes'',
%in \underbar{Seminar on Algebraic}
%\underbar{Groups ...} (by A. Borel, et al),
%Springer LNM \underbar{131}, 2nd printing 1986

%\bibitem[T]{T} M. Tadi\'c, ``The topology of the dual space of a
%reductive group over a
%local field'', Glasnik Matem., \underline{18}(1983)259-279

%\bibitem[Ta]{Ta} M. Taibleson, \underbar{Fourier} \underbar{analysis}
%\underbar{on} \underbar{local} \underbar{fields}, Princeton Univ.
%Press, Princeton, 1975

%\bibitem[Va]{Va} V. Varadarajan,
%\underline{An introduction to harmonic analysis on semisimple}
%\underline{Lie groups}
%Cambridge Studies in Adv. in Math \underline{16}, Cambridge, 1986

%\bibitem[Vig]{V} M.-F. Vign\'eras, ``Caract\'erisation des integrales
%orbitales sur un
%groupe r\'eductif p-adique'', J. Fac. Sci. Univ. Tokyo
%\underline{28}(1981)945-961

\bibitem[TW]{TW} S. Takeda and A. Wood, ``Hecke algebra correspondences
for the metaplectic group'', Trans. Amer. Math. Soc. \underbar{370}
(2018) 1101--1121

\bibitem[Zh]{Zh} Q. Zhang, ``A strong multiplicity one theorem for $SL_2$'',
Pacific J. Math. \underline{292} (2018) 283ff.

\bibitem[Wa1]{Wa1} J.-L. Waldspurger, ``Correspondance de Shimura'',
J. Math. Pure Appl. \underline{59}(1980)1-133

\bibitem[Wa2]{Wa2} -------, ``Correspondance de Shimura et quaternions'',
Forum Math \underbar{3}(1991)

\bibitem[We]{W1} A. Weil, \underline{Basic Number Theory},
Springer-Verlag, 1975

\bibitem[We2]{W2} -----, \underline{Adeles and algebraic groups},
Birkhauser, 1982


\end{thebibliography}
%% ===== end aabookrf.tex =====



\end{document}
